\documentclass[12pt]{amsart}
\usepackage[usenames,dvipsnames]{xcolor}
\usepackage{amsmath,amssymb,amsthm,graphicx,mathrsfs,url,latexsym,enumerate}
%%* \usepackage{a4wide}
\usepackage{ifthen}
\usepackage{verbatim}
\usepackage{fancyhdr}
\usepackage{url}
\usepackage{bbm}
\usepackage[sans]{dsfont}
%\usepackage{showkeys}
\usepackage{transparent}
\usepackage{amsbsy}
%\usepackage[justification=centering]{caption}
\usepackage[colorlinks=true,linkcolor=Red,citecolor=Green]{hyperref}
%%* \usepackage{wrapfig}
\usepackage{tikz}

\definecolor{bleu}{RGB}{27,88,145}
\definecolor{mauve}{RGB}{138,20,79}

%% \setlength{\textheight}{8.50in} \setlength{\oddsidemargin}{0.00in}
%% \setlength{\evensidemargin}{0.00in} \setlength{\textwidth}{6.08in}
%% \setlength{\topmargin}{0.00in} \setlength{\headheight}{0.18in}
%% \setlength{\marginparwidth}{1.0in}
%% \setlength{\abovedisplayskip}{0.2in}
%% \setlength{\belowdisplayskip}{0.2in}
%% \setlength{\parskip}{0.05in}

\newcommand{\mc}[1]{\mathcal{#1}}
\renewcommand{\Im}{\operatorname{Im}}
\renewcommand{\Re}{\operatorname{Re}}
\newcommand{\dist}{\operatorname{dist}}
\newcommand{\Id}{\operatorname{Id}}
\newcommand{\Op}{\operatorname{Op}}
\newcommand{\Ent}{\operatorname{Ent}}
\newcommand{\Var}{\operatorname{Var}}
\newcommand{\E}{\mathbb E}
\newcommand{\rank}{\operatorname{Rank}}
\newcommand{\supp}{\operatorname{supp}}
\newcommand{\vect}{\operatorname{Vect}}
\newcommand{\diag}{\operatorname{diag}}
\renewcommand{\span}{\operatorname{span}}
\newcommand{\C}{\mathbb C}
\newcommand{\CC}{\mathscr{C}}
\newcommand{\N}{\mathbb N}
\newcommand{\R}{\mathbb R}
\newcommand{\indic}{\operatorname{1\negthinspace l}}
\newcommand{\pii}{\varpi}
\newcommand{\Hess}{\operatorname{Hess}}
\newcommand{\Cl}{\operatorname{Cl}}
\newcommand{\Ran}{\operatorname{Ran}}
\newcommand{\Ker}{\operatorname{Ker}}
\newcommand{\bsigma}{\boldsymbol{\sigma}}
\newcommand{\ulambda}{\underline{\lambda}}
\newcommand{\trt}{\tilde {\operatorname{tr}}}
\newcommand{\Tr}{{\operatorname{Tr}}}
\renewcommand{\div}{\operatorname{div}}
\def\<{\langle}
\def\>{\rangle}
\newcommand{\bp}{{\it Proof. }}
\newcommand{\ep}{\hfill $\square$}

%% \newcommand{\begin{equation}}{\begin{equation}}
%% \newcommand{\end{equation}}{\end{equation}}

%% \newcommand{\begin{equation*}}{\begin{equation*}}
%% \newcommand{\end{equation*}}{\end{equation*}}

\renewcommand{\d}{\operatorname{d}}

\numberwithin{equation}{section}
\numberwithin{figure}{section}

\def\m{\mathbf{m}}
\def\s{\mathbf{s}}
\def\j{\mathbf{j}}
\def\k{\mathbf{k}}
\def\um{\underline{\m}}



%%%%%%%%%%%%%%%%%%%%%%%%%%%%%%%%%%%%%%%%%%%%%%
%%%%%%%%%%%%%%   macro loïs   %%%%%%%%%%%%%%%%
%%%%%%%%%%%%%%%%%%%%%%%%%%%%%%%%%%%%%%%%%%%%%%


\everymath{\displaystyle}
%\renewcommand{\Pr}{\Pi_\rho}
\DeclareMathOperator{\p}{\partial}
\DeclareMathOperator{\eps}{\varepsilon}
\DeclareMathOperator{\phii}{\varphi}
\newcommand{\dx}{\ dx}
\newcommand{\dy}{\ dy}
\newcommand{\dz}{\ dz}
\newcommand{\dt}{\ dt}
\newcommand{\ds}{\ ds}
\newcommand{\du}{\ du}
\newcommand{\dv}{\ dv}
%\newcommand{\vvvert}[1]{{\left\vert\kern-0.25ex\left\vert\kern-0.25ex\left\vert #1 
    %\right\vert\kern-0.25ex\right\vert\kern-0.25ex\right\vert}}
%\newcommand{\VVVert}[1]{{\vert\kern-0.25ex\vert\kern-0.25ex\vert #1 
    %\vert\kern-0.25ex\vert\kern-0.25ex\vert}}
\newcommand{\vvert}[1]{\left\lVert#1\right\rVert}
%\newcommand{\VVert}[1]{\lVert#1\rVert}
\newcommand{\lb}{\operatorname{[\negthinspace [}}
\newcommand{\rb}{\operatorname{]\negthinspace ]}}


%%%%%%%%%%%%%%%%%%%%%%%%%%%%%%%%%%%%%%%%%%%%%%
%%%%%%%%%                     %%%%%%%%%%%%%%%%
%%%%%%%%%%%%%%%%%%%%%%%%%%%%%%%%%%%%%%%%%%%%%%

% THEOREM -------------------------------------------------------

\newtheorem{theorem}{Theorem}
\newtheorem{defin}{Definition}[section]
\newtheorem{corollary}[defin]{Corollary}
\newtheorem{lemma}[defin]{Lemma}
\newtheorem{sublem}[defin]{Sub-Lemma}
\newtheorem{proposition}[defin]{Proposition}
\newtheorem{remark}[defin]{Remark}
\newtheorem{example}[defin]{Example}
\newtheorem{assumption}{Assumption}

% ALPHABET CALLIGRAPHIQUE CAPITAL -----------------------------------------------
% trois fois la lettre !

\def\aaa{{\mathcal A}}\def\bbb{{\mathcal B}}\def\ccc{{\mathcal C}}\def\ddd{{\mathcal D}}
\def\eee{{\mathcal E}}\def\fff{{\mathcal F}}\def\ggg{{\mathcal G}} \def\hhh{{\mathcal H}}
\def\iii{{\mathcal I}} \def\jjj{{\mathcal J}}\def\kkk{{\mathcal K}}\def\lll{{\mathcal L}}
\def\mmm{{\mathcal M}}\def\nnn{{\mathcal N}} \def\ooo{{\mathcal O}}\def\ppp{{\mathcal P}}
\def\qqq{{\mathcal Q}}\def\rrr{{\mathcal R}}\def\sss{{\mathcal S}}\def\ttt{{\mathcal T}}
\def\uuu{{\mathcal U}}\def\vvv{{\mathcal V}}\def\www{{\mathcal W}}\def\xxx{{\mathcal X}}
\def\yyy{{\mathcal Y}} \def\zzz{{\mathcal Z}}


% ALPHABET CALLIGRAPHIQUE CAPITAL ROND ------------------------------------------

\def\Ar{{\mathscr A}}\def\br{{\mathscr B}}\def\Cr{{\mathscr C}}\def\Dr{{\mathscr D}}
\def\Er{{\mathscr E}}\def\Fr{{\mathscr F}}\def\Gr{{\mathscr G}}\def\Hr{{\mathscr H}}
\def\Ir{{\mathscr I}}\def\Jr{{\mathscr J}}\def\Kr{{\mathscr K}}\def\Lr{{\mathscr L}}
\def\Mr{{\mathscr M}}\def\Nr{{\mathscr N}}\def\Or{{\mathscr O}}\def\Pr{{\mathscr P}}
\def\Qr{{\mathscr Q}}\def\Rr{{\mathscr R}}\def\Sr{{\mathscr S}}%\def\Tr{{\mathscr T}}
\def\Ur{{\mathscr U}}\def\Vr{{\mathscr V}}\def\Wr{{\mathscr W}}\def\Xr{{\mathscr X}}
\def\Yr{{\mathscr Y}}\def\Zr{{\mathscr Z}}

% MATHS -----------------------------------------------------
\title{Hypocoercivity and metastability of degenerate KFP equations at low temperature}

\author{Loïs Delande}
\address{CERMICS, École des Ponts Champs-sur-Marne, France; {\normalfont and} MATHERIALS, Inria Paris, Paris, France}

\begin{document}
\begin{abstract}
We consider Kramers-Fokker-Planck operators with general degenerate %%
%% Annotation 4, 4/5 on page 1 [ ]
%% Replace: "consider Kramers<SEL>-Fokker-</SEL>Planck operators"
%% Comment: "--Fokker--"
%%
%⭣ ⭣ ⭣
coefficients. We prove semiclassical %%
%⭡ ⭡ ⭡ END of annotation 4
hypocoercivity estimates for a large class of such operators. Then, we manage to prove Eyring–Kramers formulas for the bottom of the spectrum of some particular degenerate operators in the semiclassical regime, and quantify the spectral gap separating these eigenvalues from the rest of the spectrum. The main ingredient is the construction of sharp Gaussian quasimodes through an adaptation of the WKB method.
\end{abstract}

\maketitle

\section{Introduction}

\subsection{Motivations}

When considering a cloud of particles in dimension $d$, one is often led to study kinetic equations, or stochastic processes. An example of such process is given by the Langevin dynamics:
\begin{equation}\label{eq:SDE0}
\left\{
\begin{aligned}
&dx_t = v_tdt,\\
&dv_t = -\partial_xV(x_t)dt - 2v_tdt + \sqrt{2h}dB_t,
\end{aligned}
\right.
\end{equation}
where $(x_t,v_t)$ denotes the position and velocity of %%
%% Annotation 1, 1/5 on page 1 [ ]
%% Ink: "None, annot.type = <Ink>"
%% Comment: ""
%%
%⭣ ⭣ ⭣
the particles at %%
%⭡ ⭡ ⭡ END of annotation 1
time $t$, $V:\R^d\to\R$ is a smooth potential corresponding to an energetic field constraining the studied particles, $h>0$ is a semiclassical parameter (typically proportional to the temperature of the system), and $B_t$ is a Brownian motion in $\R^d$ representing random forces.

To obtain results on the long-time behavior of the solution to a stochastic differential equation (SDE), one can for example study its %%
%% Annotation 5, 5/5 on page 1 [ ]
%% Replace: "random forces.  <SEL>To obtain results on the long-time behavior of the solution to a sto- chastic differential equation (SDE), o</SEL>ne can"
%% Comment: "O"
%%
%⭣ ⭣ ⭣
generator. For a general SDE
\begin{equation}\label{eq:SDEgeneral}
dX_t = b(X_t)dt + \sigma(X_t)dB_t,
\end{equation}
with %%
%⭡ ⭡ ⭡ END of annotation 5
drift $b:\R^d\to\R^d$ and diffusion matrix $\sigma:\R^d\to\mmm_d(\R)$, it is defined as
\begin{equation*}
\lll = \frac12\sum_{i,j=1}^da_{i,j}(x)\partial_i\partial_j + \sum_{i=1}^db_i(x)\partial_i,
\end{equation*}
acting on functions on $\R^d$, denoting $(a_{i,j})_{i,j} = \sigma\sigma^T$%%
%% Annotation 2, 2/5 on page 1 [ ]
%% Ink: "None, annot.type = <Ink>"
%% Comment: ""
%%
%% Annotation 3, 3/5 on page 1 [ ]
%% Delete: "1  bi(x)∂i<SEL>,</SEL>  ai,j(x)∂i∂j +"
%% Comment: ""
%%
%⭣ ⭣ ⭣
. The study of $\lll$ is important, because under mild hypotheses, we have that the %%
%⭡ ⭡ ⭡ END of annotations 2, 3
solution to the partial differential equation (PDE)
\begin{equation}\label{eq:FeynmanKac}
\left\{
\begin{aligned}
&\partial_tu+\lll u = 0,\\
&u_{|t=0}=u_0,
\end{aligned}
\right.
\end{equation}
is given by $u(t,x) = \E(u_0(X_t)\ |\ X_0=x)$%%
%% Annotation 6, 1/3 on page 2 [ ]
%% Ink: "None, annot.type = <Ink>"
%% Comment: ""
%%
%⭣ ⭣ ⭣
, with $(X_t)_t$ %%
%⭡ ⭡ ⭡ END of annotation 6
a solution to \eqref{eq:SDEgeneral}.

Therefore, the generator of \eqref{eq:SDE0} is
\begin{equation*}
\lll_{KFP} = v\cdot \partial_x - \partial_x V\cdot \partial_v + h\Delta_v - 2v\cdot\partial_v,
\end{equation*}
where $\Delta_v$ denotes the Laplacian in the $v$ variables only. Instead of working with $\lll_{KFP}$, it is very convenient to work with a conjugation of this operator. In the following, we will denote $Q^*$ the (formal) adjoint of any operator $Q$. Observing that (formally) $\lll_{KFP}1=0$ %%
%% Annotation 8, 3/3 on page 2 [ ]
%% Highlight: "this operator. <SEL>In the following, we will denote Q∗the (formal) ad- joint of any operator Q.</SEL> Observing that"
%% Comment: ""
%%
%⭣ ⭣ ⭣
and $\lll_{KFP}^*(e^{-2(V(x) + \frac{|v|^2}{2})/h})=0$, we therefore introduce
\begin{equation*}
P_{KFP} = -e^{(V(x) + \frac{|v|^2}{2})/h}\circ h\lll_{KFP}^*\circ e^{-(V(x) + \frac{|v|^2}{2})/h}.
\end{equation*}
This %%
%⭡ ⭡ ⭡ END of annotation 8
operator is called the (semiclassical) Kramers-Fokker-Planck operator and takes the form
\begin{equation}\label{eq:standardKFP}
P_{KFP} = v\cdot h\partial_x - \partial_x V\cdot h\partial_v - h^2\Delta_v + |v|^2 - hd.
\end{equation}
It has the nice form of the sum of a (formally) skew-adjoint operator and a (formally) self-adjoint one. The skew part being a transport term, and the self-adjoint part being a harmonic oscillator in velocity. Moreover, we observe that (formally)
\begin{equation*}
P_{KFP}(e^{-(V(x) + \frac{|v|^2}{2})/h}) = P_{KFP}^*(e^{-(V(x) + \frac{|v|^2}{2})/h}) = 0.
\end{equation*}

The operator $P_{KFP}$ has been extensively studied in the literature, in particular when trying to compare it with the Witten Laplacian associated with $V$, that is
\begin{equation*}
\Delta_V = -h^2\Delta + |\nabla V|^2 - h\Delta V.
\end{equation*}
A general theory regrouping the study of %%
%% Annotation 7, 2/3 on page 2 [ ]
%% Circle: "None, annot.type = <Circle>"
%% Comment: ""
%%
%⭣ ⭣ ⭣
$P_{KFP}$ and $\Delta_V$ (and their extensions to forms) which gives a profound motivation to %%
%⭡ ⭡ ⭡ END of annotation 7
the link between the two can be found in the works of J.-M. Bismut and his hypoelliptic Laplacian, we can mention for example \cite{Bi05,BiLe05} for an overview of these concepts. We can also mention the book \cite{HeNi06} which contains a deep study of both these operators.

One of the pioneer work studying $P_{KFP}$ is \cite{HeNi04} in which the authors determined estimates regarding the return to equilibrium for the semigroup $e^{-tP_{KFP}}$ for certain potentials that are homogeneous near infinity. This result has then been generalized for example in \cite{Li09,Be22}. In \cite{HeSjSt05} the authors derived a rough localization of the small eigenvalues of $P_{KFP}$ in the semiclassical limit $h\to0$. More precise asymptotics have been obtained in \cite{HeHiSj08-2,HeHiSj11_01}. Note that these work are not restricted to the KFP operator, the authors considered a wider class of supersymmetric non-self-adjoint operators. Then, this supersymmetric assumption has been relaxed in \cite{BoLePMi22}. The goal of this paper is to obtain similar result for a KFP operator with general coefficients that does not fit the assumptions of \cite{HeHiSj08-2,HeHiSj11_01,BoLePMi22}.


\subsection{Main results}
In this paper, we consider the following stochastic differential equation
\begin{equation}\label{eq:SDE}
\left\{
\begin{aligned}
&dx_t = \alpha(x_t,v_t)dt,\\
&dv_t = \beta(x_t,v_t)dt - 4\Sigma^T\Sigma v_tdt + \sqrt{2h}dB_t,
\end{aligned}
\right.
\end{equation}
where $\alpha$ and $\beta$ are smooth functions of both $x\in\R^d$ %%
%% Annotation 9, 1/6 on page 3 [ ]
%% Circle: "None, annot.type = <Circle>"
%% Comment: ""
%%
%⭣ ⭣ ⭣
and $v\in\R^{d'}$, %%
%⭡ ⭡ ⭡ END of annotation 9
and $\Sigma\in GL_{d'}(\R)$ is a fixed invertible matrix.

Its generator is
\begin{equation*}
\lll = \alpha\cdot\p_x + \beta\cdot\p_v - 4\Sigma^T\Sigma v\cdot\p_v + h\Delta_v.
\end{equation*}
Therefore we %%
%% Annotation 10, 2/6 on page 3 [ ]
%% Circle: "None, annot.type = <Circle>"
%% Comment: ""
%%
%⭣ ⭣ ⭣
have
\begin{equation*}
-\lll^* = \alpha\cdot\p_x + \beta\cdot\p_v + \div_x\alpha + \div_v\beta - (h\Delta_v + 4\Sigma^T\Sigma v\cdot\p_v + 4\Tr(\Sigma^T\Sigma)).
\end{equation*}

We %%
%⭡ ⭡ ⭡ END of annotation 10
notice that $(h\Delta_v + 4\Sigma^T\Sigma v\cdot\p_v + 4\Tr(\Sigma^T\Sigma))e^{-2|\Sigma v|^2/h} = 0$ and thus we make the following assumption
\begin{assumption}\label{ass.skewadj}
    There exists $V\in\ccc^\infty(\R^d,\R)$ such that denoting
    \begin{equation}\label{eq:f}
    f(x,v) = V(x) + |\Sigma v|^2,
    \end{equation}
    we %%
%% Annotation 11, 3/6 on page 3 [ ]
%% Circle: "None, annot.type = <Circle>"
%% Comment: ""
%%
%⭣ ⭣ ⭣
have
    \begin{equation*}
    \lll^*(e^{-2f/h}) = 0.
    \end{equation*}
\end{assumption}
In %%
%⭡ ⭡ ⭡ END of annotation 11
other words, this assumption %%
%% Annotation 12, 4/6 on page 3 [ ]
%% Circle: "None, annot.type = <Circle>"
%% Comment: ""
%%
%⭣ ⭣ ⭣
is
\begin{equation}\label{eq:assskew}
\alpha\cdot\p_xV + 2\beta\cdot\Sigma^T\Sigma v - \frac h2(\div_x\alpha+\div_v\beta) = 0.
\end{equation}
Noticing %%
%⭡ ⭡ ⭡ END of annotation 12
also that $\lll1=0$ we can consider $P = -e^{f/h}\circ h\lll^*\circ e^{-f/h}$, this way we formally have $P(e^{-f/h}) = P^*(e^{-f/h}) = 0$ and we have the explicit formula $P = X + N$, %%
%% Annotation 13, 5/6 on page 3 [ ]
%% Ink: "this wa<SEL>y w</SEL>e formally"
%% Comment: ""
%%
%% Annotation 14, 6/6 on page 3 [ ]
%% Highlight: "divv β) <SEL>= 0.</SEL>  f/h f/h"
%% Comment: "period at end"
%%
%⭣ ⭣ ⭣
with
\begin{equation}\label{eq:generalP}
\left\{
\begin{aligned}
&X=\alpha(x,v)\cdot h\partial_x + \beta(x,v)\cdot h\partial_v + \frac h2(\div_x\alpha+\div_v\beta),\\
&N=-h^2\Delta_v + 4|\Sigma^T\Sigma v|^2 - 2h\Tr(\Sigma^T\Sigma).
\end{aligned}
\right.
\end{equation}
Where %%
%⭡ ⭡ ⭡ END of annotations 13, 14
we denoted $\Delta_v$ the Laplacian acting on $v$ only. We observe that we have the formal algebraic relations:
\begin{equation}\label{eq:adj0}
X^* = -X,\;\;N^* = N.
\end{equation}

Considering the following operators
\begin{equation}\label{eq:dV}
\d_V = h\p_x + \p_x V
\end{equation}
and
\begin{equation}\label{eq:dSigmav}
\d_{|\Sigma v|^2} = h\p_v + 2\Sigma^T\Sigma v,
\end{equation}
we observe that thanks to Assumption \ref{ass.skewadj}, \eqref{eq:generalP} can be rewritten
\begin{equation}\label{eq:generalP2}
\left\{
\begin{aligned}
&P = X+N,\\
&X = \alpha\cdot \d_V + \beta\cdot\d_{|\Sigma v|^2},\\
&N = \d_{|\Sigma v|^2}^*\d_{|\Sigma v|^2}.
\end{aligned}
\right.
\end{equation}

We want to obtain similar precise asymptotics of the bottom of the spectrum of $P$ to the ones obtained in \cite{HeHiSj11_01,BoLePMi22}. One of the main building block of \cite{HeHiSj08-2} (upon which \cite{HeHiSj11_01,BoLePMi22} relies) is the determination of resolvent estimates. To that extent, \cite{HeHiSj08-2} crucially needs that the symbol of the operator they consider is locally quadratic. We can easily observe that for general $\alpha,\beta$ and non-Morse potential, this cannot be true.

However, the kinetic structure of $P$ is very important and yields powerful tools to palliate \cite{HeHiSj08-2}. We can make use of the notion of hypocoercivity. This notion, initiated in \cite{He06}, and developed %%
%% Annotation 15, 1/1 on page 4 [ ]
%% Highlight: "be true.  <SEL>However, the kinetic structure of P is very important and yields powerful tools to palliate [17].</SEL> We can"
%% Comment: ""
%%
%⭣ ⭣ ⭣
further in \cite{Vi09,CaDoHeMiMoSc24} and references therein. Here we are in a setting which is very convenient to use a semiclassical version of \cite{DoMoSc15} and we shall follow this strategy %%
%⭡ ⭡ ⭡ END of annotation 15
to obtain the desired resolvent estimates.

Throughout the paper, we shall consider a potential satisfying

\begin{assumption}\label{ass.confin}
There exist $C>0$ and a compact  set $K\subset\R^d$ such that
\begin{equation*}
V(x)\ \geq\  -C,\;\;\;
 \vert\nabla V(x)\vert \geq \frac 1 C\;\;\;\text{and}\;\;\; \Vert\Hess V(x)\Vert_\infty \leq C,
\end{equation*}
for all $x\in\R^d\setminus K$.
\end{assumption}

\begin{lemma}\label{lem:soussolaffine}
    Let $V$ satisfying Assumption \ref{ass.confin}, then there exists $b\in\R$ such that
    \begin{equation*}
    \forall x\in\R^d,\ \ V(x) \geq \frac1C|x| + b
    \end{equation*}
    with $C>0$ given by Assumption \ref{ass.confin}. 
\end{lemma}
With this lemma (which proof is postponed to Subsection \ref{sec:proofsoussolaffine}), we hence have $e^{-V/h}\in L^2(\R^d)$ and thus $\Ker\Delta_V = \C e^{-V/h}$ knowing that $\d_V = e^{-V/h}\circ h\p_x\circ\ e^{V/h}$. Then, thanks to Assumption \ref{ass.confin}, $e^{-f/h}\in L^2(\R^{d+d'})$ and $N(e^{-f/h}) = 0$. Moreover Assumption \ref{ass.skewadj} is equivalent to $X(e^{-f/h}) = 0$ and together with the last assumption, it ensures that $e^{-f/h}\in D(P)$, thus, the formal equality $P(e^{-f/h}) = P^*(e^{-f/h}) = 0$ is now true in $L^2(\R^{d+d'})$.

We also consider the following assumption.
\begin{assumption}\label{ass.accretive}
    There exist $C>0$ and a compact set $K\subset\R^{d+d'}$ such that outside $K$%%
%% Annotation 16, 1/5 on page 5 [ ]
%% Highlight: "Assumption 3. <SEL>There exist C > 0 and a compact set K ⊂Rd+d′</SEL> <SEL>such that outside K,</SEL>  | divx"
%% Comment: ""
%%
%⭣ ⭣ ⭣
,
    \begin{equation*}
    |\div_x \alpha + \div_v\beta|\leq Cf.
    \end{equation*}
\end{assumption}
It is essential to prove the following proposition.
\begin{proposition}\label{prop:existence}
    Suppose %%
%⭡ ⭡ ⭡ END of annotation 16
Assumptions \ref{ass.skewadj}, \ref{ass.confin}, and \ref{ass.accretive} hold. For any initial condition, the Cauchy problem associated with \eqref{eq:SDE} admits %%
%% Annotation 17, 2/5 on page 5 [ ]
%% Highlight: "Proposition 1.2. <SEL>Suppose Assumptions 1, 2, and 3 hold. For any initial condition, the Cauchy problem</SEL> associated with"
%% Comment: ""
%%
%⭣ ⭣ ⭣
a solution. It is almost surely unique and determines an absolutely continuous stochastic process for all time $t\geq0$.
\end{proposition}
The %%
%⭡ ⭡ ⭡ END of annotation 17
goal of this work is to study the long time behavior of solutions to \eqref{eq:SDE}. Therefore it is crucial to have a global existence result for this process. We postpone the proof of this proposition to Subsection \ref{ssec:existence}.

In the following we consider hypoelliptic operators. More precisely, we shall assume they satisfy the standard hypoellipticity theorem for Hörmander operators \cite[Theorem 1.1]{Ho67}
\begin{assumption}\label{ass.Hormander}
    The operators $Y_{i_1},\ [Y_{i_1},Y_{i_3}],\ [Y_{i_4},[Y_{i_5},[Y_{i_6},\ldots,Y_{i_j}]]]\ldots$ where the $Y_{i_k}$ are in $\{X,\p_{v_1},\ldots,\p_{v_{d'}}\}$ span the whole tangent space %%
%% Annotation 18, 3/5 on page 5 [ ]
%% Highlight: "Assumption 4. <SEL>The operators Yi1, [Yi1, Yi3], [Yi4, [Yi5, [Yi6, . . . , Yij]]] . . . where the Yik</SEL> are in"
%% Comment: ""
%%
%⭣ ⭣ ⭣
at any point of $\R^{d+d'}$.
\end{assumption}

\begin{proposition}\label{prop:accretive}
The operator $P$ initially defined on $\ccc^\infty_c(\R^{d+d'})$ admits a %%
%⭡ ⭡ ⭡ END of annotation 18
unique maximally accretive extension that we still denote by $(P,D(P))$.
\end{proposition}
We postpone the proof of this Proposition to Subsection \ref{ssec:accretive}.

We denote by $\uuu$ the set of critical points of $V$.
\begin{assumption}\label{ass.morse}
    For any critical point $x^*\in\uuu$, there exists a neighborhood $\vvv\ni x^*$, $(t_i^{x^*})_{1\leq i\leq d}\subset \R^{*}$, $(\nu_i^{x^*})_{1\leq i\leq d}\subset \N\setminus\{0,1\}$, a $\ccc^\infty$ change %%
%% Annotation 19, 4/5 on page 5 [ ]
%% Highlight: "Assumption 5. <SEL>For any critical point x∗∈U, there exists a neighbor- hood V ∋x∗, (tx∗  i )1≤i≤d ⊂R∗, (νx∗  i )1≤i≤d ⊂N \ {0, 1},</SEL> a C∞change"
%% Comment: ""
%%
%⭣ ⭣ ⭣
of variable $U^{x^*}$ defined on $\vvv$ such that $U^{x^*}(x^*)=x^*$, $U^{x^*}$ and $d_{x^*}U^{x^*}$ are invertible and
    \begin{equation}\label{eq:morse}
    \forall x\in\vvv,\ \ V\circ U^{x^*}(x) - V(x^*) = \sum_{i=1}^dt_i^{x^*}(x_i-x_i^*)^{\nu_i^{x^*}}.
    \end{equation}
\end{assumption}
\begin{remark}
    We %%
%⭡ ⭡ ⭡ END of annotation 19
notice from \eqref{eq:morse} that any $x^*\in\uuu$ is an isolated critical point, but Assumption \ref{ass.confin} implies that $\uuu\subset K$ which we %%
%% Annotation 20, 5/5 on page 5 [ ]
%% Highlight: "Remark 1.4. <SEL>We notice from (1.13) that any x∗∈U is an isolated critical point, but Assumption 2 implies that U ⊂K</SEL> which we"
%% Comment: ""
%%
%⭣ ⭣ ⭣
recall is compact, therefore the set $\uuu$ is finite. Furthermore, if $V$ is a Morse function, $V$ satisfies Assumption \ref{ass.morse} through the Morse Lemma with %%
%⭡ ⭡ ⭡ END of annotation 20
$\nu_i^{x^*} = 2$ for all $i$.
\end{remark}
\begin{remark}[\cite{De25_1}]\label{rem:relax}
    We can relax \eqref{eq:morse} to
    \begin{equation}\label{eq:morsealt}
    \forall x\in\vvv,\ \ V\circ U^{x^*}(x) - V(x^*) = \sum_{i=1}^dt_i^{x^*}(x_i-x_i^*)^{\nu_i^{x^*}}(1+r_i(x-x^*)),
    \end{equation}
    with $r_i(x) = O(x)$ smooth.
    Indeed, let us consider
    \begin{equation*}
    \phi : (x_1,\ldots,x_d)\to(x_1(1+r_1(x))^{1/\nu_1^{x^*}},\ldots,x_d(1+r_d(x))^{1/\nu_d^{x^*}}),
    \end{equation*}
    then $d_{0}\phi = \Id$ hence it is a $\ccc^\infty$-diffeomorphism in a neighborhood of $0$. Considering now $U^{x^*}$ that satisfies \eqref{eq:morsealt}, then $\hat U^{x^*} = U^{x^*}\circ\tau_{-x^*}\circ\phi^{-1}\circ\tau_{x^*}$ satisfies \eqref{eq:morse}, where $\tau_a(x) = x-a$ and moreover we %%
%% Annotation 24, 4/4 on page 6 [ ]
%% Highlight: "= Id <SEL>hence it is a C∞-diffeomorphism in a neighborhood of 0. Considering now U x∗</SEL>that satisfies"
%% Comment: ""
%%
%⭣ ⭣ ⭣
have $d_{x^*}\hat U^{x^*} = d_{x^*}U^{x^*}$.
\end{remark}
Most of the time when the reference is clear, we will just write $U,t_i$ and %%
%⭡ ⭡ ⭡ END of annotation 24
$\nu_i$ instead of $U^{x^*},t_i^{x^*}$ and $\nu_i^{x^*}$ in order to lighten the notations.

We denote $\uuu$ the set of critical points of $V$, and we consider the partition $\uuu = \uuu^{odd}\sqcup\uuu^{even}$ where $x^*\in\uuu^{even} \iff \forall i,\ \nu_i^{x^*}\in 2\N$. Now we shall say that $x^*\in\uuu^{even}$ is of index $j\in\lb0,d\rb$ if $\sharp\{i\ |\ t_i^{x^*}<0\}=j$ and therefore $\uuu^{even}=\bigsqcup_{j=0}^d\uuu^{(j)}$ where $\uuu^{(j)}$ is the set of critical points of $V$ with even order in each direction and of index $j$, we also denote $n_0=\sharp\uuu^{(0)}$ the number %%
%% Annotation 21, 1/4 on page 6 [ ]
%% Caret: "even =  U(j)<Caret> where U(j) "
%% Comment: ","
%%
%⭣ ⭣ ⭣
of minima of %%
%⭡ ⭡ ⭡ END of annotation 21
$V$.

Notice that the critical points of $f$ are exactly the $(x^*,0)$, with $x^*$ a critical point of $V$, Moreover we have the same partition $\uuu = \uuu^{odd}\sqcup\uuu^{even}$ and the critical points have the same index. In the following we will identify those two and use $x^*$ instead of $(x^*,0)$ where it is clear which one we are really talking about ($x^*$ will mostly be denoted either $\m$ if of index $0$ or $\s$ if of index $1$).

In the following we denote
\begin{equation}\label{eq:nubar}
\overline \nu = \max_{i,x^*}\nu_i^{x^*}
\end{equation}

We define the %%
%% Annotation 22, 2/4 on page 6 [ ]
%% Caret: "  i  i,x∗νx∗<Caret>  "
%% Comment: "."
%%
%% Annotation 23, 3/4 on page 6 [ ]
%% Circle: "None, annot.type = <Circle>"
%% Comment: "."
%%
%⭣ ⭣ ⭣
function $\rho(v) = (C_\Sigma h)^{-\frac{d'}{4}} e^{-\frac{|\Sigma v|^2}{h}}$, where $C_\Sigma > 0$ is a normalization constant so that $\vvert{\rho}_{L^2(\R^{d'}_v)} = 1$ (so %%
%⭡ ⭡ ⭡ END of annotations 22, 23
we have $C_\Sigma = \frac{\pi}{2}\big(\det\Sigma^{-1}\big)^{\frac{2}{d'}}$). We can then consider the crucial matrix that will intervene in all our work
\begin{equation}\label{eq:defG}
G = \<\alpha^T\alpha\rho,\rho\>_{L^2(\R^{d'}_v)} \in \ccc^\infty(\R^d,\Mr_d(\R)).
\end{equation} 

In the following, given two operators $A$ and $B$, we shall say that $A\lesssim B$ when for all $u$, $\<Au,u\>\lesssim\<Bu,u\>$, in the formal sense, without giving too much importance to the domains for now.


\begin{assumption}\label{ass.G}
    There exists $g_1,g_2>0$ such that for all $x\in\R^d$,
    \begin{itemize}
        \item[$i)$] $g_1(h)I_{d} \leq G(x) \leq g_2(h)I_{d}$,
        \item[$ii)$] $\forall i\in\lb1,d\rb,\ \p_{x_i}G(x) \lesssim G(x)$%%
%% Annotation 25, 1/1 on page 7 [ ]
%% Highlight: "Assumption 6. <SEL>There exists g1, g2 > 0 such that for all x ∈Rd,  i) g1(h)Id ≤G(x) ≤g2(h)Id,</SEL> ii) ∀i"
%% Comment: ""
%%
%⭣ ⭣ ⭣
,
        \item[$iii)$] $\forall i,j\in\lb1,d\rb,\ \<\alpha_i^2\alpha_j^2\rho,\rho\>_{L^2(\R_v^{d'})}(x) \leq g_2(h)^2$.
    \end{itemize}
\end{assumption}

With these assumptions, we can state a first result concerning the %%
%⭡ ⭡ ⭡ END of annotation 25
rough localization of the spectrum of $P$.

\begin{theorem}\label{thm:1}
Suppose Assumptions \ref{ass.skewadj}, \ref{ass.confin}, \ref{ass.morse}, \ref{ass.G} and \ref{ass.hypocoer} hold true. There exists $h_0>0$, $c_0,c_1,c>0$, such that for all $h\in]0,h_0]$, there exists $G_h$ subspace of $L^2(\R^{d+d'})$ of dimension $n_0$ such that
$$
\forall u\in D(P)\cap G_h^\bot\ \ \Vert (P-z)u\Vert_{L^2}\geq c_1g(h)\Vert u\Vert_{L^2}
$$
for every $z\in\C$ such that $|\Re z|\leq c_0g(h)$, with $g(h)$ defined in \eqref{eq:g}. Moreover, there exists an explicit constant $c_f>0$ depending only on $f$ such that if $g(h)$ satisfies
\begin{equation}\label{eq:gthm2}
g(h)\geq e^{-\frac{\tilde c}{2h}}\mbox{ for any } \tilde c<c_f,
\end{equation}
then there exists $\lambda_\m(h)\in\C$ for all $\m\in\uuu^{(0)}$ such that $\sigma(P)\cap\{|\Re z|\leq c_0g(h)\}=\{\lambda_\m(h),\m\in\uuu^{(0)}\}$ counted with multiplicity, and for all $\m\in\uuu^{(0)},\ |\lambda_\m(h)|\leq e^{-c/h}$. And (still under \eqref{eq:gthm2}), for all $0<c_0'<c_1$
\begin{equation*}
\forall |z|>c_0'g(h),\mbox{ such that }|\Re z|\leq c_0g(h),\ \vvert{(P-z)^{-1}}_{L^2}\leq\frac{2}{c_0'g(h)}.
\end{equation*}
In addition, under Assumptions \ref{ass.accretive} and \ref{ass.Hormander}, we have Proposition \ref{prop:accretive} and thus we can extend all these results for $\Re z \leq -c_0g(h)$.
\end{theorem}

\begin{remark}
    Observe from \eqref{eq:g} that if $\alpha$ and $\beta$ have low degeneracy leading to $g_1,g_2\sim h$ in Assumption \ref{ass.G}, then $\vvert{(P-z)^{-1}} = O(h^{\frac{4}{\overline{\nu}}-3})$ in the region $\{|z|>c_0'g(h)\}\cap\{\Re z\leq c_0g(h)\}$. We notice that this resolvent estimate is not the same as the one for the Witten Laplacian obtained in \cite{De25_1} which we recall is of order $h^{\frac{2}{\overline{\nu}}-2}$. Although the estimates are different, when taking $\overline{\nu} = 2$, that is $V$ is Morse, we obtain an $h^{-1}$ for both the Fokker-Planck operator and the Witten Laplacian (which we know is sharp), therefore, even if our result may not be optimal, it still remains relevant.
\end{remark}

For the rough localization of Theorem \ref{thm:1}, we only need vague constructions around the critical points of $V$ in order to have Proposition \ref{prop:gapWitten}. In order to obtain sharp asymptotics, we have to refine our constructions. It requires the introduction of the topological definitions we recall in Section \ref{sec:labeling}.

We recall the generic assumption \eqref{ass.gener} which is useful in order to lighten the result and the proof.
\begin{equation}\label{ass.gener}\tag{Gener}
\begin{array}{l}
    (\ast)\mbox{ for any }\m\in\uuu^{(0)},\m \mbox{ is the unique global minimum of } V_{|E(\m)}\\
    (\ast)\mbox{ for all }\m\neq\m'\in\uuu^{(0)},\j(\m)\cap\j(\m')=\emptyset.
\end{array}
\end{equation}
In particular, \eqref{ass.gener} implies that $V$ uniquely attains its global minimum at $\um$. This assumption allows us to avoid some heavy constructions regarding the set $\uuu$ and lighten the definition \ref{def:quasimodes} of the quasimodes. But it seems that its not a true obstruction and that we can pursue the computations without this assumption as described in \cite[Section 6]{BoLePMi22}, \cite{No24}, in the spirit of \cite{Mi19}.

This assumption comes from \cite[(1.7)]{BoGaKl05_01} and \cite[Assumption 3.8]{HeKlNi04_01} where they were hard to handle properly. Then it changed to become \cite[Hypothesis 5.1]{HeHiSj11_01} when finally reaching the form of \cite[Assumption 4]{LePMi20}.

One can show that \eqref{ass.gener} is weaker than \cite{BoGaKl05_01}'s, \cite{HeKlNi04_01}'s and \cite{HeHiSj11_01}'s assumptions. More precisely, they supposed that $\j(\m)$ is a singleton while here we have no restriction on the size of $\j(\m)$.





The hypotheses of Theorem \ref{thm:1} hold for very general coefficients $\alpha,\beta$ and potential $V$. Then, the objective is to refine this rough localization of the eigenvalues of $P$. To that end, we follow the method of \cite[Section 3]{BoLePMi22} which is an adaptation of the WKB method. However, we were not able to make it work under the hypotheses of Theorem \ref{thm:1}. In Section \ref{sec:geoloc}, we study several situations in which we are able to give results. These are summarized under Proposition \ref{prop:ell} which was the target of our study of the WKB method. To ease the study of the aforementioned situations, we shall assume that the potential $V$ is a Morse function.
\begin{assumption}\label{ass:vraimentMorse}
    The potential $V$ is a Morse function.
\end{assumption}
We can observe what behavior arises when this assumption is not satisfied for the Witten Laplacian in \cite{De25_1}. In addition, we also impose a stronger assumption than \eqref{eq:gthm2}.
\begin{assumption}\label{ass.polyg}
    The function $g$ defined in \eqref{eq:g} is no worse than polynomial. That is, there exists $h_0,c>0$ such that for all $h\in(0,h_0]$,
    \begin{equation*}
    g(h)\geq h^c.
    \end{equation*}
\end{assumption}
From \eqref{eq:g}, $g(h) = O(h)$, therefore we in fact have $c\geq1$.

\begin{theorem}\label{thm:2}
    Suppose the hypothesis of Theorem \ref{thm:1}, Assumptions \ref{ass:vraimentMorse}, \ref{ass.polyg}, and \eqref{ass.gener} hold true. Suppose moreover that Proposition \ref{prop:ell} apply. There exists $h_0,\gamma > 0$ such that for all $h\in\ ]0,h_0]$, one has $\lambda(\um,h) = 0$ and for all $\m\neq\um$, $\lambda(\m,h)$ satisfies the following Eyring–Kramers type formula
    \begin{equation*}
    \lambda(\m,h) = v(\m)h^{\mu(\m)}e^{-2S(\m)/h}(1+O(h^\gamma)),
    \end{equation*}
    where $v(\m)$ and $\mu(\m)$ are defined in \eqref{eq:defvmu} and depend explicitly on $V$, and $S$ is the standard height function defined by \eqref{eq:defS}.
\end{theorem}

As mentioned in the paragraph after \eqref{ass.gener}, it seems that this assumption is not required in order to obtain sharp results although we did no computation if it does not hold. For Assumption \ref{ass.polyg}, we observe that it is very commonly satisfied and all examples in this article satisfy it.

In \cite{De25_1}, we obtained a very similar result for the Witten Laplacian, but we can mention some differences on the prefactor $v(\m)$ and $\mu(\m)$. The difference between the prefactor $v(\m)h^{\mu(\m)}$ of the small eigenvalues of the standard KFP operator \eqref{eq:standardKFP} and the ones of the Witten Laplacian, both with Morse potential lies in the negative eigenvalue of some non-degenerate matrix (see for example \cite{BoLePMi22} for a formula in both cases). Here it is a little bit more subtle. The degeneracies in $\alpha$ and $\beta$ induce a degeneracy in the resolution of the equations obtained by the WKB method leading to another factor when doing the Laplace method. The exponent $\mu(\m)$ in \cite{De25_1} was completely determined by the degeneracy of the potential $V$. While here, we proved Theorem \ref{thm:2} only for Morse potentials and yet, we can have various powers of $h$. This is due to the fact that degeneracies on $\alpha$ and $\beta$ affect $\mu(\m)$.

Note that the Arrhenius law $\lim_{h\to0}h\ln\lambda(\m,h) = -2S(\m)$ remains the same as usual. It is coherent with \cite[Corollary 1.2.4]{NiSaWh24} in which the authors proved this law for the standard KFP operator with very general potential. It is reasonable to assume that this very robust law (continuous in $V$ for the $C^0$ topology) would hold for a wide class of coefficients $\alpha$ and $\beta$%%
%% Annotation 26, 1/1 on page 9 [ ]
%% Highlight: "that this <SEL>very robust law (continuous in V for the C0 topology)</SEL> would hold"
%% Comment: ""
%%
%⭣ ⭣ ⭣
.






\subsection{Metastability}
The results presented here, proved in Subsection \ref{ssec:proofcor}, are adaptations of \cite[Corollary 1.5, 1.6]{BoLePMi22} to our settings. In generality, a %%
%⭡ ⭡ ⭡ END of annotation 26
result of the form of Theorem \ref{thm:2} allows to write such corollaries without much assumptions.

From Proposition \ref{prop:accretive}, $P$ is maximally accretive, therefore, for all $u_0\in L^2(\R^d)$, the following Cauchy problem
\begin{equation}\label{eq:evolutionP}
\left\{
\begin{aligned}
    &h\partial_tu + Pu = 0,\\
    &u_{|t=0} = u_0,
\end{aligned}
\right.
\end{equation}
admits a unique solution $u\in\ccc^0([0,+\infty),L^2(\R^d))\cap\ccc^1((0,+\infty),L^2(\R^d))$ denoted $u(t) = e^{-tP/h}u_0$. Theorem \ref{thm:2} gives the following result on the long time behavior of that solution.

\begin{corollary}\label{cor:return}
    In the setting of Theorem \ref{thm:2}, there exist $C,\varepsilon>0$ %%
%% Annotation 28, 2/3 on page 10 [ ]
%% Highlight: "unique solution <SEL>u ∈C0([0, +∞), L2(Rd))∩C1((0, +∞), L2(Rd)) denoted u(t) = e−tP/hu0.</SEL> Theorem 2"
%% Comment: ""
%%
%⭣ ⭣ ⭣
such that, for all $u_0\in L^2(\R^d)$ and $h$ small enough, there exists $(u_{\m,n})_{\m,n}\subset\C$ such that the solution $u(t)$ of \eqref{eq:evolutionP} satisfies
    \begin{equation}\label{eq:return1}
    \forall t\geq0,\ \ \vvert{u(t) - \sum_{\m\in\uuu^{(0)}}\sum_{n=0}^{n_0-1}u_{\m,n}t^ne^{-\lambda(\m,h)t/h}} \leq Ce^{-\varepsilon h^{c-1}t}\vvert{u_0},
    \end{equation}
    with %%
%⭡ ⭡ ⭡ END of annotation 28
$c\geq1$ given by Assumption \ref{ass.polyg}. Moreover, there exists $C>0$ such that, for all $u_0\in L^2(\R^d)$ and $h$ small enough, the solution $u(t)$ of \eqref{eq:evolutionP} %%
%% Annotation 27, 1/3 on page 10 [ ]
%% PolyLine: "−  um,ntne−λ(m,h)t<SEL>/h  ≤Ce</SEL>−εhc−1t ∥u0∥,  n=0  m∈U(0)          <SEL>     </SEL>   with c"
%% Comment: "break"
%%
%⭣ ⭣ ⭣
satisfies
    \begin{equation}\label{eq:return2}
    \forall t\geq 0,\ \ \vvert{u(t) - \frac{\<e^{-f/h},u_0\>}{\vvert{e^{-f/h}}^2}e^{-f/h}} \leq Ce^{-t\underset{\m\neq\um}{\min}\Re(\lambda(\m,h))(1-Ch)/h}\vvert{u_0}.
    \end{equation}
\end{corollary}

Another %%
%⭡ ⭡ ⭡ END of annotation 27
way to write \eqref{eq:return1} is
\begin{equation*}
u(t) = e^{-tP/h}\Pi_\Cr u_0 + O(e^{-\varepsilon t})\vvert{u_0},
\end{equation*}
while \eqref{eq:return2} is
\begin{equation*}
u(t) = e^{-tP/h}\Pi_0 u_0 + O(e^{-t\underset{\m\neq\um}{\min}\Re(\lambda(\m,h))(1-Ch)/h})\vvert{u_0},
\end{equation*}
with the $O$ being uniform in $t$ and $h$. Here, $\Pi_0$ denotes the orthogonal projector on the kernel of $P$, and $\Pi_\Cr$ denotes the spectral projector of $P$ associated with its $n_0$ exponentially small eigenvalues (recalling $n_0$ is the number of minima of $V$). It is defined %%
%% Annotation 29, 3/3 on page 10 [ ]
%% Highlight: "and h. <SEL>Here, Π0 denotes the orthogonal projector on the kernel of P, and ΠC</SEL> denotes the"
%% Comment: ""
%%
%⭣ ⭣ ⭣
as
\begin{equation*}
\Pi_\Cr = \frac{1}{2i\pi}\int_{\Cr}(z-P)^{-1}dz,
\end{equation*}
where $\Cr = \partial D(0,\frac{c_0}2g(h))$, with $c_0$ and $g(h)$ given by Theorem \ref{thm:1}.

Furthermore, %%
%⭡ ⭡ ⭡ END of annotation 29
we can describe the metastable behavior of solutions of \eqref{eq:evolutionP}.

\begin{corollary}\label{cor:times}
    In the setting of Theorem \ref{thm:2}, let $S_1\leq\cdots\leq S_{p+1}=+\infty$ denote the non-decreasing sequence of the $S(\m)$'s defined in \eqref{eq:defS} such that if $S_k = S_{k+1}$, then $\mu_k \leq \mu_{k+1}$ and let $\Pi^\leq_k$ be the spectral projector of $P$ associated with its eigenvalues of modulus less than $h^{\mu_k}e^{-2S_k/h}$. For two positive functions $t_\pm(h)$ such that $t_-(h) = O(h^\infty)$ and $t_+^{-1}(h) = o(h|\ln h|^{-1})$, we define the %%
%% Annotation 31, 2/2 on page 11 [ ]
%% Highlight: "TEMPERATURE 11  <SEL>that if Sk = Sk+1, then µk ≤µk+1 and let Π≤  k be the spectral projector of P associated with</SEL> its eigenvalues"
%% Comment: ""
%%
%⭣ ⭣ ⭣
times
    \begin{equation*}
    t_0^+ = t_+(h)\ \ \text{ and }\forall1\leq k\leq p+1,\ t^\pm_k = t_\pm(h)e^{2S_k/h}
    \end{equation*}
    (in particular $t^-_{p+1} = +\infty$). Then, for every $h$ small enough, the solution $u(t)$ of \eqref{eq:evolutionP} satisfies
    \begin{equation*}
    \forall t^+_{k-1}\leq t\leq t^-_k,\ u(t) = \Pi^\leq_ku_0 + O(h^\infty)\vvert{u_0},
    \end{equation*}
    uniformly %%
%⭡ ⭡ ⭡ END of annotation 31
with respect to $t$, $1\leq k\leq p+1$, and $u_0\in L^2(\R^d)$.
\end{corollary}

In other words, $e^{-tP/h}$ is approximately constant equal to $\Pi^\leq_k$ on the time interval $[t^+_{k-1},t^-_k]$, with fast transition around the times $t_k = h^{\mu_k}e^{2S_k/h}\in(t^-_k,t^+_k)$. In this corollary, one can take $t_-(h) = e^{-\delta/h}$ for some $\delta>0$ and $t_+(h) = \frac{|\ln h|^2}{h}$.

We can link this with the stochastic process solving \eqref{eq:SDE}. Under mild hypotheses, one can show that the probability density $\rho(t,\cdot)$ of the process $(X_t)$ solution of \eqref{eq:SDE} is solution to the problem
\begin{equation}\label{eq:probadensity}
\partial_t \rho = \lll^*\rho
\end{equation}
Therefore, recalling that $P = -e^{f/h}\circ(h\lll^*)\circ e^{-f/h}$, $u$ is solution to \eqref{eq:evolutionP} if and only if $e^{-f/h}u$ is solution to \eqref{eq:probadensity} (with adapted initial conditions). We observe that Corollary \ref{cor:return} gives
\begin{equation*}
\forall t\geq 0,\ \ \vvert{\rho(t) - \frac{\<e^{-f/h},u_0\>}{\vvert{e^{-f/h}}^2}e^{-2f/h}}_{TV} \leq Ce^{-t\underset{\m\neq\um}{\min}\Re(\lambda(\m,h))(1-Ch)/h}\vvert{e^{-f/h}}\vvert{u_0},
\end{equation*}
using that for absolutely continuous measure $\mu,\nu$, $\vvert{\mu-\nu}_{TV} = \vvert{\mu-\nu}_1$ and the %%
%% Annotation 30, 1/2 on page 11 [ ]
%% PolyLine: "∀t ≥0,  <SEL>≤</SEL>Ce  m̸=m Re(λ(m,h))(1−Ch)/h   e−f/h   ∥u0∥,  ∥e−f/h∥2 e−2f/h  <SEL>TV</SEL>       ρ(t) −⟨e−f/h, u0⟩  <SEL>     </SEL>  using that"
%% Comment: "break"
%%
%⭣ ⭣ ⭣
Cauchy-Schwarz inequality.



\subsection*{Acknowledgements}
The author is grateful to Laurent Michel for his advice through this work and to Jean-François Bony %%
%⭡ ⭡ ⭡ END of annotation 30
for helpful discussions. This work is supported by the ANR project QuAMProcs 19-CE40-0010-01.

The rest of the paper is organized as follows. In the next section, we develop the hypocoercive estimates in order to have a rough localization of the eigenvalues of $P$ as well as resolvent estimates resulting in Theorem \ref{thm:1}. In Section \ref{sec:examples}, we give a short list of examples of coefficients $\alpha$ and $\beta$ for which Theorem \ref{thm:1} apply.

Then, we want to obtain precise Eyring-Kramers law for our operator. We did not manage to obtain such result in broad generality, but we could prove these in several prescribed situations. Section \ref{sec:geoloc} focuses on the local constructions of the WKB method in the spirit of \cite{BoLePMi22}, while the purpose of Section \ref{sec:global} is to glue these local constructions to obtain globally defined cutoffs. This leads to the proof of Theorem \ref{thm:2}.



\section{Hypocoercive estimates}\label{sec:Hypocoer}

Let $\chi_\m$, $\m\in \uuu^{(0)}$ be some cutoffs in $\ccc_c^\infty(\R^d)$ such that $\chi_\m$ is supported in $B(\m,r)$ for some $r>0$ to be chosen small enough and $\chi_\m=1$ near $\m$. We then %%
%% Annotation 32, 1/1 on page 12 [ ]
%% Highlight: "Hypocoercive estimates  <SEL>Let χm, m ∈U(0) be some cutoffs in C∞  c (Rd) such that χm is sup- ported in B(m, r)</SEL> for some"
%% Comment: ""
%%
%⭣ ⭣ ⭣
consider
\begin{equation*}
V_\m(x) = \chi_\m(x) e^{-(V(x) - V(\m))/h}
\end{equation*}
and their space
\begin{equation*}
E_h = \span\{V_\m,\ \m\in\uuu^{(0)}\}.
\end{equation*}
For %%
%⭡ ⭡ ⭡ END of annotation 32
$r$ small enough, the $V_\m$ have disjoint support hence $\dim E_h = n_0$.

\begin{proposition}\label{prop:gapWitten}\cite[Proposition 2.1]{De25_1}
    Recalling $\overline{\nu}=\max_{i,x^*}\nu_i^{x^*}$, we have
    \begin{equation*}
    \exists C>0,\ \forall u\in D(\Delta_V)\cap E_h^\bot\ \ \<\Delta_Vu,u\> \geq Ch^{2-\frac2{\overline{\nu}}}\vvert{u}^2.
    \end{equation*}
\end{proposition}

As the critical points of $f$ are the $(x^*,0)$ for $x^*\in\uuu$, with the same index (when it is defined), we will identify those two and use $x^*$ instead of $(x^*,0)$ where it is clear which one we are really talking about ($x^*$ will mostly be denoted either $\m$ if of index $0$ or $\s$ if of index $1$). We also define the global quasimodes
\begin{equation*}
f_\m(x,v)=\chi_\m(x) e^{-(f(x,v)-f(\m))/h} = V_\m(x)e^{-\frac{|\Sigma v|^2}{h}},
\end{equation*}
\begin{equation*}
F_h=\span\{f_\m,\;\m\in\uuu^{(0)}\}.
\end{equation*}
Both $E_h$ and $F_h$ have dimension $n_0$ but we must notice that $E_h\subset L^2(\R^d)$ while $F_h\subset L^2(\R^{d+d'})$.

We recall the function $\rho(v)=(C_\Sigma h)^{-\frac {d'}4}e^{-\frac{|\Sigma v|^2}{h}}$ and we introduce the projector onto the kernel of $N$ defined on $L^2(\R^{d+d'})$ by
\begin{equation*}
\Pi u(x,v)=\int_{\R^{d'}}u(x,v')\rho(v')dv' \rho(v)=u_\rho(x)\rho(v),
\end{equation*}
where we denoted
\begin{equation}\label{eq:rho}
u_\rho = \< u, \rho\>_{L^2(\R_v^{d'})}.
\end{equation}
We observe that $\Pi (E_h\otimes L^2(\R^{d'})) = F_h$.

\begin{lemma}\label{lem:XPi}
    Under Assumption \ref{ass.skewadj},
    \begin{equation}\label{eq:XPi}
    X\Pi = \alpha\cdot\d_V\Pi
    \end{equation}
    and hence recalling $G = \<\alpha\alpha^T\rho,\rho\>_{L^2(\R_v^{d'})} \in \ccc^\infty(\R^d,\Mr_d(\R))$, we obtain
    \begin{equation}\label{eq:XPistarXPi}
    (X\Pi)^*(X\Pi) = \d_V^*G\d_V\Pi.
    \end{equation}
\end{lemma}

\bp Using \eqref{eq:generalP2}, and having that $\Pi$ is a projector on the kernel of $\d_{|\Sigma v|^2}$, we immediately get that $X\Pi = \alpha\cdot\d_V\Pi$. We then obtain
\begin{equation*}
\begin{aligned}
    (X\Pi)^*(X\Pi) &= \Pi(\alpha\cdot\d_V)^*(\alpha\cdot\d_V)\Pi = \d_V^*\Pi\alpha\alpha^T\Pi\d_V\\
    \Pi\alpha\alpha^T\Pi u &= \Pi\alpha\alpha^Tu_\rho(x)\rho = u_\rho(x)\<\alpha\alpha^T\rho,\rho\>_{L^2(\R_v^{d'})}\rho = G\Pi u
\end{aligned}
\end{equation*}
and hence $(X\Pi)^*(X\Pi) = \d_V^*G\d_V\Pi$.

\ep

Under Assumption \ref{ass.G}, we have a straightforward corollary of Proposition \ref{prop:gapWitten},

\begin{corollary}\label{cor:lowbound}
    $\exists C>0,\ \forall u\in D(\Delta_V)\cap F_h^\bot\ \<\d_V^*G\d_Vu,u\> \geq Ch^{2-\frac2{\overline{\nu}}}g_1(h)\vvert{u}^2$.
\end{corollary}

We now define the following auxiliary operator
\begin{equation}\label{eq:A}
A = (h^{2-\frac2{\overline{\nu}}}g_1(h) + (X\Pi)^*(X\Pi))^{-1}(X\Pi)^* = (h^{2-\frac2{\overline{\nu}}}g_1(h) + \d_V^*G\d_V)^{-1}(X\Pi)^*.
\end{equation}
This auxiliary operator is introduced in \cite{DoMoSc15} and used %%
%% Annotation 34, 2/4 on page 13 [ ]
%% Caret: " G dV u, u⟩≥<Caret>Ch2−2  We no"
%% Comment: "\forcelinebreak{}"
%%
%⭣ ⭣ ⭣
in \cite{LeSaSt20} in %%
%⭡ ⭡ ⭡ END of annotation 34
order to ease the calculus in the proof of Theorem \ref{thm:1}. This kind of %%
%% Annotation 33, 1/4 on page 13 [ ]
%% Text: "None, annot.type = <Text>"
%% Comment: "adjust to fit"
%%
%⭣ ⭣ ⭣
method to %%
%⭡ ⭡ ⭡ END of annotation 33
compute hypocoercivity was mainly introduced and used at first in \cite{Vi09}, \cite{HeNi04} %%
%% Annotation 35, 3/4 on page 13 [ ]
%% Highlight: "ν g1(h)+(XΠ)∗(XΠ))−1(XΠ)∗<SEL>=</SEL> (h2−2  ν"
%% Comment: "align with first ="
%%
%⭣ ⭣ ⭣
and \cite{He06}.
\begin{lemma}\label{lem:Aborne}
The %%
%⭡ ⭡ ⭡ END of annotation 35
operator $A$ is bounded on $L^2(\R^{d+d'})$, it satisfies $A = \Pi A$ and one has the estimate
\begin{equation*}
\Vert A\Vert_{L^2}\leq O(h^{\frac1{\overline{\nu}}-1}g_1(h)^{-\frac12})
\end{equation*}
\end{lemma}
\bp The bound is easily seen using Lemmas \ref{lem:1} and \ref{lem:2}.

\ep

The following assumption will help us prove some bounds on $A$ for the hypocoercivity result

\begin{assumption}\label{ass.hypocoer}
    We consider coefficients $\alpha$ and $\beta$ such that
    \begin{itemize}
        \item[$i)$] $\forall (x,v)\in\R^{d+d'}\ \int \alpha(x,\sqrt{h}v)e^{-2|\Sigma v|^2}\dv = 0$
        \item[$ii)$] For all $q\in\{hJ_x\alpha\alpha,hJ_v\alpha\beta,hJ_v\alpha\Sigma^T\Sigma v,h^2\Delta_v\alpha\}$, for all $i,j\in\lb1,d\rb$
        \begin{equation*}
        \Pi q_iq_j\Pi \lesssim g_2(h)^2(|\nabla V|^2+h)\Pi,
        \end{equation*}
        where we denote $J_y\gamma$ the Jacobian matrix of $\gamma$ with respect %%
%% Annotation 36, 4/4 on page 13 [ ]
%% Highlight: "∈R  α(x,  <SEL>ii) For all q ∈{hJxαα, hJvαβ, hJvαΣTΣv, h2\Deltavα}, for all i, j ∈  [[ 1, d ]]</SEL>  ΠqiqjΠ ≲g2(h)2(|∇V"
%% Comment: ""
%%
%⭣ ⭣ ⭣
to the variable $y$ and $\Delta_v\alpha$ the vector $(\Delta_v\alpha_i)_i$.
    \end{itemize}
\end{assumption}

\begin{remark}
    Let us notice that $i)$ implies that $\Pi X\Pi = \Pi\alpha\d_V\Pi = 0$ and %%
%⭡ ⭡ ⭡ END of annotation 36
thus $A = A(1-\Pi)$.
\end{remark}

%%There are several ways to verify this assumption, an easy one is to have that $\alpha$ is odd in $v$ which is the case for the Kramers-Fokker-Planck operator where $\alpha(x,v) = v$. But this is not always true, for example, the adaptive Langevin dynamic \cite{De23} does not have an odd $\alpha$ but still verifies Assumption \ref{ass.odd}.

This leads to the intermediate Lemma

\begin{lemma}\label{lem:bounds}
    Under Assumption \ref{ass.hypocoer}, there exists $C,h_0>0$ such that for all $h\in]0,h_0]$, for all $u\in L^2(\R^{d+d'})$, one %%
%% Annotation 37, 1/1 on page 14 [ ]
%% Highlight: "Lemma 2.6. <SEL>Under Assumption 9, there exists C, h0 > 0 such that for all h ∈]0, h0], for all u ∈L2(Rd+d′), one has</SEL>  |⟨ ⟩|"
%% Comment: ""
%%
%⭣ ⭣ ⭣
has
    \begin{equation}\label{eq:maj1}
    |\<AX (1-\Pi)u,u\>|\leq Ch^{\frac{2}{\overline\nu}-1}\Big(\frac{g_2(h)}{g_1(h)}\Big)^{\frac32}\Vert \Pi u\Vert\, \Vert(1- \Pi) u\Vert,
    \end{equation}
    \begin{equation}\label{eq:maj2}
    |\<AN u,u\>|\leq Ch^{\frac{2}{\overline\nu}-1}\Big(\frac{g_2(h)}{g_1(h)}\Big)^{\frac32}\Vert \Pi u\Vert\, \Vert(1- \Pi) u\Vert,
    \end{equation}
    \begin{equation}\label{eq:maj3}
    |\<X u,A u\>|\leq C\Vert(1- \Pi) u\Vert^2.
    \end{equation}
\end{lemma}
\bp
Within this %%
%⭡ ⭡ ⭡ END of annotation 37
proof, $C$ will denote a positive constant that may only depends on the dimension $d$ and $\Sigma$ and can change from line to line. Let us denote
\begin{equation}\label{eq:defR}
R = (h^{2-\frac2{\overline{\nu}}}g_1(h) + \d_V^*G\d_V)^{-1},
\end{equation}
this will makes the computations clearer. This way, $A = R(X\Pi)^*$.

Let us start with the proof of \eqref{eq:maj1}. Since $A=\Pi A$, by the Cauchy-Schwarz inequality it is sufficient to show that the operator $AX$ (or equivalently its adjoint) is bounded on $L^2$. But $X^*A^* = -X^2\Pi R$ with
\begin{equation*}
X^2\Pi = X(\alpha\cdot\d_V)\Pi = (\alpha\cdot\d_V)^2\Pi + \beta\cdot\d_{|\Sigma v|^2}\alpha\cdot\d_V\Pi
\end{equation*}
thanks to Lemma \ref{lem:XPi} and \eqref{eq:generalP2}. Using that $\d_{|\Sigma v|^2}\Pi = 0$,
\begin{equation}\label{eq:X2Pi}
X^2\Pi = (\alpha\cdot\d_V)^2\Pi + h\beta\cdot\p_v(\alpha\cdot\d_V)\Pi = (\alpha\cdot\d_V)^2\Pi + hJ_v\alpha\beta\cdot\d_V\Pi.
\end{equation}
Moreover, denoting $\d_{V,i} = h\p_{x_i} + \p_{x_i}V$, we can write
\begin{equation}\label{eq:alphadv2}
\begin{aligned}
    (\alpha\cdot\d_V)^2 &= \sum_{i,j}\alpha_i\d_{V,i}\alpha_j\d_{V,j}\\
    &= \sum_{i,j}\alpha_i\alpha_j\d_{V,i}\d_{V,j} + \sum_{i,j}\alpha_i[\d_{V,i},\alpha_j]\d_{V,j}\\
    &= -\sum_{i,j}\alpha_i\alpha_j\d_{V,i}^*\d_{V,j} + 2\sum_{i,j}\alpha_i\alpha_j\p_{x_i}V\d_{V,j} + hJ_x\alpha\alpha\cdot\d_V.
\end{aligned}
\end{equation}

Using Assumption \ref{ass.G} $iii)$, for $i,j\in\lb1,d\rb$ and $u\in L^2(\R^{d+d'})$, we have
\begin{equation*}
\begin{aligned}
    \vvert{\alpha_i\alpha_j\d_{V,i}^*\d_{V,j}\Pi Ru}^2 &= \<\Pi \alpha_i^2\alpha_j^2\Pi\d_{V,i}^*\d_{V,j}Ru,\d_{V,i}^*\d_{V,j}Ru\>\\
    &\leq Cg_2(h)^2\vvert{\d_{V,i}^*\d_{V,j}Ru}^2\\
    &\leq C\big(h^{\frac{2}{\overline{\nu}} - 1}g_1(h)^{-3/2}g_2(h)^{3/2}\big)^2
\end{aligned}
\end{equation*}
with Lemma \ref{lem:4}. Using the same arguments, we obtain
\begin{equation*}
\vvert{\alpha_i\alpha_j\p_{x_i}V\d_{V,j}\Pi Ru} \lesssim g_2(h)\vvert{\p_{x_i}V\d_{V,j}\Pi Ru}.
\end{equation*}
Moreover, with Assumption \ref{ass.hypocoer}, we have
\begin{equation*}
\vvert{hJ_x\alpha\alpha\cdot\d_V\Pi Ru} \lesssim g_2(h)(\vvert{|\nabla V|\d_V\Pi Ru} + \sqrt{h}\vvert{\d_V\Pi Ru}),
\end{equation*}
hence there just remains to control terms of the form $\p_{x_i}V\d_{V,j}R$ using Lemma \ref{lem:3}. Noticing that the non-negativity of the Laplacian %%
%% Annotation 38, 1/3 on page 15 [ ]
%% Highlight: "ΠRu∥+  hence <SEL>there just remains to control terms of the form ∂xi V dV,j R using L</SEL>emma B.3."
%% Comment: ""
%%
%⭣ ⭣ ⭣
implies that $|\nabla V|^2\leq \Delta_V + h\Delta V$, we have
\begin{equation}\label{eq:nablaVleqDeltaV}
|\nabla V|^2\leq\Delta_V + C'h
\end{equation}
for %%
%⭡ ⭡ ⭡ END of annotation 38
some $C'>0$ using Assumption \ref{ass.confin}. Thus we just have to estimate
\begin{equation*}
\<\d_V^*\Delta_V\d_VRu,Ru\>.
\end{equation*}
Notice that the $\Delta_V$ appearing is actually a scalar matrix $\Delta_VI_d$ %%
%% Annotation 39, 2/3 on page 15 [ ]
%% Highlight: "\DeltaV Id <SEL>where Id denotes</SEL> the identity"
%% Comment: ""
%%
%⭣ ⭣ ⭣
where $I_d$ denotes the identity matrix of size $d$. Consider now the relation $\Delta_V^{(1)} = \Delta_VI_d + 2h\Hess V$ where $\Delta_V^{(1)}$ denotes the Witten Laplacian acting on $1$-forms (that we identify with $\R^d$-valued functions, we refer to %%
%⭡ ⭡ ⭡ END of annotation 39
\cite{HeSj85_01} for more details about Witten Laplacians on $p$-forms). Using Assumption \ref{ass.confin}, and the commutation rule (\cite[(2.1)]{HeKlNi04_01})
\begin{equation*}
\Delta_V^{(1)}\d_V = \d_V\Delta_V,
\end{equation*}
we obtain
\begin{equation*}
\begin{aligned}
    \<\d_V^*\Delta_V\d_VRu,Ru\> &= \<\d_V^*\Delta_V^{(1)}\d_VRu,Ru\> - 2h\<\d_V^*\Hess V\d_VRu,Ru\>\\
    &\leq \<\d_V^*\d_V\Delta_VRu,Ru\> + 2Ch\vvert{\d_VRu}^2\\
    &= \vvert{\Delta_VRu}^2 + 2Ch\vvert{\d_VRu}^2.
\end{aligned}
\end{equation*}
This shows that
\begin{equation}\label{eq:pVdVR}
\vvert{|\nabla V|\d_VR} \lesssim \vvert{\Delta_VR} + \sqrt{h}\vvert{\d_VR},
\end{equation}
therefore we obtain with
\begin{equation*}
\vvert{X^2\Pi R} \lesssim h^{\frac{2}{\overline{\nu}} - 1}g_1(h)^{-3/2}g_2(h)^{3/2} + g_2(h)\vvert{\Delta_V\Pi R} + \sqrt hg_2(h)\vvert{\d_V \Pi R}.
\end{equation*}
Hence \eqref{eq:maj1} using Lemmas \ref{lem:1}, \ref{lem:3} and \ref{lem:4}.

















































Now for \eqref{eq:maj2}, because $\Pi$ is the projection onto the kernel of $N$, %%
%% Annotation 40, 3/3 on page 15 [ ]
%% Replace: "by the <SEL>Cauchy- Schwarz</SEL> inequality, we"
%% Comment: "Cauchy--Schwarz"
%%
%⭣ ⭣ ⭣
we have $N = N(1-\Pi)$ and we recall that $A = \Pi A$. Hence by the Cauchy-Schwarz inequality, we only need a bound on $AN$ or $NA^*$ (since $N$ is self-adjoint). Recall %%
%⭡ ⭡ ⭡ END of annotation 40
that $X\Pi = \alpha\cdot\d_V\Pi$ and thus
\begin{equation*}
\begin{aligned}
    NA^* &= N\alpha\cdot\d_V\Pi R = [N,\alpha\cdot\d_V]\Pi R = [N,\alpha]\cdot\d_V\Pi R\\
    &= -h^2(\Delta_v\alpha + 2J_v\alpha\p_v)\Pi\cdot\d_V R\\
    &= (-h^2\Delta_v\alpha + 4hJ_v\alpha\Sigma^T\Sigma v)\cdot\d_V\Pi R
\end{aligned}
\end{equation*}
where we recall we denoted $\Delta_v\alpha$ the vector $(\Delta_v\alpha_i)_i$. Therefore thanks to Assumption \ref{ass.hypocoer} and \eqref{eq:pVdVR},
\begin{equation*}
\vvert{NA^*} \lesssim g_2(h)\big(\vvert{\Delta_V\Pi R} + \sqrt{h}\vvert{\d_V\Pi R}\big),
\end{equation*}
which proves \eqref{eq:maj2} using the same estimates as for $X^*A^*$.

And finally for \eqref{eq:maj3}, using that $A = \Pi A (1-\Pi)$ and $\Pi X \Pi = 0$ we quickly obtain
\begin{equation*}
\begin{aligned}
    \<Xu,Au\> &= \<Xu,\Pi A(1-\Pi)u\> = \<\Pi X(1-\Pi)u,A(1-\Pi)u\>\\
    &\leq \vvert{(\Pi X)^*A}\vvert{(1-\Pi)u}^2\\
    &= \vvert{X\Pi R(X\Pi)^*}\vvert{(1-\Pi)u}^2
\end{aligned}
\end{equation*}
And we recognize the norm of $QQ^*$ with $Q = X\Pi R^{1/2}$ which is bounded by Lemmas \ref{lem:XPi} and \ref{lem:2}, hence the result.

\ep






\begin{proposition}\label{prop:hypocoer}
There exists $C,\delta_0,h_0>0$ such that for all $h\in]0,h_0]$, and for all $u\in D(P)\cap F_h^\bot$, one has
$$
\Re\,\<Pu,(1+\delta(h)(A+A^*))u\>\geq C\delta(h) \Vert u\Vert^2,
$$
where $\delta(h)=\delta_0 \frac{h}{1 + h^{\frac{4}{\overline\nu}-2}\Big(\frac{g_2(h)}{g_1(h)}\Big)^{3}}$.
\end{proposition}
\bp
For all $\delta>0$ and $u\in D(P)\cap F_h^\bot$, let us define 
\begin{equation*}
I_\delta = \Re\,\<P u,(1+\delta(A+A^*)u\>
\end{equation*}
Using the decomposition $P = X + N$, and the skew-adjointness of $X$ coming from \eqref{eq:adj0}, one gets
\begin{equation*}
I_\delta = \<N u,u\> + \delta\Re\<Pu,(A+A^*)u\>
\end{equation*}
From the spectral properties of $N$, it follows that
\begin{equation}\label{eq:coer1}
I_\delta \geq h\Vert (1-\Pi)u\Vert^2 + \delta\Re\<Pu,(A+A^*)u\>.
\end{equation}
Denoting $J = \<Pu,(A+A^*)u\>$, one has
\begin{equation*}
J = \<AXu,u\> + \<ANu,u\> + \<Xu,Au\> + \<Nu,Au\>
\end{equation*}
and since $A = \Pi A$ and $\Pi N = 0$ it follows that 
\begin{equation}\label{eq:J}
J = \<AX\Pi u,u\> + J'
\end{equation}
with 
\begin{equation}\label{eq:J'}
J' = \<AX(1-\Pi)u,u\> + \<ANu,u\> + \<Xu,Au\>.
\end{equation}
Moreover, by definition of $A$ and Lemma \ref{lem:XPi}, on $F_h^\bot$
\begin{equation}\label{eq:AXPi}
AX\Pi = (h^{2-\frac2{\overline{\nu}}}g_1(h) + \d_V^*G\d_V)^{-1}\d_V^*G\d_V\Pi \geq c_0\Pi
\end{equation}
for some $c_0>0$ using functional calculus and Corollary \ref{cor:lowbound}. Therefore, combining \eqref{eq:coer1}, \eqref{eq:J}, \eqref{eq:J'}, \eqref{eq:AXPi} and Lemma \ref{lem:bounds}, we obtain
\begin{equation*}
\begin{aligned}
    \forall u\in D(P)\cap F_h^\bot,\ \ I_\delta &\geq h\vvert{(1-\Pi)u}^2 + \delta c_0\vvert{\Pi u}^2 - C\delta\vvert{(1-\Pi)u}^2\\&\phantom{****} - C\delta h^{\frac{2}{\overline\nu}-1}\Big(\frac{g_2(h)}{g_1(h)}\Big)^{\frac32}\Vert \Pi u\Vert\, \Vert(1- \Pi) u\Vert,
\end{aligned}
\end{equation*}
thus with Young's inequality, we have
\begin{equation*}
I_\delta \geq \Big(h-C\delta-\frac{C^2}{2c_0}\delta h^{\frac{4}{\overline\nu}-2}\Big(\frac{g_2(h)}{g_1(h)}\Big)^{3}\Big)\vvert{(1-\Pi)u}^2 + \delta\frac{c_0}{2}\vvert{\Pi u}^2.
\end{equation*}
Optimizing the right hand side by taking
\begin{equation*}
\delta = \frac{2c_0h}{c_0^2 + 2c_0C + C^2h^{\frac{4}{\overline\nu}-2}\Big(\frac{g_2(h)}{g_1(h)}\Big)^{3}},
\end{equation*}
we obtain
\begin{equation*}
I_\delta \geq \delta\frac{c_0}{2}\vvert{u}^2.
\end{equation*}
Noticing that there exists $C>0$ such that for $h$ small enough,
\begin{equation*}
\frac1C\delta(h)\leq\frac{h}{1 + h^{\frac{4}{\overline\nu}-2}\Big(\frac{g_2(h)}{g_1(h)}\Big)^{3}}\leq C\delta(h),
\end{equation*}
we proved the proposition.

\ep










\subsection{Proof of Theorem \ref{thm:1}}
Let $z\in\C$, $u\in D(P)\cap F_h^\bot$ and $\delta(h)$ as in Proposition \ref{prop:hypocoer}. First with the Cauchy-Schwarz inequality, we have
\begin{equation}\label{eq:thm1}
\Re\< (P-z)u, (1+\delta(h)(A+A^*))u \> \leq \vvert{(P-z)u}\vvert{1+\delta(h)(A+A^*)}\vvert{u},
\end{equation}
then thanks to %%
%% Annotation 42, 2/3 on page 17 [ ]
%% PolyLine: "Re⟨(P−z)u, (1+δ(h)(A+A∗)<SEL>)u⟩≥</SEL>Cδ(h) ∥u∥2−Re(z⟨u,"
%% Comment: "break"
%%
%⭣ ⭣ ⭣
Proposition \ref{prop:hypocoer}, 
\begin{equation*}
\Re\< (P-z)u, (1+\delta(h)(A+A^*))u \> \geq C\delta(h)\vvert{u}^2-\Re(z\< u,(1+\delta(h)(A+A^*))u\>).
\end{equation*}
Because %%
%⭡ ⭡ ⭡ END of annotation 42
$(1+\delta(h)(A+A^*))$ is symmetric, this leads %%
%% Annotation 41, 1/3 on page 17 [ ]
%% Text: "None, annot.type = <Text>"
%% Comment: "adjust to fit"
%%
%⭣ ⭣ ⭣
to
\begin{equation*}
\Re\< (P-z)u, (1+\delta(h)(A+A^*))u \> \geq C\delta(h)\vvert{u}^2-|\Re z|\vvert{u}^2\vvert{1+\delta(h)(A+A^*)}.
\end{equation*}
Using %%
%⭡ ⭡ ⭡ END of annotation 41
%%
%% Annotation 43, 3/3 on page 17 [ ]
%% PolyLine: "Re⟨(P−z)u, (1+δ(h)(A+A∗))<SEL>u⟩≥</SEL>Cδ(h) ∥u∥2−|"
%% Comment: ""
%%
%⭣ ⭣ ⭣
that $\vvert{1+\delta(h)(A+A^*)} \leq 1 + 2\delta(h)\vvert{A}$, %%
%⭡ ⭡ ⭡ END of annotation 43
this and \eqref{eq:thm1} lead to
\begin{equation*}
\vvert{(P-z)u} \geq C\frac{\delta(h)}{1 + 2\delta(h)\vvert{A}}\vvert{u} - |\Re z|\vvert{u}.
\end{equation*}
With the expression of $\delta(h)$ and Lemma \ref{lem:Aborne}, let us denote for $h>0$
\begin{equation}\label{eq:g}
g(h) = \frac{h}{1 + h^{\frac{4}{\overline\nu}-2}\Big(\frac{g_2(h)}{g_1(h)}\Big)^{3} + h^{\frac1{\overline\nu}}g_1(h)^{-\frac12}}.
\end{equation}
We notice that $g(h) = O(h)$. We then obtain that there exists $c_0,c_1>0$ such that for $|\Re z| \leq c_0g(h)$,\begin{equation}\label{eq:hypoPi}
\forall u\in D(P)\cap F_h^\bot,\ \vvert{(P-z)u} \geq c_1g(h)\vvert{u}.
\end{equation}

And we can now deduce the second part of Theorem \ref{thm:1} from that, following the same sketch of proof as in \cite{No23}.

Let $\m\in\uuu^{(0)}$, by recalling $f_\m(x,v)=\chi_\m(x) e^{-(f(x,v)-f(\m))/h}$, and because $e^{-f/h}\in\Ker P$, we obtain
\begin{equation*}
P(f_\m)= [P,\chi_\m]e^{-(f-f(\m))/h} = h\alpha\cdot\nabla\chi_\m e^{-(f-f(\m))/h} = O(e^{-c_\m/h})
\end{equation*}
with $c_\m=\inf_{\supp\nabla\chi_\m} f-f(\m)>0$ (because $\chi\equiv1$ near $\m$).

Moreover, with a change of variable $(x,v)\mapsto(h^{\frac{1}{\nu_1}}x_1,\ldots,h^{\frac{1}{\nu_d}}x_d,h^{\frac12}v)$,
\begin{equation*}
\vvert{f_\m} = Ch^{\frac{d'}{4} + \sum_{i=1}^d\frac{1}{2\nu_i^\m}}(1+O(h^\frac{1}{\overline\nu}))
\end{equation*}
for some $C>0$. Thus, since the $(f_\m)_{\m\in\uuu^{(0)}}$ are orthogonal, we actually have :
\begin{equation}\label{eq:Pfm}
\forall u\in F_h,\ \vvert{Pu}=O(e^{-c/h})\vvert{u}
\end{equation}
for all $c<c_f$ where $c_f=\min_{\m\in\uuu^{(0)}}c_\m>0$. Furthermore, \eqref{eq:Pfm} is also true replacing $P$ by $P^*$ because $P^*(f_\m)=-P(f_\m)$. And because
\begin{equation*}
P^*Pf_\m = -X(Pf_\m) = (\alpha\cdot h\p_x + \beta\cdot h\p_v)(h\alpha\cdot\nabla\chi_\m)e^{-(f-f(\m))/h},
\end{equation*}
\eqref{eq:Pfm} is still valid replacing $P$ by $P^*P$.

We denote by $\Pi_F$ the projector on $F_h$. Let $0<c_0'\leq c_1$, $u\in D(P)$ and $z$ such that $|\Re z|\leq c_0g(h)$ and $|z|\geq c_0'g(h)$,
\begin{equation*}
\begin{aligned}
    \vvert{(P-z)u}^2 &= \vvert{(P-z)(\Pi_F+\Id-\Pi_F)u}^2\\
    &= \vvert{(P-z)(\Id-\Pi_F)u}^2+\vvert{(P-z)\Pi_F u}^2\\&\phantom{*****}+2\Re\<(P-z)(\Id-\Pi_F)u,(P-z)\Pi_F u\>,
\end{aligned}
\end{equation*}
but one has
\begin{equation*}
\vvert{(P-z)(\Id-\Pi_F)u} \geq c_1g(h)\vvert{(\Id-\Pi_F)u}
\end{equation*}
thanks to \eqref{eq:hypoPi}, and 
\begin{equation*}
\vvert{(P-z)\Pi_F u}^2 \geq (\vvert{P\Pi_F u}-\vvert{z\Pi_F u})^2 \geq \vvert{z\Pi_F u}(\vvert{z\Pi_F u}-2\vvert{P\Pi_F u}).
\end{equation*}
Assume now that $g(h)\geq e^{-c/(2h)}$ for some $c<c_f$ (this is \eqref{eq:gthm2}),
\begin{equation*}
|z|\geq c_0'g(h)\geq c_0'e^{-c/(2h)}\geq c_0'e^{-c/h},
\end{equation*}
thus using \eqref{eq:Pfm}, we get
\begin{equation*}
\vvert{(P-z)\Pi_F u}^2 \geq \frac{|z|^2}{2}\vvert{\Pi_F u}^2.
\end{equation*}
Studying each term in the scalar product, there exists $c>0$ such that
\begin{equation*}
\begin{aligned}
    (\ast) :&= \Re\<(P-z)(\Id-\Pi_F)u,(P-z)\Pi_F u\>\\
    &= \Re\big(\<P(\Id-\Pi_F)u,P\Pi_F u\>-z\<(\Id-\Pi_F)u,P\Pi_F u\> - \bar{z}\<P(\Id-\Pi_F)u,\Pi_F u\>\big)\\
    &= (1+|z|)\vvert{(\Id-\Pi_F)u}\vvert{\Pi_F u}O(e^{-c/h})\\
    &= \Big(\vvert{u}^2+|z|^2\vvert{\Pi_F u}^2+\vvert{(\Id-\Pi_F)u}^2\Big)O(e^{-c/h})
\end{aligned}
\end{equation*}
hence
\begin{equation*}
\begin{aligned}
    \vvert{(P-z)u}^2 &\geq (c_1g(h))^2\vvert{(\Id-\Pi_F)u}^2 + \frac{|z|^2}3\vvert{\Pi_F u}^2\\
    &\phantom{*******}+(\vvert{u}^2+\vvert{(\Id-\Pi_F)u}^2)O(e^{-c/h})\\
    &\geq \frac13(c_0'g(h))^2\vvert{u}^2+(\vvert{u}^2+\vvert{(\Id-\Pi_F)u}^2)O(e^{-c/h})\\
    &\geq \frac14(c_0'g(h))^2\vvert{u}^2
\end{aligned}
\end{equation*}
for $h$ small enough, using \eqref{eq:gthm2}. It leads to
\begin{equation}\label{eq:injectivePz}
\forall u\in D(P),\quad\vvert{(P-z)u} \geq \frac{c_0'}{2}g(h)\vvert{u}.
\end{equation}

By using the same arguments for $P^*$ we have the same result for it (the key point is that $e^{-f/h}$ is in the kernel of $X$ and $N$ hence it also is in $P^*$'s one). It just remains to show that $P-z$ is surjective in order to obtain the resolvent estimate, we show it the classical way, by showing that $\Ran(P-z)$ is closed and dense.

Let $u_n\in D(P)$ and $w\in L^2$ such that $(P-z)u_n\to w$ therefore $((P-z)u_n)_{n\in\N}$ is Cauchy and so is $(u_n)_{n\in\N}$ thanks to \eqref{eq:injectivePz}, hence there exists $u\in L^2$ such that $u_n\to u$. Because the convergence is also true in $\ddd'$, $(P-z)u = w$ in $\ddd'$, and since $w\in L^2$, so is $(P-z)u$, thus $u\in D(P)$ and $\Ran(P-z)$ is closed. Now to show that $\Ran(P-z)$ is dense, we use \eqref{eq:injectivePz} for $P^*$ and so $\Ker(P^*-\overline{z})=\{0\}$.

All this leads to the resolvent estimate
\begin{equation}\label{eq:resol}
\vvert{(P-z)^{-1}}\leq\frac{2}{c_0'g(h)}.
\end{equation}
Hence, $P$ has no spectrum in
\begin{equation*}
\{|\Re z|\leq c_0g(h)\}\cap\{|z|\geq c_0'g(h)\}.
\end{equation*}
Suppose now that Assumption \ref{ass.Hormander} and \ref{ass.accretive} hold true, from Proposition \ref{prop:accretive} we know that $P$ is maximally accretive and therefore $P - z$ is invertible for all $\Re z < 0$. Moreover we easily see that
\begin{equation*}
\vvert{(P-z)u}\vvert{u} \geq \Re \<(P-z)u,u\> \geq -\Re z\vvert{u}^2,
\end{equation*}
thus for all $\Re z < 0$,
\begin{equation*}
\vvert{(P-z)^{-1}} \leq \frac{1}{-\Re z}.
\end{equation*}
which extends \eqref{eq:resol}
\begin{equation*}
\forall z\in\{\Re z\leq c_0g(h)\}\cap\{|z|\geq c_0'g(h)\},\ \ \vvert{(P-z)^{-1}}\leq\frac{2}{c_0'g(h)}.
\end{equation*}

There is left to show that the spectrum within $\{|z|\leq c_0'g(h)\}$ is composed of $n_0$ eigenvalues exponentially small compared to $h^{-1}$. By denoting $D =D(0,c_0' g(h))$ the disk in $\C$ centered at $0$ of radius $c_0' g(h)$, we consider the spectral projector
\begin{equation*}
\Pi_D=\frac{1}{2i\pi}\int_{\p\negthinspace D}(z-P)^{-1}dz
\end{equation*}
the projector on the small eigenvalues. We start by proving the following lemma
\begin{lemma}\label{lem:Pi_0}
There exists $C>0$ such that $\vvert{P\Pi_D}\leq Cc_0'g(h)$.
\end{lemma}

\bp
\begin{equation*}
P\Pi_D = \frac{1}{2i\pi}\int_{\p\negthinspace D}P(z-P)^{-1}dz = \frac{1}{2i\pi}\int_{\p\negthinspace D}z(z-P)^{-1}dz,
\end{equation*}
hence
\begin{equation*}
\vvert{P\Pi_D}\leq C(c_0' g(h))^2\frac{2}{c_0'g(h)}
\end{equation*}
thanks to \eqref{eq:resol}.

\ep

Let us now prove that $\dim\Ran\Pi_0 = n_0$. We first show that $\dim\Ran\Pi_0\leq n_0$. By contradiction, let us suppose $F_h^\bot\cap\Ran\Pi_D\neq\emptyset$ and so let us take $u\in F_h^\bot\cap\Ran\Pi_D$ of norm one. Since $u\in\Ran\Pi_D$, by Lemma \ref{lem:Pi_0}, $\vvert{Pu}\leq Cc_0' g(h)$, but because $u\in F_h^\bot$ we can use \eqref{eq:hypoPi} and so $\vvert{Pu}\geq c_1g(h)$. Taking $c_0'$ low enough, we have the contradiction we aimed for and thus, $\dim\Ran\Pi_D\leq n_0$.

For the converse inequality, we have
\begin{equation*}
\Pi_D-\Id=\frac{1}{2i\pi}\int_{\p\negthinspace D}z^{-1}(z-P)^{-1}Pdz
\end{equation*}
and therefore
\begin{equation}\label{quasiortho}
\begin{aligned}
    \eps_\m&=\Pi_Df_\m-f_\m=\frac{1}{2i\pi}\int_{\p\negthinspace D}(z-P)^{-1}P(f_\m)\frac{dz}z\\
    &= O(g(h)^{-1}e^{-c/h})=O(e^{-\frac{c}{2h}})
\end{aligned}
\end{equation}
for some $c>0$, using \eqref{eq:resol}, \eqref{eq:Pfm} and the hypothesis \eqref{eq:gthm2}.

Let us suppose $\sum_{\m\in\uuu^{(0)}} a_\m\Pi_D f_\m = 0$ with $\sum_{\m\in\uuu^{(0)}}  |a_\m|^2 = 1$, since $\Pi_D f_\m = f_\m + \eps_\m$, we have for all $\m'\in\uuu^{(0)}\sum_{\m\in\uuu^{(0)}}  a_\m(\delta_{\m,\m'} + \<\eps_\m,f_{\m'}\>) = 0$ and thus for all $\m, a_\m = O(e^{-c/h})$ for some $c>0$, which is in contradiction with $\sum |a_\m|^2 = 1$. We deduce that $\dim\Ran\Pi_D\geq n_0$ and hence $\dim\Ran\Pi_D = n_0$.

This leads to
\begin{equation*}
\sigma(P)\cap\{|\Re z|\leq c_0g(h)\}=\{\lambda_\m(h),\m\in\uuu^{(0)}\}\subset D(0,\frac{c_1}{2} g(h)).
\end{equation*}

It only remains to show that $\lambda_\m(h) = O(e^{-c/h})$. Noticing that $\Ran\Pi_D$ is $P$-stable and that $(\Pi_Df_\m)_{\m\in\uuu^{(0)}}$ is one of its basis, there exists $C>0$,
\begin{equation*}
\vvert{P\Pi_D f_\m} = \vvert{\Pi_D Pf_\m} \leq C\vvert{Pf_\m} = O(e^{-c/h}).
\end{equation*}
Having that $\vvert{f_\m} = 1 + o(1)$, this yields $P_{|\Ran\Pi_D} = O(e^{-c/h})$m hence $\sigma(P_{|\Ran\Pi_D})\subset D(0,e^{-c/h})$ for some $c>0$.








































\section{Examples of hypocoercive operators}\label{sec:examples}

\subsection{Example 1 : a generalization of the adaptive Langevin dynamics.}\label{ssec:exampleAdaptive}

We can try to generalize the process considered in \cite{LeSaSt20}. The starting point is to model a lack of knowledge on the gradient of $V$ by a drift proportional to another random process, in other words, we consider the following SDE
\begin{equation}\label{eq:SDEAdaptLang}
\left\{
\begin{aligned}
&dx'_t = 2\Sigma^T\Sigma v_tdt,\\
&dv_t = -\p_{x'}V(x'_t)dt - M_t\Sigma^T\Sigma v_tdt - 4\Sigma^T\Sigma v_tdt + \sqrt{2h}dB_t,\\
&dM_t = \Gamma(x'_t,M_t,v_t)dt.
\end{aligned}
\right.
\end{equation}
Where $x'_t,v_t\in\R^d$, $M_t\in\Mr_d(\R)$ is the new variable and $\Gamma$ is to be set so that the system admits an invariant probability measure with the same form as the standard Langevin process. Therefore, denoting $f(x',v,M) = V(x') + |\Sigma v|^2 + W(M)$ with $W$ another smooth real valued function to be determined, having Assumption \ref{ass.skewadj} satisfied, that is $-h\lll^*(e^{-2f/h}) = 0$ is equivalent to \eqref{eq:assskew}, which in this case is
\begin{equation}\label{eq:skewAdapt1}
4\<M\Sigma^T\Sigma v,\Sigma^T\Sigma v\> - h\div_v(M\Sigma^T\Sigma v) -\Gamma\cdot\p_M W + h\div_M\Gamma = 0,
\end{equation}
where
\begin{equation*}
\Gamma \cdot \p_MW = \sum_{i,j}\Gamma_{i,j}\p_{M_{i,j}}W \text{ and likewise, }\div_M\Gamma = \sum_{i,j}\p_{M_{i,j}}\Gamma_{i,j}.
\end{equation*}
To have a more workable framework, we can consider the case where there exists a fixed real unitary matrix $Q$ such that $Q^T\Gamma Q$ is diagonal. Let us denote $\gamma_i$ such that the $i$-th diagonal block of $Q^T\Gamma Q$ is $\gamma_iI_{r_i}$, $r_i\in\N^*$ such that $\sum r_i = d$. In the following if $A$ is a matrix, then $\diag(A) = (A_{i,i})_i$ and if $u$ is a vector, then $\diag(u) = (u_i\delta_{i,j})_{i,j}$ using $\delta$ the Kronecker symbol. In the next, we will use $\diag$ to switch from vector to matrix and vice versa. Thus we denote $\gamma = \diag(Q^T\Gamma Q)$. This special form for $\Gamma$ induces the same for $M$, let us denote $D$ diagonal by block, which $i$-th block is $y_iI_{r_i}$, such that $M = QDQ^T$. Therefore, the $y_i$ are the new variables. That way, we denote $y = \diag(D)$.

Thus, from the last equation of \eqref{eq:SDEAdaptLang} and the equality $M = QDQ^T$, $\gamma$ has the same form as $D$ and denoting $y$ the new variable,
\begin{equation*}
dy_{i,t} = \gamma_i(x'_t,QD_tQ^T,v_t)dt.
\end{equation*}
Assuming $\Gamma$ does not depend on $y$, then \eqref{eq:skewAdapt1} becomes
\begin{equation}\label{eq:skewAdapt2}
4\<QDQ^T\Sigma^T\Sigma v,\Sigma^T\Sigma v\> - h\Tr(QDQ^T\Sigma^T\Sigma) - \sum_i\<\gamma_i, \p_{y_i}W\> = 0.
\end{equation}
Now, a natural $W$ to consider is $W(y) = |y|^2$, hence $\p_{y_i}W = 2y_i$.

Noticing now that for any $d\times d$ matrix $A$ and vector $u$ of size $d$ we have
\begin{equation*}
\Tr(A\diag(u)) = \<\diag(A),u\>,
\end{equation*}
thus $\Tr(QDQ^T\Sigma^T\Sigma) = \<\diag(Q^T\Sigma^T\Sigma Q),y\>$. Moreover, using that for any vector $u$, $Du = \diag(u) y$, we obtain
\begin{equation*}
\<QDQ^T\Sigma^T\Sigma v,\Sigma^T\Sigma v\> = \<\diag(Q^T\Sigma^T\Sigma v)y,Q^T\Sigma^T\Sigma v\>.
\end{equation*}

From \eqref{eq:skewAdapt2}, we thus obtain
\begin{equation}\label{eq:skewAdapt3}
\<4\diag((\Sigma Q)^T\Sigma v)(\Sigma Q)^T\Sigma v - h\diag((\Sigma Q)^T\Sigma Q), y\> = 2\sum_i\<\gamma_i, y_i\>.
\end{equation}

We can consider the case where $Q = I_d$ to simplify the expression of $\gamma$. Because \eqref{eq:skewAdapt3} must be true for any $y$, we get
\begin{equation*}
\gamma_i = 2\sum_{k = 1+\sum_{n\leq i-1}r_{n}}^{\sum_{n\leq i}r_{n}}(\sum_{j=1}^d(\Sigma^T\Sigma)_{k,j}v_j)^2 - \frac h2\sum_{k = 1+\sum_{n\leq i-1}r_{n}}^{\sum_{n\leq i}r_{n}}(\Sigma^T\Sigma)_{k,k},
\end{equation*}
with the convention $r_0 = 0$. When taking $r_1=d$ we have the following
\begin{equation*}
\gamma_1 = 2|\Sigma^T\Sigma v|^2 - \frac h2\Tr(\Sigma^T\Sigma)
\end{equation*}
where we recognize the last equation of \cite[(1.2)]{LeSaSt20} (up to the change $\Sigma\mapsto\frac12\Sigma$).

In the following we will assume $Q = I_d$ and all the multiplicities of the variables $y_i$ are simple, in other words, $\forall i, r_i = 1$, this way we choose $\Gamma = \diag(\gamma)$ and for $i\in\lb1,d\rb$,
\begin{equation*}
\gamma_i = 2(\sum_{j=1}^d(\Sigma^T\Sigma)_{i,j}v_j)^2 - \frac h2(\Sigma^T\Sigma)_{i,i}.
\end{equation*}

Using the decomposition $x = (x',y)$, this leads to the coefficients $\beta(x',y,v) = -\p_{x'}V(x') - \diag(y)\Sigma^T\Sigma v$ and
\begin{equation*}
\alpha(x',y,v) = \begin{pmatrix}2\Sigma^T\Sigma v\\2\diag(\Sigma^T\Sigma v)\Sigma^T\Sigma v - \frac h2\diag(\Sigma^T\Sigma)\end{pmatrix},
\end{equation*}
and we combine the variables to be in the settings of the previous sections: $x = (x',y)$. We now need to check if these coefficients satisfy the different assumptions we made.

\textbullet\ Assumption \ref{ass.skewadj}: We constructed $\alpha$ and $\beta$ from this starting point, so it is satisfied.

\textbullet\ Assumption \ref{ass.confin} and \ref{ass.morse}: Since $W$ is quadratic, as long as $V$ satisfies both of them, the global potential $\hat V(x',y) = V(x') + W(y)$ does.

\textbullet\ Assumption \ref{ass.G}: There, we need to compute the matrix $G = (G_{i,j})_{1\leq i,j\leq 2d}$,
\begin{equation*}
G_{i,j} = (C_\Sigma h)^{-\frac d2}\int_{\R^d}\alpha_i\alpha_je^{-2|\Sigma v|^2/h}dv.
\end{equation*}
For $i,j\leq d$, $G_{i,j} = 0$ for $i\neq j$ using the change of variable $v\mapsto(\Sigma^T\Sigma)^{-1} v$ and a parity argument. %%
%% Annotation 44, 1/1 on page 23 [ ]
%% Ink: "Rd Σk,iΣn<SEL>,ivkvne−2|v|2/hdv</SEL>  k,n  Z  2 det(Σ−1) X  Σ2  = 4δi,j(CΣh)−d  k,i  k<SEL>e−2|v|2/hdv</SEL>  Rd v2  k  Z  2 det(Σ−1) X  Σ2  = 4δi,j(CΣh)−d  k,i  h 4  Rd e<SEL>−2|v|2/hdv</SEL>  k  "
%% Comment: ""
%%
%⭣ ⭣ ⭣
Then
\begin{equation*}
\begin{aligned}
    G_{i,j} &= 4\delta_{i,j}(C_\Sigma h)^{-\frac d2}\int_{\R^d}(\Sigma^T\Sigma v)_i^2e^{-2|\Sigma v|^2/h}dv\\
    &= 4\delta_{i,j}(C_\Sigma h)^{-\frac d2}\det(\Sigma^{-1})\sum_{k,n}\int_{\R^d}\Sigma_{k,i}\Sigma_{n,i}v_kv_ne^{-2|v|^2/h}dv\\
    &= 4\delta_{i,j}(C_\Sigma h)^{-\frac d2}\det(\Sigma^{-1})\sum_{k}\Sigma_{k,i}^2\int_{\R^d}v_k^2e^{-2|v|^2/h}dv\\
    &= 4\delta_{i,j}(C_\Sigma h)^{-\frac d2}\det(\Sigma^{-1})\sum_{k}\Sigma_{k,i}^2\frac h4\int_{\R^d}e^{-2|v|^2/h}dv\\
    &= \delta_{i,j}h(\Sigma^T\Sigma)_{i,i}
\end{aligned}
\end{equation*}
using %%
%⭡ ⭡ ⭡ END of annotation 44
the expression of $\rho$ (see \eqref{eq:rho}).

If $i\leq d$ and $j>d$ with another parity argument, we again have that $G_{i,j} = 0$. For index greater than $d$, we have from the previous %%
%% Annotation 45, 1/4 on page 24 [ ]
%% Text: "None, annot.type = <Text>"
%% Comment: "please adjust to fit"
%%
%⭣ ⭣ ⭣
computations
\begin{equation*}
\begin{aligned}
    G_{i+d,j+d} &= (C_\Sigma h)^{-\frac d2}\int_{\R^d}(2(\Sigma^T\Sigma v)_i^2 - \frac h2(\Sigma^T\Sigma)_{i,i})(2(\Sigma^T\Sigma v)_j^2 - \frac h2(\Sigma^T\Sigma)_{j,j})e^{-2|\Sigma v|^2/h}dv\\
    &= 4(C_\Sigma h)^{-\frac d2}\int_{\R^d}(\Sigma^T\Sigma v)_i^2(\Sigma^T\Sigma v)_j^2e^{-2|\Sigma v|^2/h}dv\\
    &\phantom{********} - h(C_\Sigma h)^{-\frac d2}(\Sigma^T\Sigma)_{j,j}\int_{\R^d}(\Sigma^T\Sigma v)_i^2e^{-2|\Sigma v|^2/h}dv\\
    &\phantom{********} - h(C_\Sigma h)^{-\frac d2}(\Sigma^T\Sigma)_{i,i}\int_{\R^d}(\Sigma^T\Sigma v)_j^2e^{-2|\Sigma v|^2/h}dv\\
    &\phantom{********} + \frac{h^2}{4}(C_\Sigma h)^{-\frac d2}(\Sigma^T\Sigma)_{j,j}(\Sigma^T\Sigma)_{i,i}\int_{\R^d}e^{-2|\Sigma v|^2/h}dv\\
    &= 4(C_\Sigma h)^{-\frac d2}\int_{\R^d}(\Sigma^T\Sigma v)_i^2(\Sigma^T\Sigma v)_j^2e^{-2|\Sigma v|^2/h}dv - \frac{h^2}{4}(\Sigma^T\Sigma)_{j,j}(\Sigma^T\Sigma)_{i,i}\\
\end{aligned}
\end{equation*}
which, %%
%⭡ ⭡ ⭡ END of annotation 45
with parity arguments leads %%
%% Annotation 46, 2/4 on page 24 [ ]
%% Ink: "4(CΣh)−d  Σk1,iΣk2,iΣk3,jΣk<SEL>4,jvk1vk2vk3vk4e−2|v|2/hdv</SEL>  Rd  k1,k2,k3,k4  X  Σ2  2 det(Σ−1) Z  = 4(CΣh)−d  k,iΣ2  n,jv2  kv2  ne−2|v|2/hd<SEL>v</SEL>  Rd  k̸=n  X  2 det(Σ−1) Z  + 4(CΣh)−d  Σk,iΣn,i<SEL>Σk,jΣn,jv2  kv2  ne−2|v|2/hdv</SEL>  Rd 2  k̸=n  X  2 det(Σ−1) Z  Σ2  + 4(CΣh)−d  k,iΣ2  k,jv<SEL>4  ke−2|v|2/hdv</SEL>  Rd  k"
%% Comment: ""
%%
%% Annotation 47, 3/4 on page 24 [ ]
%% Ink: "None, annot.type = <Ink>"
%% Comment: ""
%%
%% Annotation 48, 4/4 on page 24 [ ]
%% Text: "None, annot.type = <Text>"
%% Comment: "adjust to fit"
%%
%⭣ ⭣ ⭣
to 
\begin{equation*}
\begin{aligned}
    (\ast_1) :&= 4(C_\Sigma h)^{-\frac d2}\det(\Sigma^{-1})\int_{\R^d}(\Sigma^Tv)_i^2(\Sigma^T v)_j^2e^{-2|v|^2/h}dv\\
    &= 4(C_\Sigma h)^{-\frac d2}\det(\Sigma^{-1})\int_{\R^d}\sum_{k_1,k_2,k_3,k_4}\Sigma_{k_1,i}\Sigma_{k_2,i}\Sigma_{k_3,j}\Sigma_{k_4,j}v_{k_1}v_{k_2}v_{k_3}v_{k_4}e^{-2|v|^2/h}dv\\
    &= 4(C_\Sigma h)^{-\frac d2}\det(\Sigma^{-1})\int_{\R^d}\sum_{k\neq n}\Sigma_{k,i}^2\Sigma_{n,j}^2v_k^2v_n^2e^{-2|v|^2/h}dv\\
    &\phantom{********} + 4(C_\Sigma h)^{-\frac d2}\det(\Sigma^{-1})\int_{\R^d}2\sum_{k\neq n}\Sigma_{k,i}\Sigma_{n,i}\Sigma_{k,j}\Sigma_{n,j}v_k^2v_n^2e^{-2|v|^2/h}dv\\
    &\phantom{********} + 4(C_\Sigma h)^{-\frac d2}\det(\Sigma^{-1})\int_{\R^d}\sum_{k}\Sigma_{k,i}^2\Sigma_{k,j}^2v_k^4e^{-2|v|^2/h}dv\\
    &= \frac{h^2}{4}\big(\sum_{k\neq n}\Sigma_{k,i}^2\Sigma_{n,j}^2 + 2\sum_{k\neq n}\Sigma_{k,i}\Sigma_{n,i}\Sigma_{k,j}\Sigma_{n,j} + 3\sum_{k}\Sigma_{k,i}^2\Sigma_{k,j}^2\big)\\
    &= \frac{h^2}{4}\big((\Sigma^T\Sigma)_{i,i}(\Sigma^T\Sigma)_{j,j} + 2(\Sigma^T\Sigma)_{i,j}^2\big).
\end{aligned}
\end{equation*}
Hence %%
%⭡ ⭡ ⭡ END of annotations 46, 47, 48
$G_{i+d,j+d} = \frac{h^2}{2}(\Sigma^T\Sigma)_{i,j}^2$. We can summarize these computations in
\begin{equation*}
G = \frac h2\left(\begin{array}{c|c}
    2\diag(\diag(\Sigma^T\Sigma)) & 0 \\
    \hline
    0 & h(\Sigma^T\Sigma)^{\odot 2}
\end{array}\right),
\end{equation*}
where $\odot$ is the Hadamard product. We then know using Schur's product Theorem \cite[Theorem VII]{Sc11} that $(\Sigma^T\Sigma)^{\odot 2}$ is definite positive since $\Sigma^T\Sigma$ is, which therefore proves Assumption \ref{ass.G} $i)$ and $ii)$ with $g_1(h)\asymp h^2$ and $g_2(h)\asymp h$. This also makes \eqref{eq:gthm2} true. For $iii)$:

let $i,j\in\lb1,d\rb$, with the scaling $v\mapsto\sqrt{h}v$, we observe that
\begin{equation*}
\begin{aligned}
    &\<\alpha_i^2\alpha_j^2\rho,\rho\>_{L^2(\R^d_v)} = O(h^2)\\
    &\<\alpha_{i+d}^2\alpha_j^2\rho,\rho\>_{L^2(\R^d_v)} = O(h^3)\\
    &\<\alpha_{i+d}^2\alpha_{j+d}^2\rho,\rho\>_{L^2(\R^d_v)} = O(h^4)
\end{aligned}
\end{equation*}

\textbullet\ Assumption \ref{ass.hypocoer}: Using that $\alpha$ is a polynomial of order $2$ in $v$, does not depend on the other variables, and recalling that $\beta(x',y,v) = -\p_{x'}V(x') - \diag(y)\Sigma^T\Sigma v - 4\Sigma^T\Sigma v$, $ii)$ of this assumption is easily satisfied.

Indeed, take $q\in h\{J_x\alpha\alpha,J_v\alpha\beta,J_v\alpha\Sigma^T\Sigma v,h\Delta_v\alpha\}$, we have for all $i,j\in\lb1,d\rb$,
\begin{equation*}
\Pi q_iq_j\Pi \lesssim h^2(|\nabla \hat V|^2 + h)\Pi
\end{equation*}
recalling that $\hat V$ denotes the potential in both $x'$ and $y$, using that $\Pi v_i^k\Pi = Ch^{k/2}\Pi$ for some $C>0$. Noticing that we took $g_2(h)\asymp h^2$, we have the result.

Now for $i)$, taking $j\in\lb1,d\rb$, we have $\<\alpha_j\rho,\rho\>_{L^2(\R^d_v)} = 0$ using parity and from the first set of equation of the previous point, we obtain
\begin{equation*}
\<\alpha_{j+d}\rho,\rho\>_{L^2(\R^d_v)} = 2(C_\Sigma h)^{-\frac d2}\int_{\R^d}(\Sigma^T\Sigma v)_j^2e^{-2|\Sigma v|^2/h}dv - \frac{h}{2}(\Sigma^T\Sigma)_{j,j} = 0.
\end{equation*}

\textbullet Assumption \ref{ass.Hormander}: Because $\Sigma^T\Sigma$ is invertible, the family $\p_{v_i}\alpha_j$ for $i,j\in\lb1,d\rb$ will generate $\p_{x'}$. With the same argument, we obtain $\p_y$ from $\p_{v_k}\p_{v_i}\alpha_{j+d}$, with $k\in\lb1,d\rb$.

\textbullet Assumption \ref{ass.accretive}: We first notice that $\alpha$ does not depend on $x$, then $\div_{x}\alpha = 0$. While for $\beta$, $\div_v\beta = -\Tr(\diag(y)\Sigma^T\Sigma)$ hence $|\div_v\beta| \lesssim |y| \leq f^{1/2}\leq f$ near infinity.










\subsection{Example 2 : a Langevin dynamic rescaled in time.}\label{ssec:exampleRescaled}
Another example can be found in \cite{LeLeLaHi24}, the Langevin dynamics is written with the form

\begin{equation}\label{eq:SDEg}
\left\{
\begin{aligned}
&dx_t = g(x_t)v_tdt,\\
&dv_t = -g(x_t)\p_xVdt + h\p_xgdt - g(x_t)v_tdt + \sqrt{2hg(x_t)}dB_t.
\end{aligned}
\right.
\end{equation}
And considering $f(x,v) = \frac{V(x)}{2} + \frac{|v|^2}{4}$, $P_g = -e^{f/h}\circ h\lll^*\circ e^{-f/h}$ has the form $P_g = X_g + g\Delta_{\frac{|v|^2}{4}}$ where
\begin{equation*}
X_g = hg(x)(v\cdot\p_x - \p_xV\cdot\p_v) + h\p_xg\cdot(\frac{v}{2} + h\p_v)
\end{equation*}
satisfies
\begin{equation*}
P_g(e^{-f/h}) = P_g^*(e^{-f/h}) = 0.
\end{equation*}
Let us notice that taking $g\equiv 1$ we recover the usual Fokker-Planck operator.

Looking at the details in Section \ref{sec:Hypocoer}, we see that we mostly manipulate $X$ and rather rarely $N$. So we will consider a generalization of the operator extracted from \eqref{eq:SDEg}, taht is $x,v\in\R^d$, $P = X + gN$, with
\begin{equation*}
\left\{
\begin{aligned}
&X = 2g(x)\Sigma^T\Sigma v\cdot h\p_x - g(x)\p_xV\cdot h\p_v + \frac h2\p_xg\cdot(2\Sigma^T\Sigma v + h\p_v),\\
&N = -h^2\Delta_v + 4|\Sigma^T\Sigma v|^2 - 2h\Tr(\Sigma^T\Sigma).
\end{aligned}
\right.
\end{equation*}
From these expressions, we thus consider
\begin{equation*}
\alpha(x,v) = 2g(x)\Sigma^T\Sigma v\quad \beta(x,v) = - g(x)\p_xV(x) + \frac h2\p_xg(x).
\end{equation*}

Like in \cite{LeLeLaHi24}, we assume $g$ satisfies
\begin{equation}\label{eq:assgExample3}
\exists m,M,\ \forall x\in\R^d,\ 0 < m \leq |\p_xg(x)| \lesssim g(x) \leq M.
\end{equation}
We will see that this condition ensures that all the assumptions for Theorem \ref{thm:1} hold.

Because we have $g$ in factor of $N$, we must be careful to the change it would induce in the proofs of the theorems. First, observe that if $\{g = 0\}$ is of measure $0$, then the kernel of $gN$ and $N$ are the same, and thus $\Pi$ still is the orthogonal projector on the kernel of $gN$. Therefore there is no need to change the definition of $A$, and within the proofs of the lemmas of Section \ref{sec:Hypocoer}, we only use $N$ once, when we bound $NA^*$, which means the other results of that section hold. If $g\in L^\infty$, then $\vvert{gNA^*} \leq \vvert{g}_\infty \vvert{NA^*}$. This lead us to consider both these hypothesis to be true. Now let us look at the other assumptions.

\textbullet\ Assumption \ref{ass.skewadj}: The form of $P$ is made so that this is satisfied.

\textbullet\ Assumption \ref{ass.confin} and \ref{ass.morse}: We just need $V$ to satisfy them.

\textbullet\ Assumption \ref{ass.G}: We observe from the computations of the previous example that $G = hg(x)\diag(\diag(\Sigma^T\Sigma))$. Hence with the assumption that there exists $m,M$ such that for all $x\in\R^d$, we have $0<m\leq g(x)\leq M$ and $\p_ig \lesssim g$ for all $i\in\lb1,d\rb$, then Assumption \ref{ass.G} is satisfied with both $g_1(h)\asymp h$ and $g_2(h)\asymp h$. This also makes \eqref{eq:gthm2} true.

\textbullet\ Assumption \ref{ass.hypocoer}: For $i)$, this is immediate noticing $\alpha$ is odd in $v$. For $ii)$, we need to study each term individually

\quad $\star$ Remark that $hJ_x\alpha\alpha = hg(x)\p_xg|\Sigma^T\Sigma v|^2$, hence we must assume that
\begin{equation*}
|\p_x(g^2)| \lesssim |\p_x V| + 1,
\end{equation*}
which is ensured by \eqref{eq:assgExample3}.

\quad $\star$ We then have $hJ_v\alpha\beta = -2hg(x)^2\Sigma^T\Sigma\p_xV + \frac{h^2}2\Sigma^T\Sigma\p_x(g^2)$. Both term are controlled by $h(|\p_xV| + 1)$ because $g$ is bounded for the first one, and because of the previous point for the other one.

\quad $\star$ Using the two previous points, we also have that $hJ_v\alpha\Sigma^T\Sigma v = 2hg(x)(\Sigma^T\Sigma)^2v$ is well controlled.

\quad $\star$ And lastly $h^2\Delta_v\alpha = 0$.

\textbullet\ Assumption \ref{ass.Hormander}: We have that $\p_{v_i}\alpha_j = 2g(x)(\Sigma^T\Sigma)_{j,i}$. Using that $g$ is bounded from below and $\Sigma^T\Sigma$ is non-degenerate, we can recover all the $\p_{x_i}$.

\textbullet\ Assumption \ref{ass.accretive}: We have that $\div_x\alpha + \div_v\beta = 2\p_xg\cdot\Sigma^T\Sigma$, thus we need to assume $|\p_xg|\lesssim V$ at infinity.










\subsection{Example 3 : adding a magnetic field in the usual KFP operator.}\label{ssec:exampleMagnetic}
The usual KFP equation consist in considering the coefficients
\begin{equation*}
\alpha(x,v) = 2\Sigma^T\Sigma v;\quad \beta(x,v) = - \p_xV(x).
\end{equation*}
In $d = d' = 3$, adding a magnetic field means to add a part of the form $b\wedge v\cdot\p_v$ within the stochastic equation, which means that this term will appear in $\beta$, this leads to
\begin{equation*}
\alpha(x,v) = 2\Sigma^T\Sigma v;\quad \beta(x,v) = - \p_xV(x) + b(x)\wedge\Sigma^T\Sigma v.
\end{equation*}

We now need to check if those coefficients satisfy the assumptions of the previous sections. The part in $2\Sigma^T\Sigma v\cdot h\p_x - \p_xV\cdot h\p_v$ in $P$ is the usual diffusion term of the well-known Fokker-Planck equation, so we will focus on the magnetic term $b(x)\wedge\Sigma^T\Sigma v\cdot \p_v$.

\textbullet\ Assumption \ref{ass.skewadj}: Using that for any $a,u\in\R^3$, $a\wedge u\cdot u = 0$ and since $\alpha$ does not depend on $x$, we have that this assumption is satisfied if and only if $\div_v(b\wedge\Sigma^T\Sigma v) = 0$. If we write $b = (b_i)_{1\leq i\leq3}$ and $M = (M_{i,j})_{i\leq i,j\leq3}$, we observe that
\begin{equation*}
\div_v(b\wedge Mv) = b_1(M_{2,3} - M_{3,2}) + b_2(M_{3,1} - M_{1,3}) + b_3(M_{1,2} - M_{2,1}),
\end{equation*}
hence taking $M = \Sigma^T\Sigma$, which is symmetric, this assumption is satisfied.

We can then write the full operator considered, $P = X + N$, with
\begin{equation*}
\left\{
\begin{aligned}
&X = 2\Sigma^T\Sigma v\cdot h\p_x - \p_xV\cdot h\p_v + b(x)\wedge\Sigma^T\Sigma v\cdot h\p_v,\\
&N = -h^2\Delta_v + 4|\Sigma^T\Sigma v|^2 - 2h\Tr(\Sigma^T\Sigma).
\end{aligned}
\right.
\end{equation*}

\textbullet\ Assumption \ref{ass.confin} and \ref{ass.morse}: We just need to have that $V$ satisfy those.

\textbullet\ Assumption \ref{ass.G}: Since the change in $P$ from the standard KFP operator lies in $\beta$, we have that $G = h\diag(\diag(\Sigma^T\Sigma))$, hence this assumption is satisfied with $g_1(h) = g_2(h) \propto h$.

\textbullet\ Assumption \ref{ass.hypocoer}: For $i)$ this works just like for example \ref{ssec:exampleAdaptive}, and for $ii)$, we only need to check the term $hJ_v\alpha\beta\Pi$, which is $2h\Sigma^T\Sigma\beta\Pi$. The first term appearing is $2h\Sigma^T\Sigma\p_xV\Pi$ which is directly controlled by $g_2(h)|\p_xV|\Pi$. For the other one, we need to work a bit, let $u\in L^2(\R^6)$ and $i\in\lb1,3\rb$%%
%% Annotation 49, 1/1 on page 28 [ ]
%% FreeText: "<SEL>adjust to fit</SEL>  2 Z"
%% Comment: "adjust to fit"
%%
%⭣ ⭣ ⭣
,
\begin{equation*}
\begin{aligned}
    \vvert{h(\Sigma^T\Sigma (b\wedge\Sigma^T\Sigma v))_i\Pi u}^2 &= h^2(C_\Sigma h)^{-\frac d2}\int_{\R^6}u_\rho^2(x)(\Sigma^T\Sigma (b\wedge\Sigma^T\Sigma v))_i^2e^{-2|\Sigma v|^2/h}dvdx\\
    &= h^2(C_\Sigma h)^{-\frac d2}\det(\Sigma^{-1})\\
    &\phantom{********}\times\int_{\R^6}u_\rho^2(x)(\Sigma^T\Sigma (b\wedge\Sigma^T v))_i^2e^{-2|v|^2/h}dvdx\\
    (\Sigma^T\Sigma (b\wedge\Sigma^T v))_i &= \sum_{j=1}^3(\Sigma^T\Sigma)_{i,j}(b\wedge\Sigma^T v)_j.
\end{aligned}
\end{equation*}
For %%
%⭡ ⭡ ⭡ END of annotation 49
$j\in\lb1,3\rb$, we denote $k,\ell\in\lb1,3\rb$ such that $k=j+1[3]$ and $\ell=k+1[3]$, that way we get
\begin{equation*}
\begin{aligned}
    (b\wedge\Sigma^T v)_j &= b_k(\Sigma^T v)_\ell - b_\ell(\Sigma^T v)_k\\
    &= \sum_{n=1}^3(b_k\Sigma_{n,\ell} - b_\ell\Sigma_{n,k})v_n.
\end{aligned}
\end{equation*}
Using a parity argument to get rid of terms of the form $v_nv_m$ where $n\neq m$, we %%
%% Annotation 50, 1/1 on page 29 [ ]
%% FreeText: "Z  2  <SEL>adju st to  fit</SEL>  3 X"
%% Comment: "adjust to fit"
%%
%⭣ ⭣ ⭣
obtain
\begin{equation*}
\begin{aligned}
    \vvert{h(\Sigma^T\Sigma b\wedge\Sigma^T\Sigma v)_i\Pi u}^2 &=h^2(C_\Sigma h)^{-\frac d2}\det(\Sigma^{-1})\sum_{n=1}^3\Bigg(\int_{\R^3}v_n^2e^{-2|v|^2/h}dv\\
    &\phantom{****}\times\int_{\R^3}u_\rho^2(x)\Big(\sum_{j=1}^3(\Sigma^T\Sigma)_{i,j}(b_k(x)\Sigma_{n,\ell} - b_\ell(x)\Sigma_{n,k})\Big)^2dx\Bigg)\\
    &= \frac{h^3}{4}\int_{\R^3}u_\rho^2(x)\sum_{n=1}^3\Big(\sum_{j=1}^3(\Sigma^T\Sigma)_{i,j}(b_k(x)\Sigma_{n,\ell} - b_\ell(x)\Sigma_{n,k})\Big)^2dx\\
    &= \frac{h^3}{4}\int_{\R^3}u_\rho^2(x)\sum_{n=1}^3(\Sigma^T\Sigma (b(x)\wedge\Sigma^T e_n))_i^2dx\\
\end{aligned}
\end{equation*}
denoting %%
%⭡ ⭡ ⭡ END of annotation 50
$e_n$ the $n$-th element of the canonical basis of $\R^3$. Assuming we have $|b(x)| \lesssim |\p_xV(x)| + 1$ on all $\R^3$, we will then have the assumption verified. Notice that constant or bounded magnetic fields satisfy this bound.

\textbullet\ Assumption \ref{ass.Hormander}: Using that $\p_{v_j}\alpha_i\p_{x_i} = 2(\Sigma^T\Sigma)_{i,j}\p_{x_i}$ this assumption is directly satisfied.

\textbullet\ Assumption \ref{ass.accretive}: Since we have no divergence, this assumption is trivially verified.













































\section{Geometric construction of the quasimodes}\label{sec:geoloc}

We now want to have a better view on the small eigenvalues of $P$. For this purpose, we are going to build sharp quasimodes, and so we are following the steps of \cite[section 3\&4]{BoLePMi22}. Their theorem does not apply here because $P$ does not satisfy either (1.9), (Harmo), (Hypo) or (Morse) from their paper for general $\alpha,\beta$ and $V$ we consider. Therefore, the work here is to find a way to obtain a similar result without these assumptions.

Due to several inconvenience, we will consider $V$ Morse from this section onward. However, the other assumptions from \cite{BoLePMi22} still remain false in the general case. To have an overview of what behavior can arise when $V$ is not a Morse function, we refer to \cite{De25_1} where we study the Witten Laplacian $\Delta_V$ with degenerate potentials a similar way to the one we present here.

As in \cite{BoLePMi22}, given $\s\in\vvv^{(1)}$ we look for an approximate solution to the equation $Pu = 0$ in a neighborhood $\www = \www_x\times\www_v$ of $\s$.

We look for an approximate solution of $P u = 0$ of the form
\begin{equation*}
u=\chi e^{-(f-f(\m))/h},
\end{equation*}
and we set
\begin{equation*}
\chi(x,v)=\int_0^{\ell(x,v,h)}\zeta(s/\tau)e^{-\frac{s^2}{2h}}ds
\end{equation*}
where the function $\ell\in\ccc^\infty(\www)$ has a formal classical expansion $\ell\sim\sum_{j\geq0}h^j\ell_j$. Here, $\zeta$ denotes a fixed smooth even function equal to $1$ on $[-1,1]$ and supported in $[-2,2]$, and $\tau > 0$ is a small parameter which will be fixed later.

The object of this section is to construct the function $\ell$. Here we adapt the construction made in \cite[section 3]{BoLePMi22}.

Since $P(e^{-f/h}) = 0$ we have that $P(\chi e^{-f/h}) = [P,\chi](e^{-f/h})$, but
\begin{equation*}
\begin{aligned}~
    [P,\chi] &= \alpha\cdot h\p_x\chi + \beta\cdot h\p_v\chi - h^2[\Delta_v,\chi]\\
    &= \alpha\cdot h\p_x\chi + \beta\cdot h\p_v\chi - h^2(\Delta_v \chi + 2\p_v \chi\cdot\p_v),\\
    \p_v \chi &= \p_v\ell\zeta(\ell/\tau)e^{-\frac{\ell^{2}}{2h}},\\
    \p_x \chi &= \p_x\ell\zeta(\ell/\tau)e^{-\frac{\ell^{2}}{2h}},\\
    \Delta_v \chi &= \Big(\Delta_v\ell\zeta(\ell/\tau) + \frac1\tau|\p_v\ell|^2\zeta'(\ell/\tau)-\zeta(\ell/\tau)|\p_v\ell|^2\frac{\ell}{h}\Big)e^{-\frac{\ell^{2}}{2h}}.
\end{aligned}
\end{equation*}
Therefore, if we set that $\ell_0(\s) = 0$, there exists $r$ smooth such that $r\equiv0$ near $\s$, $r$ and its derivatives are locally uniformly bounded with respect to $h$ and
\begin{equation}\label{eq:wplusr}
P(\chi e^{-f/h}) = h(w+r)e^{-(f+\frac{\ell^{2}}{2})/h},
\end{equation}
where
\begin{equation*}
w = \alpha\cdot \p_x\ell + \beta\cdot \p_v\ell + 4\Sigma^T\Sigma v\cdot\p_v\ell + \ell|\p_v\ell|^2 - h\Delta_v\ell.
\end{equation*}

If $\alpha$ and $\beta$ admits a similar classical expansion $\ell$'s : $\alpha\sim\sum_{j\geq0}h^j\alpha^j$, $\beta\sim\sum_{j\geq0}h^j\beta^j$, we can see that $w$ admits a formal classical expansion $w \sim \sum_{j\geq0} h^jw_j$. Solving formally $w=0$ and identifying the powers of $h$ leads to a system of equations
\begin{equation}\label{eq:eik}\tag{eik}
\alpha^0\cdot \p_x\ell_0 + \beta^0\cdot \p_v\ell_0 + 4\Sigma^T\Sigma v\cdot\p_v\ell_0 + \ell_0|\p_v\ell_0|^2 = 0
\end{equation}
and for $j\geq1$,
\begin{equation}\label{eq:transp}\tag{T$_j$}
\big(\alpha^0\cdot \p_x + (\beta^0+ 4\Sigma^T\Sigma v + 2\ell_0\p_v\ell_0)\cdot\p_v + |\p_v\ell_0|^2\big) \ell_j + R_j = 0
\end{equation}
where $R_j$ is a smooth polynomial of the $\p^\gamma\ell_k$ for $|\gamma|\leq2$ and $k<j$. As in \cite{BoLePMi22} and by analogy with the WKB method, we call eikonal equation the first one and transport equations the next ones.
\begin{proposition}\label{prop:ell}There exists $\ell\in\ccc^\infty(\www)$ satisfying the following
    \begin{itemize}
        \item[$i)$] The eikonal and transport equations are solved up to any order leading to
        \begin{equation}\label{eq:Pchiell}
            P(\chi e^{-f/h}) = hO(X^\infty + h^\infty)e^{-(f+\frac{\ell^{2}}{2})/h}
        \end{equation}
        denoting $X=(x,v)$ and we have $|\nabla\ell_{0}(\s)|^2\neq0$.
        \item[$ii)$] Moreover, either $|\p_v\ell_{0}(\s)|^2\neq0$ or there exists a non-zero, $h$-independent, positive semidefinite matrix $A$, such that $|\p_v\ell_{0}|^2 = 4\<A\nabla f,\nabla f\>(1+O(X-\s))$, and $\<A\nabla f,\nabla f\>$ is not identically $0$ near $\s$.
        \item[$iii)$] Furthermore, $\ell$ elliptizes $f$ around $\s$ in the following sense
        \begin{equation}\label{eq:dethess}
        \det\Hess_\s\big(f + \frac{\ell_0^{2}}{2}\big) = -\det\Hess_\s f.
        \end{equation}
        \item[$iv)$] Finally,
        \begin{equation*}
        \forall X\in\www\setminus\{\s\},\ \ \ X-\s\in\eta(\s)^\bot = 0 \Rightarrow f(X) > f(\s),
        \end{equation*}
        where $\eta(\s) = \nabla\ell_{0}(\s)$.
    \end{itemize}
\end{proposition}

\begin{remark}\label{rem:signell}
    Note that $-\ell$ solves the equations the same way $\ell$ does. For now, the sign does not matter, but we will fix it in the next section when properly constructing the quasimodes on a global setting.
\end{remark}

In the following, we will not try to solve \eqref{eq:eik} and \eqref{eq:transp} in their utmost generality due to many different difficulties. We will rather show what behavior can arise when treating with some of the degeneracy one can encounter. We manage to prove Proposition \ref{prop:ell} in all three situations we describe thereafter. Up to translations, we can also assume without loss of generality that $\s = 0$ and $V(\s) = 0$.










\subsection{Situation 1 : partially degenerate $\alpha$ and $\beta$.}\label{sub:ex1geo}

In this first example, we consider a case generalizing a bit the Section \ref{ssec:exampleAdaptive}. We assume the differential $d_{x,v}\alpha^0$ is non-zero (or equivalently from the next equation $d_{x,v}\beta^0\neq0$). We use the relation \eqref{eq:assskew}, which gives when identifying the coefficients of $h^0$:
\begin{equation*}
\alpha^0\cdot\p_xV + 2\beta^0\cdot\Sigma^T\Sigma v = 0.
\end{equation*}
Therefore, assuming $\alpha^0$ has a non-zero linear part, we can write $\alpha^0 = Mv + o(v)$ for some $M\in\Mr_{d',d}(\R)$. Hence we have $\beta^0 = M_\beta x + o(x)$, with $M_\beta\in\Mr_{d,d'}(\R)$, and at its principal order, the previous equation becomes
\begin{equation*}
(M^TH + 2\Sigma^T\Sigma M_\beta)x\cdot v = 0
\end{equation*}
denoting $H$ the Hessian of $V$ at $\s$. This gives the relation $M_\beta = -\frac12(\Sigma^T\Sigma )^{-1}M^TH$. We now denote $X = (x,v)$ and consider $\ell_0 = \xi\cdot X + O(X^2)$ for some $\xi=(\xi_x,\xi_v)\in\R^{d+d'}$. Therefore, at its principal order, \eqref{eq:eik} becomes
\begin{equation}\label{eq:eik1}
(\Lambda\xi + |\xi_v|^2\xi)\cdot X = 0,
\end{equation}
where $\Lambda = \begin{pmatrix}
    0&-\frac12HM(\Sigma^T\Sigma)^{-1}\\
    M^T&4\Sigma^T\Sigma
\end{pmatrix}$.

We therefore have that \eqref{eq:eik1} is satisfied, for a non-zero $\xi$, if and only if $\xi$ is an eigenvector of $\Lambda$ associated with the eigenvalue $-|\xi_v|^2$. Let $\mu\geq0$,
\begin{equation}\label{eq:Lambdamu}
\Lambda\xi = -\mu\xi \iff \left\{\begin{array}{l}
    \frac12HM(\Sigma^T\Sigma)^{-1}\xi_v = \mu\xi_x,\\
    M^T\xi_x+ 4\Sigma^T\Sigma\xi_v = -\mu\xi_v.
\end{array}\right.
\end{equation}

Provided that this system has a solution with $\mu\neq0$, we can solve the eikonal equation up to the first order. Note that $\mu\neq0$ implies $\xi_v\neq0$ because we want $\xi\neq0$. We now consider the following assumption
\begin{equation}\label{ass.simple}\tag{Simple}
    \exists \mu>0,\ \ -\mu\in\sigma(\Lambda)\subset\{-\mu\}\sqcup\{\Re z\geq0\}\text{ and }-\mu\text{ is simple.}
\end{equation}

We observe that this case is not always contained in \cite{BoLePMi22}. If $M^T$ has a non-trivial kernel, then the Kalman-type criterion \cite[Remark 2.5]{BoLePMi22} is not verified, which means that their assumption (Harmo) cannot be satisfied. To see this, let $\Sigma^T\Sigma = I_{d'}$ and $\eta\in\Ker(M^T)$ then with their notations $A^0 = \begin{pmatrix}0&0\\0&I_{d'}\end{pmatrix}$ and $B^T = \begin{pmatrix}0&M_\beta^T\\M^T&0\end{pmatrix}$, hence
\begin{equation*}
(\eta,0)\in\Ker(A^0)\cap\Ker(B^T)\subset\bigcap_{n=0}^{d+d'-1}\Ker(A^0(B^T)^n).
\end{equation*}

Note that as soon as $d>d'$, $M^T$ is not injective, this gives a setting where we have some information that is not present in \cite{BoLePMi22}.

Notice that when $\Sigma = \Id$, we have from \eqref{eq:Lambdamu}
\begin{equation*}
\Lambda\xi = -\mu\xi \iff \left\{\begin{array}{l}
    \xi_v = \frac{-1}{4+\mu}M^T\xi_x,\\
    HMM^T\xi_x = -2\mu(4+\mu)\xi_x.
\end{array}\right.
\end{equation*}
Hence \eqref{ass.simple} is equivalent to the same assumption with $\Lambda$ replaced by $HMM^T$ which may be easier to verify in practice. For example with Section \ref{ssec:exampleAdaptive}, we observe that in this case
\begin{equation*}
M = \begin{pmatrix}2\Sigma^T\Sigma&0\\0&0\end{pmatrix},
\end{equation*}
hence taking $\Sigma^T\Sigma = I_{d'}$ leads to
\begin{equation*}
HMM^T = \begin{pmatrix}4\Hess_\s V&0\\0&0\end{pmatrix},
\end{equation*}
which satisfies \eqref{ass.simple} by definition of $\s$ being a saddle point of $V$.

We now look for $\ell_0$ that admits a decomposition $\ell_0\sim\sum_{j\geq0}\ell_{0,j}$ with $\ell_{0,j}\in\ppp^j_{hom}(X)$, where $\ppp^j_{hom}(X)$ denotes the set of homogeneous polynomials in the $X$ variables of degree $j$. Upon assuming $\alpha^0$ and $\beta^0$ admit similar decompositions, so does $w_0$. Moreover we have, recalling $\nu_1=2$, for all $j\geq1$
\begin{equation*}
w_{0,j} = ((\Lambda^T + 2A\Pi_\xi)X\cdot\nabla + \mu)\ell_{0,j} + R_{0,j}
\end{equation*}
where $A$ is the projector on the velocity: $A:(x,v)\in\R^{d+d'}\mapsto v\in\R^{d'}$, $\Pi_\xi$ the one on $\xi$, from now on we denote $\mu = |\xi_v|^2 = |\p_v\ell_{0,1}|^2$, and $R_{0,j}$ is a smooth polynomial of the $\p^\gamma\ell_{0,k}$ for $|\gamma|\leq1$ and $k<j$.

Under Assumption \eqref{ass.simple}, we have that $\lll_0 = (\Lambda^T + 2A\Pi_\xi)X\cdot\nabla + \mu$ is an automorphism of $\ppp^j_{hom}(X)$ for all $j$. In order to prove this, denoting $\Upsilon = \Lambda^T + 2A\Pi_\xi$, it is sufficient to prove that $\sigma(\Upsilon)\subset\{\Re z\geq0\}$ thanks to \cite[Lemma A.1]{BoLePMi22} (this lemma gives the result for $\{\Re z>0\}$ but we can easily extend it to $\{\Re z\geq0\}$).

In a basis of $\C^{d+d'}$ adapted to $\xi$ in which $\Lambda$ is upper triangular, we have that only the first entry of its diagonal has negative real part, being $-\mu$. Moreover in that same basis, $2\Pi_\xi A$ has zeros outside its first row, and the first element of that row is $2\mu$. Hence all the eigenvalues of $\Upsilon^T$, and thus of $\Upsilon$, have non-negative real part.

This way we solved \eqref{eq:eik} at infinite order: there exists $\ell_0\sim\sum_{j\geq0}\ell_{0,j}$ such that
\begin{equation}\label{eq:resoleikinf}
\alpha^0\cdot \p_x\ell_0 + \beta^0\cdot \p_v\ell_0 + 4\Sigma^T\Sigma v\cdot\p_v\ell_0 + \ell_0|\p_v\ell_0|^2 = O(X^\infty).
\end{equation}

In order to give a proper definition of $\ell_0\sim\sum_{j\geq0}\ell_{0,j}$, we can use a Borel procedure, which makes the sum converge in $\ccc^\infty(\www)$ and such that $\ell_0$ still solves \eqref{eq:resoleikinf}.

Let us now solve the transport equations \eqref{eq:transp}. We observe that these equations are of the form
\begin{equation*}
\lll\ell_j + R_j = 0
\end{equation*}
with $\lll = \lll_0 + \lll_>$ where $\lll_>(p) = O(X^{j+1})$ for $p\in\ppp^j_{hom}(X)$. Using the same argument as for \eqref{eq:eik}, we can thus solve all the transport equations. And with another Borel procedure in the $h$ variable, we can make sense of $\ell\sim\sum_{j\geq0}h^j\ell_j$ in $\ccc^\infty(\www)$. This leads to
\begin{equation*}
P(\chi e^{-f/h}) = hO(X^\infty + h^\infty)e^{-(f+\frac{\ell^2}{2})/h}.
\end{equation*}

The proof of Proposition \ref{prop:ell} $iii)$, works just like for \cite[Lemma 3.3]{BoLePMi22}. Denoting $\hat H = \begin{pmatrix}H&0\\0&2\Sigma^T\Sigma\end{pmatrix}$, and $E = 1 + \hat H^{-1}\Pi_\xi$, we have that $\Hess_\s(f+\frac{\ell_0^2}{2}) = \hat HE$. Moreover $E \equiv1$ on the hyperplane $\xi^\bot$ and $\<E\xi,\xi\> = -\vvert{\xi}^2$, hence $\det(\hat HE) = -\det \hat H >0$. This also proves $iv)$.










\subsection{Situation 2 : $\alpha$ and $\beta$ with no linear part.}\label{ssec:nolin}

Let us now consider a case where $\alpha$ and $\beta$ are degenerate in all directions. For this purpose we study the %%
%% Annotation 51, 1/1 on page 34 [ ]
%% Ink: "None, annot.type = <Ink>"
%% Comment: ""
%%
%⭣ ⭣ ⭣
operator
\begin{equation*}
P =  (v^2-h)\cdot h\p_x - 2v\p_xV\cdot h\p_v - h\p_xV + \Delta_{\frac{|v|^2}{4}}
\end{equation*}
acting %%
%⭡ ⭡ ⭡ END of annotation 51
on $(x,v)\in\R^{1+1}$, and with $\Sigma = \frac12$. Although this model does not quite fit in any example of hypocoercive operators given in the previous section, one can easily check that this operator is hypocoercive, with $g_1(h) \asymp g_2(h) \asymp h^2$ (where $g_1$ and $g_2$ are defined in Assumption \ref{ass.G}).

Since we are one-dimensional, we can write $\p_xV = V'$, we will use both notations but mostly the first one, to emphasize the dependence in $x$ and to see how it can be generalized to a multi dimensional case.

This operator obviously does not fall in the framework of \cite{BoLePMi22}, in particular the Kalman-type condition \cite[Remark 2.5]{BoLePMi22} is not satisfied since in their notations, $A^0 = \begin{pmatrix}0&0\\0&1\end{pmatrix}$ and $B = 0$. Recall the equation $w=0$, being
\begin{equation}\label{eq:w=0situ2}
(v^2-h)\cdot \p_x\ell - 2v\p_xV\cdot \p_v\ell + v\cdot\p_v\ell + \ell|\p_v\ell|^2 - h\Delta_v\ell = 0.
\end{equation}
And the eikonal equation is
\begin{equation*}
v^2\cdot \p_x\ell_0 - 2v\p_xV\cdot \p_v\ell_0 + v\cdot\p_v\ell_0 + \ell_0|\p_v\ell_0|^2 = 0.
\end{equation*}

Working with homogeneous polynomials just like in the previous example we have the principal order of the eikonal equation that is
\begin{equation*}
v\xi_v + \xi\cdot X|\xi_v|^2 = 0
\end{equation*}
recalling $\ell_{0,1}(X) = \xi\cdot X$, $\xi\in\R^2$. We notice that if $\xi_v\neq0$, then we must have $\xi_x=0$. But another purpose of these constructions is to change the nature of the saddle point $\s$ of $V$ into a minimum of $V + \frac{\ell^2}{2}$, which cannot be achieved if $\xi_x=0$. Therefore, we must set $\xi_v=0$, and thus $\xi_x$ remains free.

Before going further in the resolution of the equations, we recall we have $\ell\sim\sum_{j\geq0}h^j\ell_j$, $\ell_j$ $h$-independent, $\ell_j\sim\sum_{k\geq0}\ell_{j,k}$, $\ell_{j,k}\in\ppp^k_{hom}(X)$. We now also consider $\ell_{j,k} = \sum_{a+b=k}\ell_{j,a,b}$ with $\ell_{j,a,b}\in\R x^av^b$. Note that this decomposition is unique. This leads to the same decomposition for $w$ and identifying the monomials that leads to $w_{j,a,b}=0$ %%
%% Annotation 52, 1/1 on page 35 [ ]
%% Circle: "None, annot.type = <Circle>"
%% Comment: ","
%%
%⭣ ⭣ ⭣
gives
\begin{equation}\label{eq:decompjab}
\begin{aligned}
    &v^2\p_x\ell_{j,a+1,b-2} - \p_x\ell_{j-1,a+1,b} - 2v\sum_{k=1}^a\p_x^{k+1}V(0)\frac{x^k}{k!}\p_v\ell_{j,a-k,b}\\
    &\phantom{********}+ v\p_v\ell_{j,a,b} - \p_v^2\ell_{j-1,a,b+2} + p_{j,a,b} = 0
\end{aligned}
\end{equation}
where %%
%⭡ ⭡ ⭡ END of annotation 52
$p_{j,a,b}$ contains the terms coming from $|\p_v\ell|^2\ell$, that is
\begin{equation*}
p_{j,a,b} = \sum\ell_{j_1,a_1,b_1}\p_v\ell_{j_2,a_2,b_2}\p_v\ell_{j_3,a_3,b_3}
\end{equation*}
with the sum being on the indexes $j_i,a_i,b_i$, $i\in\lb1,3\rb$ such that
\begin{equation*}
j_1 + j_2 + j_3 = j;\ a_1 + a_2 + a_3 = a;\ b_1 + b_2 + b_3 = b.
\end{equation*}
All this with the convention that $\ell_{j',a',b'} = 0$ if $(j',a',b')\notin\N^3$. The goal now is to find monomials $\ell_{j,a,b}$ that solves \eqref{eq:decompjab} for all $j,a,b\in\N$.

From the structure of the eikonal equation, we have the following lemma.
\begin{lemma}\label{lem:situ2}
    Any function $\ell$ solution of \eqref{eq:w=0situ2} is odd with respect to $v$.
\end{lemma}
We are the one building an $\ell$ solution so we could just look for one that is odd with respect to $v$, but we can actually prove it must be so. Therefore it is not needed to prove this result to continue, this is why we postpone the proof in the appendix \ref{ssec:lemma}.

This property in particular implies that $\p_v\ell = O(v)$ hence plugging $v=0$ in the definition of $w$ we get $\p_x\ell_{|v=0} + \p_v^2\ell_{|v=0} = 0$, which gives for all $j,a\in\N$
\begin{equation*}
\p_x\ell_{j,a+1,0} + \p_v^2\ell_{j,a,2} = 0.
\end{equation*}
We observe that integrating with respect to $v$ this equation gives
\begin{equation}\label{eq:linkell2}
v\p_x\ell_{j,a+1,0} + \p_v\ell_{j,a,2} = 0,
\end{equation}
which will be useful later.

Note also that thanks to the structure of $w$, having that $\ell$ is even in $v$ directly implies that the odd terms in $w$ are $0$ without further assumptions. Moreover, using that $\p_v\ell = O(v)$, we therefore have that for $b=0$ \eqref{eq:decompjab} is automatically satisfied.

In order to determine the rest of the terms in the expansion of $\ell$, the aim is to solve \eqref{eq:decompjab} in increasing order of $j$, namely solve the eikonal equation first and then the series of the transport ones.

For the eikonal equation, which correspond to $j=0$, we have
\begin{equation}\label{eq:decomp0ab}
v^2\p_x\ell_{0,a+1,b-2} - v\sum_{k=1}^a\theta_kx^k\p_v\ell_{0,a-k,b} + v\p_v\ell_{0,a,b} + p_{0,a,b}(x,v) = 0,
\end{equation}
where we denote $\theta_k = \frac{2}{k!}\p_x^{k+1}V(0)$. Let us now solve \eqref{eq:decomp0ab} by induction on $b$. We already know that $b=0$ gives a trivial equation and so does $b\in2\N+1$, hence there remains to consider $b\in2\N^*$. For $b=2$ we have
\begin{equation}\label{eq:decomp0a2-1}
v(v\p_x\ell_{0,a+1,0} + \p_v\ell_{0,a,2}) -v\sum_{k=1}^a\theta_kx^k\p_v\ell_{0,a-k,2} + p_{0,a,2}(x,v) = 0.
\end{equation}
Observe that the first term is zero thanks to \eqref{eq:linkell2}. Note also that because $\ell_{0,a,b} = 0$ for all odd $b$ thanks to Lemma \ref{lem:situ2}, and $\ell_{0,0,0} = 0$, we have $a=0\Rightarrow p_{0,a,2}(x,v)=0$, so for $a=0$, \eqref{eq:decomp0a2-1} is automatically satisfied. And for $a=1$ we get
\begin{equation*}
- v\theta_{1}x\p_v\ell_{0,0,2} + \ell_{0,1,0}|\p_v\ell_{0,0,2}|^2 = 0,
\end{equation*}
using again \eqref{eq:linkell2} which gives $\ell_{0,1,0} = \xi_x x$ and $\p_v\ell_{0,0,2} = -\xi_xv$, we obtain
\begin{equation*}
v^2\theta_{1}x\xi_x + \xi_xx|\xi_xv|^2 = 0
\end{equation*}
leading to $|\xi_x|^2 = -\theta_{1}$ (note that because $\s$ is a saddle point of $V$ according to our definition, $\theta_1 = 2\p_x^{2}V(\s)<0$, hence the negative sign makes sense). Noticing that $\ell$ solves $w=0$ if and only if $-\ell$ do so, we can thus choose $\xi_x =\sqrt{-\theta_1}$ for now.

For $a\geq2$, writing
\begin{equation*}
p_{0,a,2}(x,v) = \sum_{c + d + e = a}\ell_{0,c,0}\p_v\ell_{0,d,2}\p_v\ell_{0,e,2},
\end{equation*}
we observe that (since $\ell_{0,0,0} = 0$)
\begin{equation*}
p_{0,a,2}(x,v) = \ell_{0,a,0}|\p_v\ell_{0,0,2}|^2 + 2\ell_{0,1,0}\p_v\ell_{0,0,2}\p_v\ell_{0,a-1,2} + q_{0,a,2}(x,v)
\end{equation*}
where $q_{0,a,2}$ only contains terms of the form $\p_v\ell_{0,a',2}$ with $a'<a-1$ and $\ell_{0,a'',0}$ with $a''<a$. Therefore we can write \eqref{eq:decomp0a2-1} as
\begin{equation*}
-v\theta_1x\p_v\ell_{0,a-1,2} + \ell_{0,a,0}|\p_v\ell_{0,0,2}|^2 + 2\ell_{0,1,0}\p_v\ell_{0,0,2}\p_v\ell_{0,a-1,2} = R_{0,a,2},
\end{equation*}
with $R_{0,a,2}$ a smooth polynomial of $\p_v\ell_{0,a',2}$ with $a'<a-1$ and $\ell_{0,a'',0}$ with $a''<a$. Using another time \eqref{eq:linkell2}, we can say that $R_{0,a,2}$ is a smooth polynomial of $\ell_{0,a'',0}$ with $a''<a$ exclusively, and that the previous equation can be rewritten
\begin{equation}\label{eq:decomp0a2}
\lll \ell_{0,a,0} = R_{0,a,2}.
\end{equation}
where $\lll = -v^2\theta_1(x\p_x+1)$ which is invertible over $\ppp^{a}_{hom}(X)$ for all $a\geq2$, then we can solve \eqref{eq:decomp0a2} for all $a\geq2$ by a direct induction. It determines all the $\ell_{0,a,0}$ as well as the $\ell_{0,a,2}$ thanks to \eqref{eq:linkell2}.

In conclusion, solving \eqref{eq:decomp0ab} for all $a\geq0$ and $b\leq2$ determines $\ell_{0,a,b}$ for all $a\geq0$, $b\leq2$.

Now, let $b\geq4$ and assume we have constructed the $\ell_{0,a,b'}$ for all $a\in\N$ and $b'<b$, then \eqref{eq:decomp0ab} can be rewritten
\begin{equation*}
v\p_v\ell_{0,a,b} = R_{0,a,b}
\end{equation*}
with $R_{0,a,b}$ a smooth polynomial of the $\p^\delta\ell_{0,a',b'}$ with $|\delta|\leq1$ and $b'\leq b$ or $b'=b$ but $a'<a$. Because $b\neq0$, $v\p_v$ is invertible over $\R x^av^b$ and we can construct by a direct induction on $a\in\N$ all the $\ell_{0,a,b}$.

For the transport equations, we can observe they have a structure very similar to the eikonal one. Let us solve them by induction on $j$. Let $j\geq1$ and assume we have constructed $\ell_{k,a,b}$ for all $k<j$, $a,b\in\N$. Recall that choosing $b=0$ in \eqref{eq:decompjab} gives $0=0$. For $b=2$ we have
\begin{equation*}
v(v\p_x\ell_{j,a+1,0} + \p_v\ell_{j,a,2}) - v\sum_{k=1}^a\theta_kx^k\p_v\ell_{j,a-k,2} + p_{j,a,2}(x,v) = R_{j,a,2}
\end{equation*}
with $R_{j,a,2} = \p_x\ell_{j-1,a+1,2} + \p_v^2\ell_{j-1,a,4}$. Here again, the first term is zero thanks to \eqref{eq:linkell2}. Consider $a=0$, we obtain
\begin{equation*}
\ell_{j,0,0}|\p_v\ell_{0,0,2}|^2 = \tilde R_{j,0,2}
\end{equation*}
with $\tilde R_{j,0,2}$ a smooth polynomial of the $\p^\delta\ell_{k,a',b'}$ for $|\delta|\leq1$ and $k<j$. Because $|\p_v\ell_{0,0,2}|^2 \neq 0$ this determines $\ell_{j,0,0}$. Let now $a\in\N^*$, assume that the $\ell_{j,a',0}$ are constructed for all $a'<a$. Then by \eqref{eq:linkell2}, the $\ell_{j,a',2}$ are also determined for $a'<a-1$, and we have
\begin{equation}\label{eq:decompja2}
-v\theta_1x\p_v\ell_{j,a-1,2} + \ell_{j,a,0}|\p_v\ell_{0,0,2}|^2 + 2\ell_{0,1,0}\p_v\ell_{0,0,2}\p_v\ell_{j,a-1,2} = \tilde R_{j,a,2}
\end{equation}
where $\tilde R_{j,a,2}$ is a smooth polynomial of the $\p^\delta\ell_{k,a',b'}$ for $|\delta|\leq1$ and $k<j$ or $k=j$ and either $a'<a$ with $b=0$, or $a'<a-1$. Just like for the eikonal equation, using \eqref{eq:linkell2}, we obtain
\begin{equation*}
\lll\ell_{j,a,0} = \tilde R_{j,a,2}
\end{equation*}
recalling $\lll = -v^2\theta_1(x\p_x+1)$ which is invertible over $\ppp^a_{hom}(X)$. By induction we solve \eqref{eq:decompja2} for all $a\in\N$.

Finally for $b\geq4$, assume we have constructed the $\ell_{j,a,b'}$ for all $a\in\N$ and $b'<b$, then \eqref{eq:decompjab} can be rewritten
\begin{equation*}
v\p_v\ell_{j,a,b} = R_{j,a,b}
\end{equation*}
with $R_{j,a,b}$ a smooth polynomial of the $\p^\delta\ell_{k,a',b'}$ with $|\delta|\leq1$, $k<j$ and either $b'\leq b$ or $b'=b$ but $a'<a$. Because $b\neq0$, $v\p_v$ is invertible over $\R x^av^b$ and we can construct by a direct induction on $a\in\N$ all the $\ell_{j,a,b}$.

Using Borel procedures to give sense to $\ell_j\sim\sum_{k\geq0}\ell_{j,k}$ and $\ell\sim\sum_{j\geq0}h^j\ell_j$ in $\ccc^\infty(\www)$, we have solved $w = O(X^\infty + h^\infty)$ and thus we obtain
\begin{equation*}
P(\chi e^{-f/h}) = hO(X^\infty + h^\infty)e^{-(f+\frac{\ell^2}{2})/h}.
\end{equation*}

Note that for this example, $|\p_v\ell_0(\s)|^2 = 0$, but Proposition \ref{prop:ell} $ii)$ is satisfied because $|\p_v\ell_0|^2 = - 2v^2\p_x^2V(\s)(1+O(X))$. For the last two items, we have
\begin{equation*}
(f+\frac{\ell_0^2}{2})(\s) = f(\s) + \frac12\p_x^2V(\s)x^2 + \frac{v^2}{4} + \frac{\ell_{0,1,0}^2}{2} + \ell_{0,1,0}\ell_{0,0,2} + O(v^4 + x^4).
\end{equation*}
Note that $\ell_{0,1,0}\ell_{0,0,2} = O(|x|v^2) = O(|x|^3 + |v|^3)$, and recall that $\ell_{0,1,0} = \xi_xx$, $|\xi_x|^2 = -\theta_1$ and $\theta_1 = -2\p_x^2V(\s)$. Combined together we have
\begin{equation*}
(f+\frac{\ell_0^2}{2})(\s) = f(\s) + \frac12|\p_x^2V(\s)|x^2 + \frac{v^2}{4} + O(|(x,v)|^3),
\end{equation*}
which proves $iii)$ and $iv)$.










\subsection{Situation 3 : adding a magnetic source.}

Here we consider the operator
\begin{equation*}
P = 4\Sigma^T\Sigma v\cdot h\p_x - 2\p_xV\cdot h\p_v + b(x)\wedge \Sigma^T\Sigma v\cdot h\p_v + \Delta_{|\Sigma v|^2}
\end{equation*}
acting on $(x,v)\in\R^{3+3}$. We will show that this operator satisfies all the hypotheses made in \cite{BoLePMi22} except possibly for $(1.9)$ when $b$ is not bounded. First observe that (Confin), (Gibbs) and (Morse) are already implied by our assumptions, hence there only remains to look at (Harmo) and (Hypo).

For (Harmo), we use \cite[Corollary 2.4 and Remark 2.5]{BoLePMi22}: we have to check %%
%% Annotation 54, 2/2 on page 39 [ ]
%% Polygon: "None, annot.type = <Polygon>"
%% Comment: ","
%%
%⭣ ⭣ ⭣
that
\begin{equation*}
\bigcap_{n=0}^{d-1}\Ker(A^0(B^T)^n) = \{0\}
\end{equation*}
where %%
%⭡ ⭡ ⭡ END of annotation 54
with their notation we have $A^0 = \begin{pmatrix}0&0\\0&I_3\end{pmatrix}$ and $B = \begin{pmatrix}0&4\Sigma^T\Sigma\\-2H&\tilde B\end{pmatrix}$, we recall $H$ denotes the Hessian of $V$ at $\s=0$ and here $\tilde B = \p_v(b\wedge\Sigma^T\Sigma v)(0)$. Now, let $(x,v)\in\Ker(A^0)\cap\Ker(A^0B^T)$. From the first kernel, we have that $v=0$, hence the second one gives us $4\Sigma^T\Sigma x = 0$ which implies $x=0$ since $\Sigma$ is invertible.

For (Hypo), we need to look at the dynamic they denoted $c^0(e^{tb^0\cdot\nabla}(x,v))$ for $x$ outside a neighborhood of $E = \{b^0 = 0\}\cap\{c^0 = 0\}$. In our case, we have
\begin{equation*}
b^0 = \begin{pmatrix}
    4\Sigma^T\Sigma v\\
     - 2\p_xV + b(x)\wedge \Sigma^T\Sigma v
\end{pmatrix}
\end{equation*}
and $c^0 = 4|\Sigma^T\Sigma v|^2$.

Let $(x_0,v_0)$ outside a neighborhood of $E$. We consider two cases. Either $v_0\neq0$, then around $(x_0,v_0)$, $c^0$ is uniformly bounded from below, and we obtain that (Hypo) is satisfied. Or $v_0 = 0$ and necessarily, $x_0$ is far enough from a critical point of $V$, therefore $e^{tb^0\cdot\nabla}(x_0,v_0) = e^{- 2t\p_xV(x_0)\cdot\p_v}(x_0,v_0)$. Thus for small time $t>0$, this flow pushes $(x_0,v_0)$ outside a neighborhood of $\R^3_x\times\{0\}$. We can then use the uniform lower bound on $c^0$ to conclude.

This allows us to resolve \eqref{eq:eik} and \eqref{eq:transp} exactly the same way as in \cite{BoLePMi22} since they do not require any assumption at infinity for these local constructions. Their construction of $\ell$ gives the proof of Proposition \ref{prop:ell} $i)$, $ii)$ and $iii)$. For $iv)$ %%
%% Annotation 53, 1/2 on page 39 [ ]
%% Highlight: "Lemma 4.1].  <SEL>Remark 4.4. In [5], the slow growth assumptions of (1.9) are useful to determine resolvent estimates and rough spectrum localization. If one manage to prove such results without their method, for example using the hypocoercive estimates of Section 2, then satisfying only (Harmo) and (Hypo) of [5] is enough to apply their sharp construction of Gauss</SEL>- <SEL>ian quasimodes, resulting in a precise description and Eyring-Kramers law for the bottom of the spectrum of the operator.</SEL>  "
%% Comment: "gobally fix remark style to italic lightface and upright number"
%%
%⭣ ⭣ ⭣
see \cite[Lemma 4.1]{BoLePMi22}.

\begin{remark}
    In \cite{BoLePMi22}, the slow growth assumptions of $(1.9)$ are useful to determine resolvent estimates and rough spectrum localization. If one manage to prove such results without their method, for example using the hypocoercive estimates of Section \ref{sec:Hypocoer}, then satisfying only (Harmo) and (Hypo) of \cite{BoLePMi22} is enough to apply their sharp construction of Gaussian quasimodes, resulting in a precise description and Eyring-Kramers law for the bottom of the spectrum of the operator.
\end{remark}











































\section{Global construction}\label{sec:global}

To %%
%⭡ ⭡ ⭡ END of annotation 53
construct proper quasimodes, we need the notions introduced in Definition \ref{def:labeling} and the labeling given after this. We shall also suppose the generic assumption \eqref{ass.gener} holds true, and we recall it
\begin{equation*}\tag{Gener}
\begin{array}{l}
    (\ast)\mbox{ for any }\m\in\uuu^{(0)},\m \mbox{ is the unique global minimum of } V_{|E(\m)},\\
    (\ast)\mbox{ for all }\m\neq\m'\in\uuu^{(0)},\j(\m)\cap\j(\m')=\emptyset.
\end{array}
\end{equation*}
Given $\m\in\uuu^{(0)}\setminus\{\um\}$, one has $\bsigma(\m) = \sigma_i$ for a certain $i\geq2$. Hence, since $\sigma_{i-1} > \sigma_i$, there exists a unique connected component of $\{V<\sigma_{i-1}\}$ containing $\m$, we denote $E_-(\m)$ that set.

Then we follow the construction of \cite[Section 4]{BoLePMi22}. Notice that in our setting of a kinetic operator, we have two ways of constructing the geometric setup, either by defining the objects on $\R^d_x$ and then tensorizing them by $\R^{d'}_v$, or directly defining the objects on $\R^{d+d'}$. Here we describe the construction on $\R^{d+d'}$. We refer to \cite[Section 3.1]{No23} for a very thorough description of these constructions and a justification of the equivalence between this method and the tensorization one, justifying Definition \ref{def:labeling} can extend from $V$ to $f$.

We recall $\um$ is the unique global minimum of $f$ (the uniqueness is implied by \eqref{ass.gener}). Now let us consider some arbitrary $\m\in\uuu^{(0)}\setminus\{\um\}$. For every $\s\in\j(\m)$, for any $\tau,\delta>0$, we define the sets $\bbb_{\s,\tau,\delta}$, $\ccc_{\s,\tau,\delta}$ and $E_{\m,\tau,\delta}$ by
\begin{equation*}
\bbb_{\s,\tau,\delta} = \{f\leq f(\s)+\delta\}\cap\{X\in\R^{d+d'},|\eta(\s)\cdot(X - \s)|\leq\tau\},
\end{equation*}
\begin{equation}\label{eq:ccc}
\ccc_{\s,\tau,\delta} \mbox{ the connected component of } \bbb_{\s,\tau,\delta} \mbox{ containing } \s
\end{equation}
and
\begin{equation*}
E_{\m,\tau,\delta} = \big(E_-(\m)\cap\{f<f(\j(\m))+\delta\}\big)\setminus\bigcup_{\s\in\j(\m)}\ccc_{\s,\tau,\delta},
\end{equation*}
where $\eta(\s) = \nabla\ell_{\s,0}(\s)\neq0$ with $\ell_{\s}$ defined in Proposition \ref{prop:ell}. We recall Proposition \ref{prop:ell} $iv)$,
\begin{equation*}
\forall X\in\www\setminus\{\s\},\ \ \ X-\s\in\eta(s)^\bot = 0 \Rightarrow f(X) > f(\s).
\end{equation*}
For $\tau,\delta>0$ small enough, this leads to a partition of $E_{\m,\tau,\delta} = E^+_{\m,\tau,\delta}\sqcup E^-_{\m,\tau,\delta}$ where $E^+_{\m,\tau,\delta}$ is defined as the connected component of $E^+_{\m,\tau,\delta}$ containing $\m$. Let us notice that any path connecting $E^+_{\m,\tau,\delta}$ to $E^-_{\m,\tau,\delta}$ within $\{f<f(\j(\m))+\delta\}$ shall cross at least one $\ccc_{\s,\tau,\delta}$ because of Definition \ref{def:labeling} of the separating saddle points.

We can now define, for $h>0$ and $\tau,\delta$ small enough, the function $\theta_\m$ on the sublevel set $E_-(\m)\cap\{f<f(\j(\m)) + 3\delta\}$ as follows. On the disjoint open sets $E^+_{\m,3\tau,3\delta}$ and $E^-_{\m,3\tau,3\delta}$, we define
\begin{equation*}
\theta_\m(X)=\left\{\begin{aligned}
    1\mbox{ for } X\in E^+_{\m,3\tau,3\delta},\\
    -1\mbox{ for } X\in E^-_{\m,3\tau,3\delta}.
\end{aligned}\right.
\end{equation*}
In addition, for $X\in\ccc_{\s,3\tau,3\delta}$ we set
\begin{equation*}
\theta_\m(X) = C_{\s,h}^{-1}\int_0^{\ell_\s(X)}\zeta(r/\tau)e^{-\frac{r^{2}}{2 h}}dr,
\end{equation*}
where the function $\ell_\s$ is the one constructed in the previous section and its sign (see Remark \ref{rem:signell}) is chosen so that there exists a neighborhood $\www$ of $\s$ such that $E(\m)\cap\www$ is included in the half plane $\{\eta(\s)\cdot(X - \s)>0\}$. We recall that $\zeta\in\ccc_c^\infty(\R,[0,1])$ is even and satisfies $\zeta = 1$ on $[-1,1]$ and $\zeta=0$ outside $[-2,2]$, and we set the normalizing constant
\begin{equation*}
C_{\s,h}=\frac 12\int_{-\infty}^{+\infty}\zeta(r/\tau)e^{-\frac{r^{2}}{2h}}dr.
\end{equation*}
Therefore, since, for every $\tau>0$ and then $\delta>0$ small enough, the sets $E^+_{\m,3\tau,3\delta}$, $E^-_{\m,3\tau,3\delta}$ and $\ccc_{\s,3\tau,3\delta}$ for $\s\in\j(\m)$ are mutually disjoint, $\theta_\m$ is well defined on their reunion which forms $E_-(\m)\cap\{V<V(\j(\m))+3\delta\}$. Moreover, on a small neighborhood of the common boundary between $\ccc_{\s,3\tau,3\delta}$ and $E_{\m,3\tau,3\delta}$, we have that $|\eta(\s)\cdot(X - \s)|\geq\frac52\tau$ or in other words, $|\ell_{\s,0,1}|\geq\frac52\tau$. Using that $\ell_\s = \ell_{\s,0,1} + O(|x-\s|^2 + h)$ and having that $|x-\s| = O(\delta)$ in this neighborhood, we thus obtain that for every $\tau > 0$ small and then $\delta,h>0$ small enough, $|\ell_\s| \geq 2\tau$ in a neighborhood of the boundary between $\ccc_{\s,3\tau,3\delta}$ and $E_{\m,3\tau,3\delta}$. This shows that $\theta_\m$ is $\ccc^\infty$ on $E_-(\m)\cap\{V<V(\j(\m))+3\delta\}$.

Note also that there exists $\gamma,\eps>0$ such that
\begin{equation*}
\begin{aligned}
    C_{\s,h} &= \frac 12\int_{-\infty}^{+\infty}\zeta(r/\tau)e^{-\frac{r^{2}}{2h}}dr = \int_0^{+\infty}\zeta(r/\tau)e^{-\frac{r^{2}}{2h}}dr\\
    &= \int_0^{+\infty}e^{-\frac{r^{2}}{2h}}dr + \int_0^{+\infty}(\zeta(r/\tau)-1)e^{-\frac{r^{2}}{2h}}dr\\
    &= \sqrt{\frac{\pi h}{2}} + O\big(\int_\gamma^{+\infty}e^{-\frac{r^{2}}{2h}}dr\big)\\
    &= \sqrt{\frac{\pi h}{2}}(1+O(e^{-\eps/h})).
\end{aligned}
\end{equation*}
Hence
\begin{equation}\label{eq:ch}
\exists\eps>0,\ C_{\s,h}^{-1} = \sqrt{\frac{2}{\pi h}}(1+O(e^{-\eps/h})).
\end{equation}

We now want to extend $\theta_\m$ to a cutoff defined on $\R^{d+d'}$. Considering a smooth function $\chi_\m$ such that
\begin{equation*}
\chi_\m(X) = \left\{\begin{aligned}
    &1\mbox{ for } X\in E_-(\m)\cap\{f \leq f(\j(\m)) + 2\delta\},\\
    &0\mbox{ for } X\in\R^{d+d'}\setminus\big(E_-(\m)\cap\{f < f(\j(\m)) + 3\delta\}\big),
\end{aligned}\right.
\end{equation*}
we have that $\chi_\m\theta_\m$ belongs to $\ccc_c^\infty(\R^{d+d'},[-1,1])$ and
\begin{equation*}
\supp(\chi_\m\theta_\m)\subset E_-(\m)\cap\{f < f(\j(\m)) + 3\delta\}.
\end{equation*}

\begin{defin}\label{def:quasimodes}
For $\tau>0$ and then $\delta,h>0$ small enough, we define the quasimodes 
\begin{equation*}
\left\{\begin{array}{l}
    \psi_{\um}(X)=2e^{-\frac{f(X)-f(\um)}{h}}\\
    \psi_{\m}(X)=\chi_\m(X)(\theta_\m(X)+1)e^{-\frac{f(X)-f(\m)}{h}}\quad\mbox{ for }\m\in\uuu^{(0)}\setminus\{\um\}.
\end{array}\right.
\end{equation*}
And at the same time, we define the normalized quasimodes for $\m\in\uuu^{(0)}$ by
\begin{equation*}
\phii_{\m}=\frac{\psi_{\m}}{\vvert{\psi_{\m}}}.
\end{equation*}
\end{defin}

From this definition, we have the following Lemma which gives us a first relationship between the quasimodes
\begin{lemma}\label{lem:supppsi}Let $\m\neq\m'\in\uuu^{(0)}$,

    If $\bsigma(\m) = \bsigma(\m')$ and $\j(\m)\cap\j(\m')=\emptyset$ then $\supp(\psi_\m)\cap\supp(\psi_{\m'}) = \emptyset$.

    If $\bsigma(\m) > \bsigma(\m')$, then
    \begin{itemize}
        \item[$\star$] either $\supp(\psi_\m)\cap\supp(\psi_{\m'}) = \emptyset$,
        \item[$\star$] or $\psi_\m = 2e^{-(f-f(\m))/h}$ on $\supp(\psi_{\m'})$.
    \end{itemize}
\end{lemma}

The proof is the same as for \cite[Lemma 4.4]{BoLePMi22}. We recall it for the reader's convenience.

\bp Because $E(\m)$ is a connected component of $\{f<\bsigma(\m)\}$, the boundary of $\overline{E(\m)}$ is made of non-critical points of $f$ (points where $\nabla f\neq0$) and separating saddle points $\s\in\j(\m)$ by the definition of a separating saddle point. Therefore, from the definition of $\psi_\m$, we see that for all $\eps>0$,
\begin{equation*}
\supp\psi_\m\subset \overline{E^+_{\m,3\tau,3\delta}\cup\bigcup_{\s\in\j(\m)}\ccc_{\s,3\tau,3\delta}}\subset\overline{E(\m)}+B(0,\eps)
\end{equation*}
for $\tau,\delta$ small enough. Now if $\bsigma(\m) = \bsigma(\m')$ then necessarily $E(\m)\cap E(\m') = \emptyset$. If it was not the case, because they are critical components of $\{f\leq\bsigma(\m)\}$, they would be the same which is in a contradiction with the construction of $E$. When in addition $\j(\m)\cap\j(\m')=\emptyset$, then
\begin{equation*}
\overline{E(\m)}\cap\overline{E(\m')} = \p E(\m)\cap\p E(\m') = \j(\m)\cap\j(\m') = \emptyset,
\end{equation*}
hence with $\eps$ sufficiently small we have $\supp(\psi_\m)\cap\supp(\psi_{\m'}) = \emptyset$. If $\bsigma(\m) > \bsigma(\m')$ then either $\m'\notin E(\m)$ in which case we have with the above that $\supp(\psi_\m)\cap\supp(\psi_{\m'}) = \emptyset$. And if $\m'\in E(\m)$ then $\overline{E(\m')}\subset E_-(\m')\subset E(\m)$ but $\chi_\m\theta_\m\equiv1$ on a neighborhood of $\overline{E(\m)}\setminus\bigcup_{\s\in\j(\m)}\ccc_{\s,3\tau,3\delta}$ and $\j(\m)\cap\supp\psi_{\m'}=\emptyset$ hence the last result.

\ep

We denote
\begin{equation}\label{eq:defmu}
\mu(\m) = 1 + \inf_{\s\in\j(\m)}b_{\ell_\s},
\end{equation}
where $b_{\ell_\s}$ is defined in \eqref{eq:defab}. It is a power that will appear in the following, and we denote $\tilde\j(\m)\subset\j(\m)$ the set of saddle points that satisfy the infimum, in other words
\begin{equation}\label{eq:deftildej}
\s_0\in\tilde\j(\m) \iff b_{\ell_\s} = \min_{\s\in\j(\m)}b_{\ell_\s}.
\end{equation}

\begin{proposition}\label{prop:interactionmatrix}
Suppose that Proposition \ref{prop:ell} holds, under Assumption \eqref{ass.gener}, there exists $C>0$ such that for $\tau>0$ and then $\delta,h>0$ small enough, for every $\m,\m'\in\uuu^{(0)}$,
\begin{itemize}
    \item[$i)$] $\<\phii_\m,\phii_{\m'}\> = \delta_{\m,\m'}+O(e^{-C/h}),$
    \item[$ii)$] $\<P\phii_\m,\phii_\m\> = \frac1{2\pi}\sum_{\s\in\tilde\j(\m)}\frac{(\det\Hess_\m f)^{\frac12}}{|\det\Hess_\s f|^{\frac12}}a_{\ell_\s}h^{\mu(\m)}e^{-2S(\m)/h}(1+O(h))$%%
%% Annotation 55, 1/1 on page 43 [ ]
%% Text: "None, annot.type = <Text>"
%% Comment: "adjust to fit"
%%
%⭣ ⭣ ⭣
, with %%
%⭡ ⭡ ⭡ END of annotation 55
$\mu$ defined in \eqref{eq:defmu} and $a_{\ell_\s}$ defined in \eqref{eq:defab}. Where $S(\m) = f(\j(\m))-f(\m) = V(\j(\m))-V(\m)$ for $\m\neq\um$ and $S(\um)=+\infty$ as denoted in \eqref{eq:defS}.
    \item[$iii)$] $\vvert{P\phii_\m}^2 = O(h^\infty)\<P\phii_\m,\phii_\m\>$
    \item[$iv)$] $\vvert{P^*\phii_\m}^2 = O(h^{-c})\<P\phii_\m,\phii_\m\>$ for some $c\in\R$.
\end{itemize}
\end{proposition}

\bp We will follow the proof of \cite[Proposition 5.1]{BoLePMi22} and will not explain every arguments that did not change from their proof. In this proof we will use several Laplace method (LM), they are all justified thanks to Proposition \ref{prop:ell} $iii)$.

Noticing $f$ uniquely attains its global minimum at $\m$ on $\supp\psi_\m$, by using a LM applied to $2f$ we obtain for $\m\neq\um$
\begin{equation*}
\begin{aligned}
    \vvert{\psi_\m}^2 &= \int_{\R^{d+d'}}\chi_\m^2(\theta_\m+1)^2e^{-2\frac{f-f(\m)}{h}}\\
    &= \chi_\m^2(\m)(\theta_\m(\m)+1)^2\frac{(h\pi)^{\frac {d+d'}2}}{(\det\Hess_\m f)^{\frac12}}(1+O(h)).
\end{aligned}
\end{equation*}
This leads to
\begin{equation}\label{eq:eclpsi}
\vvert{\psi_\m} = 2\frac{(h\pi)^{\frac {d+d'}4}}{(\det\Hess_\m f)^{\frac14}}(1+O(h)).
\end{equation}
Now, notice that this result also holds for $\m = \um$. Then, the proof of $i)$ is exactly the same as in \cite{BoLePMi22} using Lemma \ref{lem:supppsi}.





For $ii)$, Using the computation done after the proof of \cite[Proposition 5.1 $i)$]{BoLePMi22} of $\<P\psi_\m,\psi_\m\>$ we have 
\begin{equation*}
\begin{aligned}
    \<P\psi_\m,\psi_\m\> &= h^2\sum_{\s\in\j(\m)}C_{\s,h}^{-2}\int_{\ccc_{\s,3\tau,3\delta}}\chi_\m^2\zeta(\ell_\s/\tau)^2|\p_v\ell_\s|^2e^{-2\big(f+\frac{\ell_\s^{2}}{2}-f(\m)\big)/h}\\&\phantom{***************}+O(e^{-2(S(\m)+2\delta)/h}).
\end{aligned}
\end{equation*}
Thanks to Proposition \ref{prop:ell} $iii)$, we have that $\s$ is the unique minima of $f+\frac{\ell^{2}_{\s,0}}{2s}-f(\m)$ on $\ccc_{\s,3\tau,3\delta}$ and
\begin{equation*}
\big(f+\frac{\ell^{2}_{\s,0}}{2}-f(\m)\big)(\s) = f(\s) - f(\m) = f(\j(\m)) - f(\m) = S(\m).
\end{equation*}
Moreover, thanks to Remark \ref{rem:relax} and Proposition \ref{prop:ell} $iii)$, $f+\frac{\ell^{2}_{\s,0}}{2}$ satisfies Assumption \ref{ass.morse}, thus we can use a LM and we obtain
\begin{equation*}
\begin{aligned}
    \<P\psi_\m,\psi_\m\> &= h^2\sum_{\s\in\j(\m)}C_{\s,h}^{-2}\chi_\m^2(\s)\zeta(\ell_{\s,0}(\s)/\tau)^2|\p_v\ell_{\s,0}(\s)|^2\frac{(h\pi)^{\frac {d+d'}2}}{(\det\Hess_\s f)^{\frac12}}\\
    &\phantom{************}\times e^{-2S(\m)/h}(1+O(h))\\
\end{aligned}
\end{equation*}
in the case where $|\p_v\ell_{\s,0}(\s)|^2\neq0$. Otherwise, thanks to Proposition \ref{prop:ell} and a LM, we replace $|\p_v\ell_{\s,0}(\s)|^2$ by $h\div(A\nabla f)(\s)$, with $A$ given by Proposition \ref{prop:ell} (we have an extra factor $2$ because we use the LM with $\phii = 2f$). For the example of Subsection \ref{ssec:nolin}, this gives $-{2}{\p_xV(\s)}h$ (recall $d=d'=1$ in this case). In the following, we will use
\begin{equation}\label{eq:defab}
\begin{aligned}
    &a_{\ell_\s} = |\p_v\ell_{\s,0}(\s)|^2 \text{ and } b_{\ell_\s} = 0 \text{ if } |\p_v\ell_{\s,0}(\s)|^2 \text{ is non-zero, and}\\
    &a_{\ell_\s} = \div(A\nabla f)(\s) \text{ and } b_{\ell_\s} = 1 \text{ otherwise.}
\end{aligned}
\end{equation}
Therefore we can write
\begin{equation*}
\begin{aligned}
    \<P\psi_\m,\psi_\m\> &= h^2\sum_{\s\in\j(\m)}C_{\s,h}^{-2}\chi_\m^2(\s)\zeta(\ell_{\s,0}(\s)/\tau)^2a_{\ell_\s}h^{b_{\ell_\s}}\frac{(h\pi)^{\frac {d+d'}2}}{(\det\Hess_\s f)^{\frac12}}\\
    &\phantom{************}\times e^{-2S(\m)/h}(1+O(h)).\\
\end{aligned}
\end{equation*}
Thanks to \eqref{eq:ch}, this leads to 
\begin{equation}\label{eq:Ppsipsi}
\begin{aligned}
    \<P\psi_\m,\psi_\m\> &= \sum_{\s\in\j(\m)}a_{\ell_\s}\frac{2(h\pi)^{\frac {d+d'}2}}{\pi(\det\Hess_\s f)^{\frac12}}h^{1+b_{\ell_\s}}e^{-2S(\m)/h}(1+O(h))\\
    &= \sum_{\s\in\tilde\j(\m)}a_{\ell_\s}\frac{2\pi^{\frac {d+d'}2-1}}{(\det\Hess_\s f)^{\frac12}}h^{\mu(\m) + \frac {d+d'}2}e^{-2S(\m)/h}(1+O(h)).
\end{aligned}
\end{equation}
Combining \eqref{eq:Ppsipsi} and \eqref{eq:eclpsi}, we obtain
\begin{equation*}
\<P\phii_\m,\phii_\m\> = \frac1{2\pi}\sum_{\s\in\tilde\j(\m)}\frac{(\det\Hess_\m f)^{\frac12}}{|\det\Hess_\s f|^{\frac12}}a_{\ell_\s}h^{\mu(\m)}e^{-2S(\m)/h}(1+O(h)),
\end{equation*}
this proves $ii)$.





Let us now prove $iii)$. Starting as in \cite{BoLePMi22}'s work, we have that
\begin{equation}\label{eq:Dvpsim}
\vvert{P\psi_\m}^2 = \vvert{P(\theta_\m e^{-(f-f(\m))/h})}^2_{L^2(\supp\chi_\m)} + O(e^{-2(S(\m)+2\delta))/h}).
\end{equation}
And on $\supp\chi_\m$, $P(\theta_\m e^{-(f-f(\m))/h})$ is supported in $\bigcup_{\s\in\j(\m)}\ccc_{\s,3\tau,3\delta}$, thus using \eqref{eq:Pchiell}, we obtain with a %%
%% Annotation 57, 2/2 on page 45 [ ]
%% Text: "None, annot.type = <Text>"
%% Comment: "adjust to fit"
%%
%⭣ ⭣ ⭣
LM
\begin{equation*}
\begin{aligned}
    \vvert{P(\theta_\m e^{-(f-f(\m))/h})}^2_{L^2(\ccc_{\s,3\tau,3\delta})} &= \int_{\ccc_{\s,3\tau,3\delta}} O(X^\infty + h^\infty)e^{-2\big(f-f(\m)+\frac{\ell_\s^{2}}{2}\big)/h}\\
    &= O(h^\infty)e^{-2S(\m)/h}\\
    &= O(h^\infty)\<P\psi_\m,\psi_\m\>_{L^2(\ccc_{\s,3\tau,3\delta})}.
\end{aligned}
\end{equation*}
Combining %%
%⭡ ⭡ ⭡ END of annotation 57
this results and \eqref{eq:Dvpsim}, we proved point $iii)$.





For $iv)$, notice that around $\s\in\j(\m)$%%
%% Annotation 56, 1/2 on page 45 [ ]
%% Highlight: "point iii).  <SEL>2 )/h</SEL>  For iv), notice that around s ∈j(m), <SEL>P ∗(θme−f/h) = h( ˆw+r)e−(f+ ℓ2s</SEL>  with  "
%% Comment: "display"
%%
%⭣ ⭣ ⭣
, $P^*(\theta_\m e^{-f/h}) = h(\hat w + r)e^{-(f+\frac{\ell_\s^{2}}{2})/h}$ with %%
%⭡ ⭡ ⭡ END of annotation 56
\begin{equation*}
\hat w = - \alpha\cdot\p_x\ell_\s - \beta\cdot\p_v\ell_\s + 4\Sigma^T\Sigma v\cdot\p_v\ell_\s + \ell_\s|\p_v\ell_\s|^2 - h\Delta_v\ell_\s,
\end{equation*}
therefore with a LM and using the analog of \eqref{eq:Dvpsim} with $P^*$ instead of $P$, we obtain the result.

\ep

\subsection{Proof of Theorem \ref{thm:2} and graded matrices}\label{ssec:proofthm2}
The proof is the same as in \cite{BoLePMi22} we refer to it and to \cite{LePMi20} for the details, let us mention the main arguments. From now on, we relabel the minima by increasing saddle height:
\begin{equation}\label{eq:decreasSmu}
\forall j\in\lb1,n_0-1\rb,\left\{\begin{aligned}
    &S(\m_j) \leq S(\m_{j+1}),\\
    &\mu(\m_{j+1}) \leq \mu(\m_j) \text{ if }S(\m_j) = S(\m_{j+1})
\end{aligned}\right.
\end{equation}
Thus $\m_{n_0} = \um$ because $S(\um) = +\infty$ and $\um$ is the only minima that has this property by construction. We also denote for shortness
\begin{equation*}
\forall j\in\lb1,n_0\rb,\ \ S_j = S(\m_j),\ \ \phii_j = \phii_{\m_j},\ \mbox{ and } \ \tilde{\lambda}_j = \<P\phii_j,\phii_j\>.
\end{equation*}
Therefore we have
\begin{equation}\label{eq:Pmm'}
\forall j,k\in\lb1,n_0\rb,\ \ \<P\phii_j,\phii_k\> = \delta_{j,k}\tilde{\lambda}_j.
\end{equation}
This statement is obvious when $j=k$. Now if $j > k$, from Lemma \ref{lem:supppsi} either $\supp\phii_j\cap\supp\phii_k = \emptyset$ or $\phii_j = c_he^{-(f-f(\m_j)/h}$ on $\supp\phii_k$, $c_h$ being a normalizing constant (or the same swapping $j$ and $k$). Using that $P(e^{-f/h}) = P^*(e^{-f/h}) = 0$, we see that \eqref{eq:Pmm'} is indeed true. Let $\Cr = \partial D(0,\frac{c_0}2g(h))$. We now define
\begin{equation*}
\Pi_\Cr = \frac{1}{2i\pi}\int_{\Cr}(z-P)^{-1}dz
\end{equation*}
the spectral projector on the small eigenvalues of $P$, where $c_0$ is given by Theorem \ref{thm:1}, $g(h)$ is defined in \eqref{eq:g}. Observe that thanks to Theorem \ref{thm:1}, $\vvert{\Pi_\Cr}\leq \frac{c_0}{c_0'}$. Hence, denoting $u_j = \Pi_\Cr\phii_j$ for $j\in\lb1,n_0\rb$ (we notice that $u_{n_0} = \phii_{\um}$), we obtain the following proposition

\begin{proposition}
    Under Assumption \ref{ass.polyg}, there exists $c>0$ such that for every $j,k\in\lb1,n_0\rb$ and every $h>0$ small enough, one has
    \begin{equation}\label{eq:quasiortho}
    \<u_j,u_k\> = \delta_{j,k} + O(e^{-c/h})
    \end{equation}
    and
    \begin{equation}\label{eq:interaction}
    \<Pu_j,u_k\> = \delta_{j,k}\tilde{\lambda}_j + O\Big(h^\infty\sqrt{\tilde{\lambda}_j\tilde{\lambda}_k}\Big).
    \end{equation}
\end{proposition}

\bp Using that
\begin{equation*}
\Pi_\Cr-1 = \frac{1}{2i\pi}\int_{\Cr}(z-P)^{-1}P\frac{dz}{z}
\end{equation*}
we have that for $u\in D(P)$, $\vvert{(\Pi_\Cr-1)u}\leq\sup_{\Cr}\Vert(z-P)^{-1}\Vert\vvert{Pu}$, and thus
\begin{equation*}
\<u_j,u_k\> = \<\phii_j,\phii_k\> + g(h)^{-1}O\big(\Vert P\phii_j\Vert + \Vert P\phii_k\Vert\big)
\end{equation*}
using the resolvent estimate given by Theorem \ref{thm:1}. Using now Assumption \ref{ass.polyg}, we obtain $\<u_j,u_k\> = \delta_{j,k} + O(e^{-c/h})$ thanks to Proposition \ref{prop:interactionmatrix}. Therefore, thanks to Theorem \ref{thm:1} and \eqref{eq:Pmm'}, we then %%
%% Annotation 58, 1/1 on page 47 [ ]
%% Caret: " δj,k˜λj + O<Caret>  □  "
%% Comment: "\qedhere "
%%
%⭣ ⭣ ⭣
obtain
\begin{equation*}
\begin{aligned}
    \<Pu_j,u_k\> &= \<P\phii_j,\phii_k\> + \<P(\Pi_\Cr-1)\phii_j,\phii_k\> + \<P\Pi_\Cr\phii_j,(\Pi_\Cr-1)\phii_k\>\\
    &= \delta_{j,k}\tilde{\lambda}_j + g(h)^{-1}O\big(\Vert P\phii_j\Vert\Vert P^*\phii_k\Vert + \Vert P\phii_j\Vert\Vert P\phii_k\Vert\big)\\
    &= \delta_{j,k}\tilde{\lambda}_j + O\Big(h^\infty\sqrt{\tilde{\lambda}_{j}\tilde{\lambda}_{k}}\Big).
\end{aligned}
\end{equation*}

\ep

This %%
%⭡ ⭡ ⭡ END of annotation 58
grants to the interaction matrix a graded structure we will develop below. Then we use the Gram-Schmidt process to transform the basis $(u_{n_0-j+1})_{1\leq j\leq n_0}$ into an orthonormal basis $(e_{n_0-j+1})_{1\leq j\leq n_0}$ of $\Ran\Pi_\Cr$. Moreover, thanks to \eqref{eq:quasiortho} we have that
\begin{equation*}
\forall j\in\lb1,n_0\rb,\ \ e_j = u_j + O(e^{-c/h}),
\end{equation*}
see \cite[Lemma 4.11]{LePMi20} for the details. And thus, using \eqref{eq:interaction} along with \cite[Proposition 4.12]{LePMi20}, we have that
\begin{equation}\label{eq:Pee}
\forall j\in\lb1,n_0\rb,\ \ \<Pe_j,e_k\> = \delta_{j,k}\tilde{\lambda}_j + O\Big(h^\infty\sqrt{\tilde{\lambda}_j\tilde{\lambda}_k}\Big).
\end{equation}
Using its graded structure, we can now compute the eigenvalues of the matrix
\begin{equation}\label{eq:matrix}
M := (\<Pe_j,e_k\>)_{j,k} = {P}_{|\Ran\Pi}.
\end{equation}
We recall the results stated in \cite{LePMi20}:

We denote by $\Dr_0(E)$ the set of invertible and diagonalizable complex matrices of an Euclidean space $E$.

\begin{defin}\cite[Definition A.1]{LePMi20}
    Let $\Er = (E_j)_{1\leq j \leq p}$ be a sequence of vector spaces $E_j$ of (finite) dimension $r_j>0$, let $E = \bigoplus_{j=1}^pE_j$ and let $\tau = (\tau_i)_{2\leq i\leq p}\in(\R_+^*)^{p-1}$. Suppose that $(h,\tau)\mapsto \mmm_h(\tau)$ is a map from $(0,1]\times(\R_+^*)^{p-1}$ to the set of complex matrices on $E$.

    We say that $\mmm_h(\tau)$ is an $(\Er,\tau,h)$-graded matrix if there exists $\mmm'\in\Dr_0(E)$ independent of $(h,\tau)$ such that $\mmm_h(\tau) = \Omega(\tau)(\mmm' + O(h))\Omega(\tau)$ with $\Omega(\tau)$ and $\mmm'$ such that
    \begin{itemize}
        \item $\mmm' = \diag(M_j,1\leq j\leq p)$ with $M_j\in\Dr_0(E_j)$,
        \item $\Omega(\tau) = \diag(\eps_j(\tau)I_{r_j},1\leq j\leq p)$ with $\eps_1(\tau) = 1$ and $\eps_j(\tau) = \prod_{k=2}^j\tau_k$ for  $j\geq2$.
    \end{itemize}
\end{defin}

\begin{theorem}\cite[Theorem A.4]{LePMi20}\label{thm:graded}
    Suppose that $\mmm_h(\tau)$ is $(\Er,\tau,h)$-graded. Then, there exists $\tau_0, h_0>0$ such that for all $0<\tau_j<\tau_0$ and $h\in(0,h_0]$, one has
    \begin{equation*}
    \sigma(\mmm_h(\tau))\subset\bigsqcup_{j=1}^p\eps_j(\tau)^2(\sigma(M_j) + O(h)).
    \end{equation*}
    Moreover, for any eigenvalue $\lambda$ of $M_j$ with multiplicity $m_j(\lambda)$, there exists $K>0$ such that, denoting $D_j(\lambda) = \{z\in\C,\ |z-\eps_j(\tau)^2\lambda| < \eps_j(\tau)^2Kh\}$, one has
    \begin{equation*}
    n(D_j(\lambda);\mmm_h(\tau)) = m_j(\lambda),
    \end{equation*}
    where $n(D_j;\mmm_h(\tau))$ is the rank of the Riesz projector associated with $\mmm_h(\tau)$ with contour $D_j$. Finally, there exists $C>0$, such that for all $z\in\C\setminus\bigcup_{j=1}^p\bigcup_{\lambda\in\sigma(M_j)}D_j(\lambda)$,
    \begin{equation}\label{eq:resolventGraded}
    \vvert{(\mmm_h(\tau) - z)^{-1}} \leq Cd(z,\sigma(\mmm_h(\tau)))^{-1}.
    \end{equation}
\end{theorem}

We want to apply this theorem to \eqref{eq:matrix}, we therefore need to show that it is graded. Let us first notice that
\begin{equation*}
M = \begin{pmatrix}M'&0\\0&0\end{pmatrix}
\end{equation*}
with $M'\in\Mr_{n_0-1}(\R)$ because by construction $e_{n_0}$ is the ground state of $P$. Now, for all $j\in\lb1,n_0-1\rb$, let us define
\begin{equation}\label{eq:defvmu}
v_j = \frac1{2\pi}\sum_{\s\in\tilde\j(\m)}\frac{(\det\Hess_\m f)^{\frac12}}{|\det\Hess_\s f|^{\frac12}}a_{\ell_\s},\ \mbox{ and }\ \mu_j = 1 + b_{\ell_\s} = \mu(\m_j)
\end{equation}
for some $\s\in\tilde\j(\m_j)$ (recalling $\tilde \j$ is defined in \eqref{eq:deftildej}), with $a_{\ell_\s}$ and $b_{\ell_\s}$ defined in \eqref{eq:defab}. Therefore we have that
\begin{equation}\label{eq:lamdbav}
\tilde\lambda_j = v_jh^{\mu_j}e^{-2S_j/h}(1 + O(h)).
\end{equation}
Because of \eqref{eq:decreasSmu}, there exists a partition $J_1\sqcup\ldots\sqcup J_p$ of $\lb1,n_0-1\rb$ such that for all $k\in\lb1,p\rb$, there exists $\iota(k)\in\lb1,n_0-1\rb$ such that
\begin{equation*}
\forall j\in J_k,\ \ S_j = S_{\iota(k)},\ \mu_j = \mu_{\iota(k)}
\end{equation*}
and
\begin{equation*}
\forall 1\leq k< k'\leq p,\ \ S_{\iota(k)} < S_{\iota(k')}\ \text{ or }\ S_{\iota(k)} = S_{\iota(k')}\ \text{ and }\  \mu_{\iota(k')} < \mu_{\iota(k)}
\end{equation*}
Hence, using \eqref{eq:Pee} and \eqref{eq:lamdbav} we have that $\big(h^{\mu_1}e^{-2S_1/h}\big)^{-1}M'$ is $(\Er,\tau,h)$-graded with $\Er = (\R^{J_k})_{1\leq k\leq p}$ and $\tau_k = h^{\frac{\mu_{\iota(k)}-\mu_{\iota(k-1)}}{2}} e^{-\frac{S_{\iota(k)}-S_{\iota(k-1)}}{h}}$ for $k\geq 2$. We can now apply Theorem \ref{thm:graded} and we obtain
\begin{equation*}
\sigma(M') \subset \bigsqcup_{k=1}^ph^{\mu_1}e^{-2S_1/h}\eps_k^2(\sigma(M_k) + O(h))
\end{equation*}
with $M_k = \diag(v_j, j\in J_k)$, $h^{\mu_1}e^{-2S_1/h}\eps_k^2 = h^{\mu_{\iota(k)}}e^{-2S_{\iota(k)}/h}$ and corresponding multiplicities. All this leads to the announced result.


































\subsection{Proof of Corollaries \ref{cor:return} and \ref{cor:times}}\label{ssec:proofcor}We follow the proof of \cite[Corollaries 1.5 and 1.6]{BoLePMi22}.

Recall that $\Cr = \partial D(0,\frac{c_0}2g(h))$ and
\begin{equation*}
\Pi_\Cr = \frac{1}{2i\pi}\int_{\Cr}(z-P)^{-1}dz.
\end{equation*}
We already obtained that $\vvert{\Pi_\Cr} \leq \frac{c_0}{c_0'}$. Moreover, recall \eqref{eq:resol} that is
\begin{equation*}
\vvert{(P-z)^{-1}}\leq\frac{2}{c_0'g(h)},
\end{equation*}
for $\{\Re z \leq c_0g(h)\}\cap\{|z|\geq c_0'g(h)\}$. Since the operator valued function $(P-z)^{-1}(1-\Pi_\Cr)$ is holomorphic on $\{\Re z \leq c_0g(h)\}$, the maximum principle yields
\begin{equation}\label{eq:resol2}
\vvert{(P-z)^{-1}(1-\Pi_\Cr)}\leq\frac{2}{c_0'g(h)},
\end{equation}
for $\Re z \leq c_0g(h)$

The solution of \eqref{eq:evolutionP} can be written
\begin{equation}\label{eq:developEvol}
u(t) = e^{-tP/h}u_0 = e^{-tP/h}\Pi_\Cr u_0 + e^{-tP/h}(1-\Pi_\Cr)u_0.
\end{equation}
Let $Q:\Im(1-\Pi_\Cr)\to\Im(1-\Pi_\Cr)$ be the operator $P$ restricted to the Hilbert space $\Im(1-\Pi_\Cr)$. Since $P$ is maximally accretive from Proposition \ref{prop:accretive}, so is $Q$ and thus $\vvert{e^{-tQ/h}}\leq1$. Moreover, \eqref{eq:resol2} shows that $\vvert{(Q-z)^{-1}}\leq Cg(h)^{-1}$ for some $C>0$ and $\Re z \leq c_0g(h)$. To estimate the last term in \eqref{eq:developEvol}, we use a Gearhardt-Pr\"uss type inequality with an explicit bound. More precisely, \cite[Proposition 2.1]{HeSj10_01} (see also \cite{HeSj21_01,HeSjVi24}) gives, for some $C>0$ and all $t\geq0$,
\begin{equation*}
\vvert{e^{-tQ/h}} \leq \Big(1 + 2\frac{c_0}2g(h)\sup_{\Re z = \frac{c_0}2g(h)}\vvert{(Q-z)^{-1}}\Big)e^{-t\frac{c_0}2g(h)/h} \leq Ce^{-t\frac{c_0}2\frac{g(h)}{h}}.
\end{equation*}
Therefore, denoting $\varepsilon = c_0/2$ and using Assumption \ref{ass.polyg}, there exists $C>0$ such that for all $t\geq0$,
\begin{equation}\label{eq:semigroupQ}
\vvert{e^{-tQ/h}(1-\Pi_\Cr)u_0} \leq Ce^{-\varepsilon h^{c-1} t}\vvert{u_0},
\end{equation}
with $c\geq1$ given by Assumption \ref{ass.polyg}. Noticing now that $P_{|\Im\Pi_\Cr}$ is a matrix of size $n_0$ whose eigenvalues are the $\lambda(\m,h)$'s, using \eqref{eq:developEvol} and \eqref{eq:semigroupQ}, we obtain \eqref{eq:return1} from the usual formula for the exponential of a matrix applied to $e^{-tP/h}\Pi_\Cr u_0$.

Let us now prove \eqref{eq:return2}. Denoting $\Pi_0$ the orthogonal projector onto the kernel of $P$, we need to prove
\begin{equation*}
e^{-tP/h} = \Pi_0 + O(e^{-t\underset{\m\neq\um}{\min}\Re(\lambda(\m,h))(1-h^N)/h}),
\end{equation*}
with a $O$ uniform in $t$ and $h$. Writing
\begin{equation*}
e^{-tP/h} - \Pi_0 = e^{-tP/h}\Pi_\Cr - \Pi_0 + e^{-tP/h}(1-\Pi_\Cr),
\end{equation*}
it remains to bound $e^{-tP/h}\Pi_\Cr - \Pi_0$ thanks to \eqref{eq:semigroupQ}. Using that
\begin{equation*}
\Pi_0u = \<\frac{e^{-f/h}}{\vvert{e^{-f/h}}},u\>\frac{e^{-f/h}}{\vvert{e^{-f/h}}},
\end{equation*}
which is a bounded operator, we have $\Pi_0 = e^{-tP/h}\Pi_0$ and it suffices to show that for all $N\in\N$,
\begin{equation*}
\exists C>0,\ \ \vvert{e^{-tP/h}(\Pi_\Cr-\Pi_0)} \leq Ce^{-t\underset{\m\neq\um}{\min}\Re(\lambda(\m,h))(1-Ch)/h},
\end{equation*}
or in other words,
\begin{equation*}
\exists C>0,\ \ \vvert{e^{-tM'/h}} \leq Ce^{-t\underset{\m\neq\um}{\min}\Re(\lambda(\m,h))(1-Ch)/h},
\end{equation*}
with $M'$ the matrix of $P$ in the orthonormal basis $(e_j)_{1\leq j\leq n_0-1}$ of $\Ran(\Pi_\Cr-\Pi_0)$ defined in the preceding subsection. We already showed, using Theorem \ref{thm:graded} that
\begin{equation*}
\sigma(M') \subset \bigsqcup_{\m\in\uuu^{(0)}\setminus\{\um\}} D\big(v(\m)h^{\mu(\m)}e^{-2S(\m)/h},Kh^{\mu(\m)+1}e^{-2S(\m)/h}\big),
\end{equation*}
with $v$ and $\mu$ defined in \eqref{eq:defvmu} and some constant $K>0$. Therefore, using the functional calculus representation of $e^{-tM'/h}$ along with the resolvent estimate of Theorem \ref{thm:graded}, we obtain
\begin{equation*}
\begin{split}
    \vvert{e^{-tM'/h}} &= O\Big(\sup_{\m\neq\um}e^{-tv(\m)h^{\mu(\m)}e^{-2S(\m)/h}(1-\frac K2h)/h}\Big)\\
    &= O\Big(e^{-t\underset{\m\neq\um}{\min}\Re(\lambda(\m,h))(1-\frac K2h)/h}\Big).
\end{split}
\end{equation*}

Let us now prove Corollary \ref{cor:times}. Recall that, in the proof of Theorem \ref{thm:2}, we introduced an application $\iota:\lb1,p+1\rb\to\lb1,n_0\rb$ such that for all $1\leq k<k'\leq p+1$,
\begin{equation*}
h^{\mu_{\iota(k)}}e^{-S_{\iota(k)}/h} > h^{\mu_{\iota(k')}}e^{-S_{\iota(k')}/h}.
\end{equation*}
In the following we will omit the function $\iota$ and just write $\mu_k,S_k$ to avoid heavy expressions.

For $R>1$ we define the balls
\begin{equation*}
\forall k\in\lb1,p\rb,\ \ \ D_k = D\big((R+R^{-1})h^{\mu_k}e^{-2S_k/h},Rh^{\mu_k}e^{-2S_k/h}\big)
\end{equation*}
and $D_{p+1} = D(0,R^{-1}h^{\mu_p}e^{-2S_{p}/h})$. For $R$ fixed large enough and every $h$ small enough, each exponentially small eigenvalue of $P$ belongs to exactly one of the disjoint sets $D_k$ from Theorem \ref{thm:2}. Moreover, $\partial D_k$ is at distance of order $h^{\mu_k}e^{-2S_k/h}$ (resp. $h^{\mu_p}e^{-2S_{p}/h}$) from the spectrum of $P$ for $k\in\lb1,p\rb$ (resp. $k=p+1$). Using the resolvent estimate of $P$ on the image of $\Pi_\Cr$ given by Theorem \ref{thm:graded} and \eqref{eq:resol2} to control the contribution of on the image of $1-\Pi_\Cr$, we get
\begin{equation}\label{eq:resol3}
\forall z\in \partial D_k,\ \ \vvert{(P-z)^{-1}}\leq \left\{
\begin{aligned}
    &Ch^{-\mu_k}e^{2S_k/h}&&\text{for }k\in\lb1,p\rb,\\
    &Ch^{-\mu_p}e^{2S_{p}/h}&&\text{for }k=p+1,
\end{aligned}
\right.
\end{equation}
using Assumption \ref{ass.polyg}. In particular, the spectral projector associated with the eigenvalues of order $h^{\mu_k}e^{-2S_k/h}$,
\begin{equation*}
\Pi_k = \frac{1}{2i\pi}\int_{\partial D_k}(z-P)^{-1}dz,
\end{equation*}
is well-defined and satisfies $\vvert{\Pi_k}\leq C$.

We can now decompose
\begin{equation}\label{eq:decomProj}
e^{-tP/h}\Pi_\Cr = \sum_{k=1}^{p+1}e^{-tP/h}\Pi_k.
\end{equation}
For $k\in\lb1,p\rb$ and $0\leq t\leq t_k^-$, \eqref{eq:resol3} and $t_k^-e^{-2S_k/h} = O(h^\infty)$ %%
%% Annotation 59, 1/1 on page 52 [ ]
%% Highlight: "(5.21) and <SEL>e−t+  k hµk−1e−2Sk/h/R = O(h∞)</SEL> (using that"
%% Comment: "display"
%%
%⭣ ⭣ ⭣
imply
\begin{equation}\label{eq:estim1}
\begin{split}
    e^{-tP/h}\Pi_k &= \frac{1}{2i\pi}\int_{\partial D_k}e^{-tz/h}(z-P)^{-1}dz\\
    &= \frac{1}{2i\pi}\int_{\partial D_k}(z-P)^{-1}dz + \frac{1}{2i\pi}\int_{\partial D_k}(e^{-tz/h} - 1)(z-P)^{-1}dz\\
    &= \Pi_k + \int_{\partial D_k}O(t|z|/h)\vvert{(P-z)^{-1}}dz\\
    &= \Pi_k + O(th^{\mu_k-1}e^{-2S_k/h})\\
    &= \Pi_k + O(h^\infty).
\end{split}
\end{equation}
using that $\mu_k\geq1$. On the contrary, for $t_k^+\leq t$, \eqref{eq:resol3} and $e^{-t_k^+h^{\mu_k-1}e^{-2S_k/h}/R} = O(h^\infty)$ (using that $\mu_k\leq2$ from its definition) give
\begin{equation}\label{eq:estim2}
\begin{split}
    e^{-tP/h}\Pi_k &= \frac{1}{2i\pi}\int_{\partial D_k}e^{-tz/h}(z-P)^{-1}dz\\
    &= O\Big(\int_{\partial D_k}e^{-t\Re z/h}\vvert{(P-z)^{-1}}dz\Big)\\
    &= O\big(e^{-tR^{-1}h^{\mu_k}e^{-2S_k/h}/h}\big)\int_{\partial D_k}\vvert{(P-z)^{-1}}dz\\
    &= O(h^\infty).
\end{split}
\end{equation}
Lastly, %%
%⭡ ⭡ ⭡ END of annotation 59
$e^{-tP/h}\Pi_{p+1} = \Pi_{p+1}$ since $\Pi_{p+1}$ is the rank-one spectral projector associated with the eigenvalue $0$. On the other hand, since $e^{-\varepsilon t_0^+} = O(h^\infty)$, \eqref{eq:semigroupQ} becomes
\begin{equation}\label{eq:estim3}
\vvert{e^{-tP/h}(1-\Pi_\Cr)} = O(h^\infty),
\end{equation}
for $t\geq t_0^+$. Summing up, Corollary \ref{cor:times} is a direct consequence of the formulas \eqref{eq:resol3} and \eqref{eq:decomProj} with the relation $\Pi_k^\leq = \sum_{j=k}^{p+1}\Pi_j$ and the estimates \eqref{eq:estim1}, \eqref{eq:estim2} and \eqref{eq:estim3}.


















\appendix
\section{Labeling of the critical points}

In this section, we consider a smooth function $W\in\ccc^\infty(\R^n,\R)$ satisfying Assumption \ref{ass.morse} that we recall thereafter. We denote by $\uuu$ the set of its critical points.
\begin{assumption}\label{ass:morsedegen2}
    For any critical point $x^*\in\uuu$, there exists a neighborhood $\vvv\ni x^*$, $(t_i^{x^*})_{1\leq i\leq n}\subset \R^{*}$, $(\nu_i^{x^*})_{1\leq i\leq n}\subset \N\setminus\{0,1\}$, a $\ccc^\infty$ change of variable $U^{x^*}$ defined on $\vvv$ such that $U^{x^*}(x^*)=x^*$, $U^{x^*}$ and $d_{x^*}U^{x^*}$ are invertible and
    \begin{equation*}
    \forall x\in\vvv,\ \ W\circ U^{x^*}(x) - W(x^*) = \sum_{i=1}^nt_i^{x^*}(x_i-x_i^*)^{\nu_i^{x^*}}.
    \end{equation*}
\end{assumption}
We consider the partition $\uuu = \uuu^{odd}\sqcup\uuu^{even}$ where $x^*\in\uuu^{even} \iff \forall i,\ \nu_i^{x^*}\in 2\N$. Now we shall say that $x^*\in\uuu^{even}$ is of index $j\in\lb0,n\rb$ if $\sharp\{i\ |\ t_i^{x^*}<0\}=j$ and therefore $\uuu^{even}=\bigsqcup_{j=0}^n\uuu^{(j)}$ where $\uuu^{(j)}$ is the set of critical points of $V$ with even order in each direction and of index $j$, we also denote $n_0 = \sharp\uuu^{(0)}$ the number of minima of $W$. In the following when speaking of saddle point, we only refer to critical points of index $1$, that is elements of $\uuu^{(1)}$, and we will mostly denote them by the letter $\s$ and its variations.

In works in the lineage of \cite{BoLePMi22}, for Morse functions, the important critical points to study are the minima as they generates the eigenvalues and some saddle points because they are the only ones such that $B(\s,r)\cap\{W < W(\s)\}$ has exactly two connected components (for $r>0$ small enough, where $B(\s,r)\subset\R^n$ denotes the open ball of center $\s$ and radius $r$), meaning that if a process starts from one side of this set, it has to cross the level set $\{W = W(\s)\}$ in order to go the other side. One can show that for any other critical point, this set is connected therefore the process need not to climb the landscape given by $W$ in order to cross wells. We need to prove a similar result for functions that are not Morse anymore, but satisfy Assumption \ref{ass:morsedegen2}.

\begin{proposition}\label{prop:saddle}
    Let $\s\in\uuu$, for $r>0$ small enough, the set $B(\s,r)\cap\{W < W(\s)\}$ has exactly two connected components if $\s\in\uuu^{(1)}$ and one otherwise.
\end{proposition}

\subsection{Proof of Proposition \ref{prop:saddle}}

We first need to prove the following lemma, adapted from \cite[Lemma 3.1]{AsBoMi22}, but before that, we define a family of paths. Let $a=(a_1,\ldots,a_n)\in[1,+\infty)^n$, $x,y\in\R^n$, we denote
\begin{equation*}
\begin{array}{c|ccc}
     \gamma^a_{x,y}:&[0,1]&\to&\R^n\\
     &t&\mapsto&(t^{a_i}(x_i-y_i))_i + y.
\end{array}
\end{equation*}
We see that $\gamma^a_{x,y}(0) = y$ and $\gamma^a_{x,y}(1) = x$, and when $a\equiv 1$, this is just the standard linear parametrization of the segment $[y,x]$.
\begin{lemma}\label{lem:path}
    Let $\phii$ be a local smooth diffeomorphism of $\R^n$ defined in a neighborhood of $b\in\R^n$ and $a=(a_1,\ldots,a_n)\in[1,+\infty)^n$. Then there exists $r_b>0$ such that for all $0<r<r_b$,
    \begin{equation*}
    \forall x\in\phii(B(b,r)),\ \gamma^a_{x,\phii(b)}([0,1])\subset\phii(B(b,r)).
    \end{equation*}
\end{lemma}

\bp
In the following, we will denote $\underline a = \inf_ia_i$ and $\overline a = \sup_ia_i$. We have to show that for all $x\in B(0,r)$ and $t\in[0,1]$, the point $\gamma^a_{\phii(b+x),\phii(b)}(t)$ belongs to $\phii(B(b,r))$. In other words,
\begin{equation*}
\forall t\in[0,1],\ g(t) = |\phii^{-1}(\gamma^a_{\phii(b+x),\phii(b)}(t))-b|^2\in[0,r^2).
\end{equation*}
First, we see that $g(0) = 0$ and $g(1) = |x|^2 < r^2$. Then, there exists $E$ a set large enough such that
\begin{equation*}
\begin{aligned}
    g(t) &= |\phii^{-1}(\gamma^a_{\phii(b+x),\phii(b)}(t))-b|^2\\
    &\leq \vvert{d\phii^{-1}}_{L^\infty(E)}^2 |\gamma^a_{\phii(b+x),\phii(b)}(t)-\phii(b)|^2\\
    &= \vvert{d\phii^{-1}}_{L^\infty(E)}^2 |t^{a_i}(\phii(b+x)_i-\phii(b)_i)|^2\\
    &\leq \vvert{d\phii^{-1}}_{L^\infty(E)}^2\vvert{d\phii}_{L^\infty(B(b,r))}^2 t^{2\underline a}|x|^2.
\end{aligned}
\end{equation*}
Thus there exists $C>0$ such that for all $\eps>0$ and all $t\leq\eps$,
\begin{equation*}
g(t) \leq C\eps^{2\underline a}|x|^2.
\end{equation*}
Choosing $\eps$ such that $C\eps^{2\underline a} \leq 1$, we have that $g(t) < r^2$ for $t\in[0,\eps]$.
Moreover, the Taylor formula implies
\begin{equation*}
\begin{aligned}
    g'(t) &= 2\< \p_t(\phii^{-1}\circ\gamma^a_{\phii(b+x),\phii(b)})(t),\phii^{-1}\circ\gamma^a_{\phii(b+x),\phii(b)}(t) - b\>\\
    &= 2\< d_{\gamma^a_{\phii(b+x),\phii(b)}(t)}\phii^{-1}({\gamma^a_{\phii(b+x),\phii(b)}}'(t)),\phii^{-1}\circ\gamma^a_{\phii(b+x),\phii(b)}(t) - b\>\\
    &= 2\< d_{\phii(b)+O(t^{\underline a}x)}\phii^{-1}\big((a_it^{a_i-1}d_b\phii_i(x))_i + O(t^{\underline a-1}x^2)\big),\\&\phantom{************}d_{\phii(b)}\phii^{-1}\big((t^{a_i}d_b\phii_i(x))_i\big) + O(t^{2\underline a}x^2)\>.
\end{aligned}
\end{equation*}
Here we denoted $\phii_i$ the $i$-th component of $\phii$. Denoting $A = d_{\phii(b)}\phii^{-1}$ and $\Gamma_1(t) = \diag(t^{a_i-1})$, $\Gamma_2(t) = \diag(a_it^{a_i-1})$, we observe that
\begin{equation*}
g'(t) = 2t\<A\Gamma_2(t)A^{-1}x,A\Gamma_1(t)A^{-1}x\> + \vvert{A\Gamma_1(t)A^{-1}x}O(t^{\underline a}x^2) + O(t^{2\underline a-1}x^3).
\end{equation*}
But using that $A$ is invertible and thus $A^*A$ is positive definite, we have that
\begin{equation*}
\begin{aligned}
    \<A^*A\Gamma_2(t)A^{-1}x,\Gamma_1(t)A^{-1}x\> &\gtrsim \<\Gamma_2(t)A^{-1}x,\Gamma_1(t)A^{-1}x\>\\
    &\gtrsim t^{2\overline a-2}\<(AA^*)^{-1}x,x\>\\
    &\gtrsim t^{2\overline a-2}|x|^2.
\end{aligned}
\end{equation*}
Thus we have for $\eps$ defined above
\begin{equation*}
\exists 0<c_{\eps}<1,r_{\eps}>0,\forall |x|<r_{\eps},\forall t\geq\eps,\ \ g'(t) \geq c_{\eps}|x|^2
\end{equation*}
and it leads to
\begin{equation*}
\forall t\geq\eps,\ g(t) \leq (1-(1-t)c_{\eps})|x|^2 < r_{\eps}^2.
\end{equation*}

\ep

We can now prove the proposition.

\bp
Up to a translation, we can consider that $\s=0$ and $W(\s)=0$. Thanks to Assumption \ref{ass:morsedegen2}, we know that there exists $U, (t_i)_i\subset\R^*, (\nu_i)_i\subset\N\setminus\{0,1\}$ such that $U$ is a smooth diffeomorphism in a neighborhood of $0$, $U(0) = 0$ and for $x$ in that neighborhood,
\begin{equation*}
W\circ U (x) = \sum_{i=1}^nt_ix_i^{\nu_i}.
\end{equation*}

We denote $X_0 = \{W<0\}$, let $\overline \nu = \sup\negthinspace_i\nu_i$, $a = (\overline\nu/\nu_i)_i\in[1,+\infty)^n$ and $x\in U^{-1}(X_0\cap B(0,r))$ with $0<r<r_0$, $r_0$ given by Lemma \ref{lem:path}. We see that
\begin{equation*}
\forall s\in[0,1],\ W\circ U (\gamma^a_{x,0}(s)) = \sum_{i=1}^nt_i(s^{\frac{\overline\nu}{\nu_i}}x_i)^{\nu_i} = s^{\overline\nu}W\circ U(x),
\end{equation*}
hence
\begin{equation}\label{eq:a}\tag{a}
\gamma^a_{x,0}((0,1])\subset U^{-1}(X_0).
\end{equation}
Moreover, applying Lemma \ref{lem:path} to $\phii = U^{-1}$ and $b=0$, we have
\begin{equation}\label{eq:b}\tag{b}
\gamma^a_{x,0}((0,1])\subset U^{-1}(B(0,r)),
\end{equation}
thus, combining \eqref{eq:a} and \eqref{eq:b}, we have
\begin{equation}\label{eq:c}\tag{c}
\gamma^a_{x,0}((0,1])\subset U^{-1}(X_0\cap B(0,r)).
\end{equation}
Now, let us notice that
\begin{equation}\label{eq:d}\tag{d}
\gamma^a_{x,0}(s)\xrightarrow[s\to0]{}0
\end{equation}
and
\begin{equation}\label{eq:e}\tag{e}
\forall \eta \ll 1,\ U^{-1}(X_0\cap B(0,r))\cap B(0,\eta) = U^{-1}(X_0)\cap B(0,\eta)
\end{equation}
because $U^{-1}$ is an open map since it is a diffeomorphism. Thus, up to proving that
\begin{equation}\label{eq:effe}\tag{f}
U^{-1}(X_0)\cap B(0,\eta) \mbox{ is connected}
\end{equation}
we have
\begin{equation*}
X_0\cap B(0,r) \mbox{ is connected}
\end{equation*}
combining \eqref{eq:c}, \eqref{eq:d}, \eqref{eq:e} and \eqref{eq:effe}, because connectedness is invariant by diffeomorphism. Therefore it just remains to prove \eqref{eq:effe}. Now let us show that if $0\notin\uuu^{(1)}$, then $U^{-1}(\{W<0\})\cap B(0,\eta) = \{x\in B(0,\eta),\ W\circ U(x)<0\}$ is connected for $\eta>0$ small enough.

We consider the case where $0\in\uuu^{odd}$ which means that there is at least one odd $\nu$. Without any loss of generality, we can consider that $\nu_1$ is odd and $t_1<0$ (in the case $t_1>0$, we just replace $\eta/2$ by $-\eta/2$ in the following).

Start from a point $(x_1,\ldots,x_n)$ in $\{x\in B(0,\eta),\ W\circ U(x)<0\}$, the goal is to connect it to $(\eta/2,0,\ldots,0)$ via a path that remains within the set. In the following, by saying that we link $a$ to $b$ we mean that we create a segment from one to another (so a path of the form $t\mapsto a + t(b-a)$), one can check that each time we do that, the whole path stays in $\{x\in B(0,\eta),\ W\circ U(x)<0\}$.

The first step of the path is to link all the positive $t_ix_i^{\nu_i}$ to $0$, hence either $x_1=0$ or $t_1x_1^{\nu_1}<0$ and all other non-zero contributions to $W\circ U$ are negative ones. Then we halve all the $x_i$ (so we link $x_i$ to $x_i/2$), this way, the ending point is in $B(0,\eta/2)$ and we know that its first coordinate is non-negative. Now we link $x_1$ to $\eta/2$, because we were in $B(0,\eta/2)$ and we had $x_1\geq0$, we indeed remain in $B(0,\eta)$. And at last, we link all the other coordinates to $0$, this way, we finally reached the aimed point.

For $0\in\uuu^{even}$ the setting is extremely similar to the well-known Morse case and the proof is exactly the same.

Recalling that for the sets we considered throughout the proof (open subsets of the Euclidean space), being connected and arc-connected is equivalent, we have proven the proposition.

\ep











































\subsection{Labeling in the generic case}\label{sec:labeling}

Now, among the saddle points of $W$ near a given minima $\m$, not all are important for the study of the eigenvalue associated with $\m$. Heuristically, the crucial ones are those of minimum height (that means they minimize $W(\s)$) such that a process stuck around $\m$ needs to cross in order to fall into a well of lower energy (that is, a well associated with some $\m'$ such that $W(\m') < W(\m)$).

\begin{defin}\label{def:labeling}
    We say that $\s\in\uuu$ is a separating saddle point if, for every $r>0$ small enough, $\{x\in B(\s,r),\ W(x)<W(\s)\}$ is composed of two connected components that are contained in two different connected components of $\{x\in\R^n,\ W(x)<W(\s)\}$. Hence, from Proposition \ref{prop:saddle} we have that the set of these points is a subset of $\uuu^{(1)}$, we will then denote it by $\vvv^{(1)}$, we also denote by separating saddle value of $W$ a point in $V(\vvv^{(1)})$.

    We say that $E\subset\R^n$ is a critical component of $W$ if there exists $\sigma\in V(\vvv^{(1}))$ such that $E$ is a connected component of $\{W\leq\sigma\}$ and $\p\negthinspace E\cup \vvv^{(1)}\neq\emptyset$.
\end{defin}

\begin{figure}[h]
    \centering
    \includegraphics[width=\linewidth]{separating.png}
    \caption{Example of a separating and %%
%% Annotation 62, 3/3 on page 57 [ ]
%% Delete: "a non<SEL>-</SEL> separating saddle"
%% Comment: ""
%%
%⭣ ⭣ ⭣
a non-separating saddle %%
%⭡ ⭡ ⭡ END of annotation 62
point}
    \label{fig:separating}
\end{figure}

Consider a potential $W$ having the level sets of Figure \ref{fig:separating}, with $\m_1,\m_2$ minima, $\s_1,\s_2$ saddle points such %%
%% Annotation 61, 2/3 on page 57 [ ]
%% Text: "None, annot.type = <Text>"
%% Comment: "please double check that figure is within margins"
%%
%⭣ ⭣ ⭣
that
\begin{equation*}
W(\m_1) < W(\m_2) < W(\s_1) < W(\s_2).
\end{equation*}
And %%
%⭡ ⭡ ⭡ END of annotation 61
the capital letters denote connected components of some subsets the following %%
%% Annotation 60, 1/3 on page 57 [ ]
%% Text: "None, annot.type = <Text>"
%% Comment: "please adjust to fit"
%%
%⭣ ⭣ ⭣
way
\begin{equation*}
A_1\sqcup A_2 = \{W > W(\s_2)\},\ B = \{W(\s_1) < W < W(\s_2)\},\ C_1\sqcup C_2 = \{W < W(\s_1)\}.
\end{equation*}
Observe %%
%⭡ ⭡ ⭡ END of annotation 60
that $\s_1$ is a separating saddle point but not $\s_2$. A process starting near $\m_2$ can go to a well of lower energy by crossing $\s_1$ and going to $\m_1$. However, a process starting near $\m_1$ cannot reach below $\m_1$ by crossing $\s_2$ because both connected components of $\{x\in B(\s_2,r),\ W(x)<W(\s_2)\}$ are in the same connected component of $\{x\in \R^n,\ W(x)<W(\s_2)\}$, and so the process would just go back to $\m_1$.

Let us now recall the labeling procedure. Under the assumptions \ref{ass.confin} and \ref{ass:morsedegen2}, we have that $W(\vvv^{(1}))$ is finite. We denote $N = \sharp W(\vvv^{(1})) + 1$ and by $\sigma_2 > \sigma_3 > \ldots > \sigma_N$ the elements of $W(\vvv^{(1}))$, for convenience, we also introduce a fictive infinite saddle value $\sigma_1 = +\infty$. Starting from $\sigma_1$, we recursively associate to each $\sigma_i$ a finite family of local minima $(\m_{i,j})_j$ and a finite family of critical component $(E_{i,j})_j$:
\begin{itemize}
    \item Let $X_{\sigma_1} = \{x\in\R^n,\ W(x) < \sigma_1\} = \R^n$. We let $\m_{1,1}$ be any global minimum of $V$ and $E_{1,1} = \R^n$. In the following we will write $\um = \m_{1,1}$.
    \item Next, we consider $X_{\sigma_2} = \{x\in \R^n,\ W(x) < \sigma_2\}$. This is the union of its finitely many connected components. Exactly one contains $\um$ and the other components are denoted by $E_{2,1},\ldots,E_{2,N_2}$. They are all critical and, in each component $E_{2,j}$ we pick up a point $\m_{2,j}$ which is a global minimum of $W_{|E_{2,j}}$.
    \item Suppose now that the families $(\m_{k,j})_j$ and $(E_{k,j})_j$ have been constructed until rank $k = i-1$. The set $X_{\sigma_i} = \{x\in\R^n,\ W(x) < \sigma_i\}$ has again finitely many connected components and we label $E_{i,j}$, $j\in\lb1,N_i\rb$, those of these components that do not contain any $\m_{k,j}$ for $k<i$. They are all critical and, in each $E_{i,j}$, we pick up a point $\m_{i,j}$ which is a global minimum of $W_{|E_{i,j}}$.
\end{itemize}
At the end of this procedure, all the minima have been labeled. Throughout, we denote by $\s_1$ a fictive saddle point such that $W(\s_1) = \sigma_1 = \infty$ and, for any set $A$, $\ppp(A)$ denotes the power set of $A$. From the above labeling, we define two mappings
\begin{equation*}
E:\uuu^{(0)}\to\ppp(\R^n)\ \ \mbox{ and }\ \ \j:\uuu^{(0)}\to\ppp(\vvv^{(1)}\cup\{\s_1\})
\end{equation*}
as follows: for every $i\in\lb1,N\rb$ and every $j\in\lb1,N_i\rb$,
\begin{equation*}
E(\m_{i,j}) = E_{i,j},
\end{equation*}
and
\begin{equation*}
\j(\um) = \{\s_1\}\ \ \mbox{ and }\ \ \j(\m_{i,j}) = \p\negthinspace E_{i,j}\cap \vvv^{(1)}\ \mbox{ for }i\geq2.
\end{equation*}
In particular, we have $E(\um) = \R^n$ and, for all $i\in\lb1,N\rb$, $j\in\lb1,N_i\rb$, one has $\emptyset\neq\j(\m_{i,j})\subset\{W=\sigma_i\}$. We then define the mappings
\begin{equation*}
\bsigma:\uuu^{(0)}\to W(\vvv^{(1)})\cup\{\sigma_1\}\ \ \mbox{ and }\ \ S:\uuu^{(0)}\to(0,+\infty],
\end{equation*}
by
\begin{equation}\label{eq:defS}
\forall \m\in\uuu^{(0)},\ \ \bsigma(\m) = W(\j(\m))\ \ \mbox{ and }\ \ S(\m) = \bsigma(\m) - W(\m),
\end{equation}
where, with a slight abuse of notation, we have identified $W(\j(\m))$ with its unique element. Note that $S(\m) = \infty$ if and only if $\m = \um$. 

We now consider the following generic assumption in order to lighten the result and the %%
%% Annotation 64, 2/2 on page 59 [ ]
%% Highlight: "METASTABILITY OF <SEL>DEGENERATE KFP EQUATI</SEL>ONS AT"
%% Comment: "AU: please suggest shorter title for running head"
%%
%⭣ ⭣ ⭣
proof.
\begin{equation}\label{ass.gener2}\tag{Gener}
\begin{array}{l}
    (\ast)\mbox{ for any }\m\in\uuu^{(0)},\m \mbox{ is the unique global minimum of } W_{|E(\m)}\\
    (\ast)\mbox{ for all }\m\neq\m'\in\uuu^{(0)},\j(\m)\cap\j(\m')=\emptyset.
\end{array}
\end{equation}
In %%
%⭡ ⭡ ⭡ END of annotation 64
particular, \eqref{ass.gener2} implies that $W$ uniquely attains its global minimum at $\um$. This assumption allows us to avoid some heavy constructions regarding the set $\uuu$ and lighten the definition of the quasimodes, see \cite{Mi19} for the general setting. But it seems that its not a true obstruction and that we can pursue the computations without this assumption as described in \cite[Section 6]{BoLePMi22} and \cite{No24}, in the spirit of \cite{Mi19}.

This assumption comes from \cite[(1.7)]{BoGaKl05_01} and \cite[Assumption 3.8]{HeKlNi04_01} where they were hard to handle properly. Then it changed to become \cite[Hypothesis 5.1]{HeHiSj11_01} when finally reaching the form of \cite[Assumption 4]{LePMi20}.

One can show that \eqref{ass.gener2} is weaker than \cite{BoGaKl05_01}'s, \cite{HeKlNi04_01}'s and \cite{HeHiSj11_01}'s assumptions. More precisely, they supposed that $\j(\m)$ is a singleton while here we have no restriction on the size of $\j(\m)$.


































\section{Some technical results}

For the four following lemmas, by $Q = O(r(h))$ (for $Q$ operator and $r$ some positive real function) we mean there exists $C>0$ such that $\vvert{Q} \leq Cr(h)$.

\begin{lemma}\label{lem:1}
For $s\geq0$
\begin{equation*}
(h^{2-\frac2{\overline{\nu}}}g_1(h) + \d_V^*G\d_V)^{-s} = O((h^{2-\frac2{\overline{\nu}}}g_1(h))^{-s}).
\end{equation*}
\end{lemma}
\bp We use that $\d_V^*G\d_V\geq0$ and functional calculus.

\ep

\begin{lemma}\label{lem:2}
    $\alpha\cdot \d_V\Pi(h^{2-\frac2{\overline{\nu}}}g_1(h) + \d_V^*G\d_V)^{-1/2} = O(1)$%%
%% Annotation 63, 1/2 on page 59 [ ]
%% Highlight: "V G <SEL>dV )−1/2u, (h2−</SEL>2  V"
%% Comment: "multline-like break across arguments (between comma)"
%%
%⭣ ⭣ ⭣
.
\end{lemma}
\bp
\begin{equation*}
\begin{split}
    (\ast_1) :&= \vvert{\alpha\cdot \d_V\Pi(h^{2-\frac2{\overline{\nu}}}g_1(h) + \d_V^*G\d_V)^{-1/2}u}^2\\
    &= \<(\alpha\cdot \d_V\Pi)^*\alpha\cdot \d_V\Pi(h^{2-\frac2{\overline{\nu}}}g_1(h) + \d_V^*G\d_V)^{-1/2}u,(h^{2-\frac2{\overline{\nu}}}g_1(h) + \d_V^*G\d_V)^{-1/2}u\>\\
    &= \<\d_V^*G\d_V\Pi(h^{2-\frac2{\overline{\nu}}}g_1(h) + \d_V^*G\d_V)^{-1/2}u,(h^{2-\frac2{\overline{\nu}}}g_1(h) + \d_V^*G\d_V)^{-1/2}u\>\\
    &\leq \vvert{\d_V^*G\d_V(h^{2-\frac2{\overline{\nu}}}g_1(h) + \d_V^*G\d_V)^{-1}}\vvert{u}^2,
\end{split}
\end{equation*}
having %%
%⭡ ⭡ ⭡ END of annotation 63
that $\d_V^*G\d_V$ commutes with $(h^{2-\frac2{\overline{\nu}}}g_1(h) + \d_V^*G\d_V)^{-1/2}$ from functional calculus, hence the result.

\ep

In this Lemma we can replace $\alpha\Pi$ by $G^{1/2}$ and obtain the same result.

\begin{lemma}\label{lem:3}
    $\d_V(h^{2-\frac2{\overline{\nu}}}g_1(h) + \d_V^*G\d_V)^{-1/2} = O(g_1(h)^{-1/2})$%%
%% Annotation 66, 2/2 on page 60 [ ]
%% Highlight: ") u    <SEL>ν g1(h) + d∗  ν g1(h) + d∗</SEL>  = <SEL>⟨G−1G1/2 dV (h2−2  V G dV )−1/2u, G1/2 dV (h2−2  V G dV )</SEL>−1/2u⟩  ≤"
%% Comment: "multline-like break between arguments (after comma)"
%%
%⭣ ⭣ ⭣
.
\end{lemma}
\bp
\begin{equation*}
\begin{split}
    (\ast_2) :&= \vvert{\d_V(h^{2-\frac2{\overline{\nu}}}g_1(h) + \d_V^*G\d_V)^{-1/2}u}^2\\
    &= \vvert{G^{-1/2}G^{1/2}\d_V(h^{2-\frac2{\overline{\nu}}}g_1(h) + \d_V^*G\d_V)^{-1/2}u}^2\\
    &= \<G^{-1}G^{1/2}\d_V(h^{2-\frac2{\overline{\nu}}}g_1(h) + \d_V^*G\d_V)^{-1/2}u,G^{1/2}\d_V(h^{2-\frac2{\overline{\nu}}}g_1(h) + \d_V^*G\d_V)^{-1/2}u\>\\
    &\leq g_1(h)^{-1}\vvert{G^{-1/2}\d_V(h^{2-\frac2{\overline{\nu}}}g_1(h) + \d_V^*G\d_V)^{-1/2}u}^2\\
    &\leq g_1(h)^{-1}\vvert{u}^2
\end{split}
\end{equation*}
using %%
%⭡ ⭡ ⭡ END of annotation 66
Assumption \ref{ass.G}, and a result similar to Lemma \ref{lem:2}.

\ep

\begin{lemma}\label{lem:4}
    For $i\in\lb1,d\rb$, denoting $\d_{V,i} = h\p_{x_i} + \p_iV$ so that $\d_V = (\d_{V,i})_i$, we have for all $i,j\in\lb1,d\rb$
    \begin{equation*}
    \d_{V,i}^*\d_{V,j}(h^{2-\frac2{\overline{\nu}}}g_1(h) + \d_V^*G\d_V)^{-1} = O(h^{\frac{2}{\overline{\nu}} - 1}g_1(h)^{-3/2}g_2(h)^{1/2}).
    \end{equation*}
\end{lemma}
\bp
In the following, we denote $R = (h^{2-\frac2{\overline{\nu}}}g_1(h) + \d_V^*G\d_V)^{-1}$ for clarity.

Using the commutation rules
\begin{equation*}
[\d_{V,i}^*,\d_{V,j}] = -2h\p_{i,j}^2V;\quad [\d_{V,i}^*,R] = R[\d_V^*G\d_V,\d_{V,i}^*]R
\end{equation*}
we have
\begin{equation*}
\begin{split}
    \d_{V,i}^*\d_{V,j}R &= [\d_{V,i}^*,\d_{V,j}]R + \d_{V,j}\d_{V,i}^*R\\
    &= -2h\p_{i,j}^2VR + \d_{V,j}R\d_{V,i}^* + \d_{V,j}[\d_{V,i}^*,R]\\
    &= -2h\p_{i,j}^2VR + \d_{V,j}R\d_{V,i}^* + \d_{V,j}R[\d_V^*G\d_V,\d_{V,i}^*]R.
\end{split}
\end{equation*}
The last commutator gives
\begin{equation*}
\begin{split}
    [\d_V^*G\d_V,\d_{V,i}^*] &= [\d_V^*,\d_{V,i}^*]G\d_V + \d_V^*[G,\d_{V,i}^*]\d_V + \d_V^*G[\d_V,\d_{V,i}^*]\\
    &= \d_V^*h\p_iG\d_V + 2h\d_V^*G(\p_{i,k}^2V)_{k}
\end{split}
\end{equation*}
since $[\d_{V,i}^*,\d_{V,j}^*] = 0$ for all $i,j$. Using Assumptions \ref{ass.confin} and \ref{ass.G} $ii)$ we %%
%% Annotation 65, 1/2 on page 60 [ ]
%% Highlight: "  R1/2 d∗      <SEL>+ Ch  V G1/2     G1/2 dV R1/2     R1/</SEL>2    <SEL>  dV,j R1/2     R1/2 d∗  + 2Ch    dV,j R1/2     R1/2 d∗  V G1/2     G1/2   ∥R∥.</SEL>  "
%% Comment: "shift left"
%%
%⭣ ⭣ ⭣
have
\begin{equation*}
\begin{split}
    \vvert{\d_{V,i}^*\d_{V,j}R} &\leq 2Ch\vvert{R} + \vvert{\d_{V,j}R^{1/2}}\vvert{R^{1/2}\d_{V,i}^*}\\
    &\phantom{*****} + Ch\vvert{\d_{V,j}R^{1/2}}\vvert{R^{1/2}\d_V^*G^{1/2}}\vvert{G^{1/2}\d_VR^{1/2}}\vvert{R^{1/2}}\\
    &\phantom{*****} + 2Ch\vvert{\d_{V,j}R^{1/2}}\vvert{R^{1/2}\d_V^*G^{1/2}}\vvert{G^{1/2}}\vvert{R}.
\end{split}
\end{equation*}
Here %%
%⭡ ⭡ ⭡ END of annotation 65
we used that
\begin{equation*}
\begin{split}
    \vvert{R^{1/2}\d_V^*\p_iG\d_VR^{1/2}} &= \sup_{u\neq0}\frac{\<R^{1/2}\d_V^*\p_iG\d_VR^{1/2}u,u\>}{\vvert{u^2}}\\
    &= \sup_{u\neq0}\frac{\<\p_iG\d_VR^{1/2}u,\d_VR^{1/2}u\>}{\vvert{u^2}}\\
    &\lesssim \sup_{u\neq0}\frac{\<G\d_VR^{1/2}u,\d_VR^{1/2}u\>}{\vvert{u^2}}\\
    &= \vvert{R^{1/2}\d_V^*G\d_VR^{1/2}}\\
    &\leq \vvert{R^{1/2}\d_V^*G^{1/2}}\vvert{G^{1/2}\d_VR^{1/2}}
\end{split}
\end{equation*}

Now with Lemmas \ref{lem:2} and \ref{lem:3}, we %%
%% Annotation 67, 1/1 on page 61 [ ]
%% Line: "obtain   d∗  <SEL>   ≲</SEL>h  hence"
%% Comment: "break"
%%
%⭣ ⭣ ⭣
obtain
\begin{equation*}
\vvert{\d_{V,i}^*\d_{V,j}R} \lesssim h^{\frac{2}{\overline{\nu}} - 1}g_1(h)^{-1} + g_1(h)^{-1} + h^{\frac{1}{\overline{\nu}}}g_1(h)^{-1} + h^{\frac{2}{\overline{\nu}} - 1}g_1(h)^{-3/2}g_2(h)^{1/2}.
\end{equation*}
hence %%
%⭡ ⭡ ⭡ END of annotation 67
the result.

\ep

Therefore, we also have
\begin{equation*}
\Delta_V(h^{2-\frac2{\overline{\nu}}}g_1(h) + \d_V^*G\d_V)^{-1} = O(h^{\frac{2}{\overline{\nu}} - 1}g_1(h)^{-3/2}g_2(h)^{1/2}).
\end{equation*}










\subsection{Proof of Lemma \ref{lem:soussolaffine}}\label{sec:proofsoussolaffine}
Up to a translation, let us consider that $V\geq 0$. According to Assumption \ref{ass.confin}, we also have that there exists $R>0$ such that for all $|x|\geq R$, $|\nabla V(x)| \geq \frac1C$.

For $x_0\in\R^d$, we consider the maximal solution to the Cauchy problem
\begin{equation}\label{1eq:edo}
\left\{
\begin{aligned}
&\dot x(t) = -\nabla V(x(t)),
\\&x(0) = x_0.
\end{aligned}
\right.
\end{equation}
We then see that
\begin{equation*}
\frac{d}{dt}(V(x(t))) = -|\nabla V(x(t))|^2.
\end{equation*}

Let us first show that for $\lambda>0$, $x_0\in\R^d$ such that $V(x_0) \leq\lambda$ and $|x_0| > R + \lambda C$, we have a solution to \eqref{1eq:edo} defined on $\R_+$ and for all $t\geq0$, $|x(t)|\geq R$. Consider $t$ the supremum such that this is true, and by contradiction assume that $t<\infty$. Then
\begin{equation*}
\begin{aligned}
    |x(t) - x(0)| &= \big|\int_0^t-\nabla V(x(s))ds\big|\ \leq \int_0^t|\nabla V(x(s))|ds\\
    &\leq C\int_0^t|\nabla V(x(s))|^2ds = C(V(x_0) - V(x(t)) \leq \lambda C,
\end{aligned}
\end{equation*}
which leads to $|x(t)|\leq |x_0| + \lambda C$. But we also have that
\begin{equation*}
|x(t)| \geq |x_0| - |x(t) - x_0| > R + \lambda C - \lambda C = R
\end{equation*}
this is a contradiction with the maximality of $t$.

Now assume that such an $x_0$ exists, therefore we %%
%% Annotation 69, 2/2 on page 62 [ ]
%% Highlight: "d(V (x(s))) <SEL>=</SEL> V (x0)−"
%% Comment: "align with first ="
%%
%⭣ ⭣ ⭣
have
\begin{equation*}
0\leq V(x(t)) = V(x_0) + \int_0^td(V(x(s))) = V(x_0) - \int_0^t|\nabla V(x(s))|^2ds \leq \lambda - \frac1{C^2}t,
\end{equation*}
but %%
%⭡ ⭡ ⭡ END of annotation 69
because $x(t)$ is global, this is absurd for $t$ big enough. This leads to
\begin{equation}\label{1eq:sublevel}
\{V\leq\lambda\}\subset \overline B(0,R+\lambda C)
\end{equation}
where the right-hand side term is the closed ball centered at $0$ of radius $R+\lambda C$. Now let $x\in\R^d$ be such that $|x| > R$, therefore there exists $\eps>0$ such that $\lambda = \frac1C(|x| - R - \eps) > 0$. Hence we can write $|x| = R + \lambda C + \eps$, thus by \eqref{1eq:sublevel}, we know that $V(x) > \lambda = \frac1C|x| + a$, with $a = -\frac{R+\eps}{C}$. Therefore, there exists $\tilde a\in\R$ such that
\begin{equation}\label{1eq:borneinfV}
\forall |x| > R,\ \ V(x) > \frac1C|x| + \tilde a.
\end{equation}
Because $|x|\leq R$ is compact, we can find $b\in\R$ such that \eqref{1eq:borneinfV} is true on the whole space.










\subsection{Proof of Proposition \ref{prop:existence}}\label{ssec:existence}

We use \cite[Theorem 1.7]{Kh12} which states that a Lyapunov function for the associated deterministic differential equation ensures the result announced in the proposition. The system we have to study is
\begin{equation*}
\left\{
\begin{aligned}
&dx_t = \alpha(x_t,v_t)dt,\\
&dv_t = \beta(x_t,v_t)dt - 4\Sigma^T\Sigma v_tdt.
\end{aligned}
\right.
\end{equation*}
Consider the function $t\mapsto f(x_t,v_t)$%%
%% Annotation 68, 1/2 on page 62 [ ]
%% Highlight: "V +2β·ΣTΣvt−8|ΣTΣvt|2 <SEL>≤</SEL>h  2(divx"
%% Comment: "align with ="
%%
%⭣ ⭣ ⭣
,
\begin{equation*}
\frac{d(f(x_t,v_t))}{dt} = \alpha\cdot\p_xV + 2\beta\cdot\Sigma^T\Sigma v_t - 8|\Sigma^T\Sigma v_t|^2 \leq \frac h2 (\div_x\alpha + \div_v\beta) \lesssim f(x_t,v_t).
\end{equation*}
Here %%
%⭡ ⭡ ⭡ END of annotation 68
we used \eqref{eq:assskew} and Assumption \ref{ass.accretive}. We define
\begin{equation*}
L:t\mapsto\sqrt{f(x_t,v_t) - \min V + 1}.
\end{equation*}
It is clear that this is a Lyapunov function satisfying $L' \lesssim L$. Using Assumption \ref{ass.confin}, it also satisfies the other required hypotheses to apply \cite[Theorem 1.7]{Kh12}, namely: $\lim_{R\to\infty}\inf_{|x,v|\geq R}L = \infty$ and $L$ is globally Lipschitz. For the first one, it is a direct consequence of Lemma \ref{lem:soussolaffine}, and the global Lipschitz property is induced by the boundedness of the Hessian of $V$%%
%% Annotation 70, 1/1 on page 63 [ ]
%% Highlight: "V .  <SEL>□</SEL>  "
%% Comment: "remove new lines before \end{proof}"
%%
%⭣ ⭣ ⭣
.

\ep










\subsection{Proof of Proposition \ref{prop:accretive}}\label{ssec:accretive}

The %%
%⭡ ⭡ ⭡ END of annotation 70
idea is to mimic the proof of \cite[Theorem 15.1]{He13}.

Let $h>0$ be fixed. To show that $P$ admits a maximal accretive extension, it is first necessary to show that it admits an accretive extension, this comes from the skew-adjointness of $X$, as well as from the non-negativity of $N$. It therefore remains to show the maximal side, for that we use the criterion which tells us that $P$ is maximal accretive if $T=P+(2h\Tr(\Sigma^T\Sigma)+1)\Id$ has a dense image.

Let $u\in L^2(\R^{d+d'})$ such that
\begin{equation}\label{eq:uTphii}
\forall \phii\in\ccc_c^\infty(\R^{d+d'}),\ \< u,T\phii\> =0.
\end{equation}
We then must show that $u=0$. As $P$ is real, we can assume also is $u$. We split $X= \underbrace{\alpha\cdot h\p_x + \beta\cdot h\p_v}_{=X_1}\underbrace{+\frac h2(\div_x\alpha+\div_v\beta)}_{=X_0}$ where we can see $X_1$ is a homogeneous differential operator of order 1 and $X_0$ is a mere $\ccc^\infty$ function.

Under Assumption \ref{ass.Hormander}, $T^*$ is hypoelliptic and using that thanks to \eqref{eq:uTphii}, $T^*u=0$ in $\ddd'$, we have that $u\in\ccc^\infty(\R^{d+d'})$. We then consider $\zeta$ a $\ccc^\infty_c(\R)$ function satisfying $0\leq\zeta\leq1$, $\zeta = 1$ on $[-1,1]$ and $\zeta=0$ outside of $[-2,2]$ and we denote for $k\in\R_+^*$, $\zeta_k(x,v) = \zeta\Big(\frac{f(x,v)}{k}\Big)$ where $f$ is defined in \eqref{eq:f}. We thus have
\begin{equation*}
\begin{split}
    \<u,T\zeta_k^2 u\> &= \<u,X\zeta_k^2 u\> + \<\zeta_k u,(N+(2h\Tr(\Sigma^T\Sigma)+1)\Id)\zeta_k u\> + \<u, [N,\zeta_k]\zeta_k u\>
\end{split}
\end{equation*}
\begin{equation*}
\begin{split}
    \<u, [N,\zeta_k]\zeta_k u\> &= -h^2\<\p_v\zeta_k,2u\p_v(\zeta_k u)-\p_v(u\zeta_k u)\>\\
    &= -h^2\<\p_v\zeta_k,u^2\p_v\zeta_k+2\zeta_k(u\p_v u-\p_vu\ u)\>\\
    &= -h^2\vvert{u\p_v\zeta_k}^2
\end{split}
\end{equation*}
\begin{equation*}
\begin{split}
    \<\zeta_k u,(N+(2h\Tr(\Sigma^T\Sigma)+1)\Id)\zeta_k u\> &= h^2\vvert{\p_v(\zeta_k u)}^2 + \vvert{2|\Sigma^T\Sigma v|\zeta_k u}^2 + \vvert{\zeta_k u}^2
\end{split}
\end{equation*}
\begin{equation*}
\begin{split}
    \<u,X\zeta_k^2 u\> &= \<\zeta_k u,X\zeta_k u\> + \<u,[X,\zeta_k]\zeta_k u\>\\
    &= \<u,[X_1,\zeta_k]\zeta_k u\>\\
    &= \frac{h^2}{2k}\<u,(\div_x\alpha + \div_v\beta)\zeta'(f/k)\zeta_k u\>
\end{split}
\end{equation*}
where we used Assumption \ref{ass.skewadj} for the last estimate and the skew-adjointness of $X$ the line above. Using \eqref{eq:uTphii} with $\phii = \zeta_k^2u$, we obtain
\begin{equation*}
\begin{split}
    (\ast):&= h^2\vvert{\p_v(\zeta_k u)}^2 + \vvert{2|\Sigma^T\Sigma v|\zeta_k u}^2 + \vvert{\zeta_k u}^2\\
    &= h^2\vvert{u\p_v\zeta_k}^2 - \frac{h^2}{2k}\<u,(\div_x\alpha + \div_v\beta)\zeta'(f/k)\zeta_k u\>.
\end{split}
\end{equation*}
Using that $\p_v\zeta_k = \frac2k\Sigma^T\Sigma v\zeta'(f/k)$ and
\begin{equation*}
h^2\vvert{\p_v(\zeta_k u)}^2 + \vvert{2|\Sigma^T\Sigma v|\zeta_k u}^2-\frac{h^2}{k^2}\vvert{2|\Sigma^T\Sigma v|\zeta'(f/k)u}^2\geq0
\end{equation*}
for $k$ large enough, we %%
%% Annotation 72, 2/2 on page 64 [ ]
%% Highlight: "we have <SEL>1  k(divx α+divv β)ζ′(f/k) → 0.</SEL> Hence, with"
%% Comment: "display"
%%
%⭣ ⭣ ⭣
have
\begin{equation*}
\vvert{\zeta_k u}^2 \leq - \frac{h^2}{2k}\<u,(\div_x\alpha + \div_v\beta)\zeta'(f/k)\zeta_k u\>.
\end{equation*}
Therefore, using that $\supp\zeta' \cap (-1,1) = \emptyset$, we have $\frac1k(\div_x\alpha + \div_v\beta)\zeta'(f/k)\to0$. Hence, with the dominated convergence theorem, $u=0$ because $\zeta_ku\to u$%%
%⭡ ⭡ ⭡ END of annotation 72
.

\ep










\subsection{Proof of Lemma \ref{lem:situ2}}\label{ssec:lemma}

To this purpose, we define a valuation $\gamma$ over the monomials in the variables $h,x$ and $v$ by: $\gamma(ch^jx^av^b) = 2j+a+b$ for all $c\in\R^*$. We now prove by induction the %%
%% Annotation 71, 1/2 on page 64 [ ]
%% Caret: "ing property<Caret>  (P ) The t"
%% Comment: ":"
%%
%⭣ ⭣ ⭣
following property
\begin{equation*}\label{HR}\tag{$\Pr_n$}
\text{The terms of valuation }n\text{ in }\ell\text{ are even in }v.
\end{equation*}
We %%
%⭡ ⭡ ⭡ END of annotation 71
start at $n=1$, it has been proven when looking at the principal order of the eikonal equation that assuming $\xi_x \neq 0$ implies $\xi_v = 0$. Let $n\geq2$, assume $\Pr_m$ is true for all $m<n$, then look at \eqref{eq:decompjab} for $2j+a+b = n$ and $b=2b'+1$ is odd. We know that $\p_v\ell = O(v+h)$ and $\ell = O(x + v^2 + h)$. Thus, since the valuation $\gamma$ is multiplicative ($\gamma(uv) = \gamma(u)\gamma(v)$), all the factors of terms appearing in $p_{j,a,2b'+1}$ come from monomials of valuation less than $n$, hence they are even in $v$ by the hypothesis. Then, $\p_v\ell$ will give odd terms and $\ell$ even ones, therefore $p_{j,a,2b'+1}$ is even in $v$, but it is of the form $h^jx^av^{2b'+1}$ thus $p_{j,a,2b'+1} = 0$. Moreover, for all $k\in\lb1,a\rb$,
\begin{equation*}
\gamma(\ell_{j,a-k,2b'+1}) \leq \gamma(\ell_{j,a+1,2b'-1}) = \gamma(\ell_{j-1,a+1,2b'+1}) = n-1.
\end{equation*}
By the induction hypothesis, the monomials $\ell_{j,a-k,2b'+1}, \ell_{j,a+1,2b'-1}, \ell_{j-1,a+1,2b'+1}$ are $0$ and we obtain
\begin{equation}\label{eq:vdvmoinsdv2}
v\p_v\ell_{j,a,2b'+1} - \p_v^2\ell_{j-1,a,2b'+3} = 0
\end{equation}
for all $j,a,b'\in\N$ such that $2j+a+2b'+1 = n$. Therefore %%
%% Annotation 74, 2/2 on page 65 [ ]
%% Caret: "taking j = 0<Caret> we have  "
%% Comment: ","
%%
%⭣ ⭣ ⭣
taking $j=0$ we %%
%⭡ ⭡ ⭡ END of annotation 74
have
\begin{equation*}
\forall a,b'\in\N,a+2b'+1=n,\quad v\p_v\ell_{0,a,2b'+1} = 0.
\end{equation*}
Using now a quick induction on $j$ along %%
%% Annotation 73, 1/2 on page 65 [ ]
%% Caret: "g with (B.5)<Caret> we complete"
%% Comment: ","
%%
%⭣ ⭣ ⭣
with \eqref{eq:vdvmoinsdv2} we %%
%⭡ ⭡ ⭡ END of annotation 73
complete the induction on $n$%%
%% Annotation 88, 14/14 on page 66 [ ]
%% Highlight: "<SEL>References</SEL>  [1] M"
%% Comment: "make author titles non small caps globally"
%%
%⭣ ⭣ ⭣
.

\ep

\vfill
\begin{thebibliography}{10}

\bibitem{AsBoMi22}
{\sc M.~Assal, %%
%⭡ ⭡ ⭡ END of annotation 88
J.-F. Bony, and L.~Michel}, {\em Metastable diffusions with degenerate drifts}.
\newblock Online first, 2024.

\bibitem{Be22}
{\sc M.~Ben~Said}, {\em {K}ramers–{F}okker–{P}lanck operators with homogeneous potentials}, Mathematical Methods in the Applied Sciences, 45 (2022), pp.~914--927.

\bibitem{Bi05}
{\sc J.-M. Bismut}, {\em The hypoelliptic {L}aplacian on the cotangent bundle}, J. Amer. Math. Soc., 18 (2005), pp.~379--476.

\bibitem{BiLe05}
{\sc J.-M. Bismut and G.~Lebeau}, {\em The hypoelliptic {L}aplacian and analytic torsion. ({L}aplacien hypoelliptique et torsion analytique.)}, Comptes Rendus Mathematique, 341 (2005), pp.~113--118.

\bibitem{BoLePMi22}
{\sc J.-F. Bony, D.~Le~Peutrec, and L.~Michel}, {\em {Eyring-Kramers law for Fokker-Planck type differential operators}}, J. Eur. Math. Soc.,  (2024).
\newblock Published online first.

\bibitem{BoGaKl05_01}
{\sc A.~Bovier, V.~Gayrard, and M.~Klein}, {\em Metastability in reversible diffusion processes. {II}. {P}recise asymptotics for small eigenvalues}, J. Eur. Math. Soc., 7 (2005), pp.~69--99.

\bibitem{CaDoHeMiMoSc24}
{\sc K.~Carrapatoso, J.~Dolbeault, F.~H\'erau, S.~Mischler, C.~Mouhot, and C.~Schmeiser}, {\em Special macroscopic modes and hypocoercivity}, Journal of the European Mathematical Society,  (2024).

\bibitem{De25_1}
{\sc L.~Delande}, {\em Sharp spectral gap for some degenerate {W}itten {L}aplacians}.
\newblock arXiv:2410.21899, 2025.

\bibitem{DoMoSc15}
{\sc J.~Dolbeault, C.~Mouhot, and C.~Schmeiser}, {\em Hypocoercivity for linear kinetic equations conserving mass}, Trans. Amer. Math. Soc., 367 (2015), pp.~3807--3828.

\bibitem{He13}
{\sc B.~Helffer}, {\em Spectral theory and its applications}, vol.~139 of Cambridge Studies in Advanced Mathematics, Cambridge University Press, Cambridge, 2013.

\bibitem{HeKlNi04_01}
{\sc B.~Helffer, M.~Klein, and F.~Nier}, {\em Quantitative analysis of metastability in reversible diffusion processes via a {W}itten complex approach}, Mat. Contemp., 26 (2004), pp.~41--85.

\bibitem{HeNi06}
{\sc B.~Helffer and F.~Nier}, {\em Quantitative analysis of metastability in reversible diffusion processes via a {W}itten complex approach: the case with boundary}, M\'em. Soc. Math. Fr. (N.S.), 26 (2006).

\bibitem{HeSj85_01}
{\sc B.~Helffer and J.~Sj\"ostrand}, {\em Puits multiples en m\'ecanique semi-classique. {IV}. \'{E}tude du complexe de {W}itten}, Comm. Partial Differential Equations, 10 %%
%% Annotation 75, 1/14 on page 66 [ ]
%% Highlight: "pp. 245--340. <SEL>[14]</SEL> , From"
%% Comment: "B. Helffer and J. Sj\"ostrand, (not \bysame)"
%%
%% Annotation 87, 13/14 on page 66 [ ]
%% Delete: "10 (1985), <SEL>pp.</SEL> 245--340. [14]"
%% Comment: ""
%%
%⭣ ⭣ ⭣
(1985), pp.~245--340.

\bibitem{HeSj10_01}
\bysame, {\em From resolvent %%
%⭡ ⭡ ⭡ END of annotations 75, 87
bounds to semigroup %%
%% Annotation 76, 2/14 on page 66 [ ]
%% Highlight: "arXiv:1001.4171, (2010). <SEL>[15]</SEL> , Improving"
%% Comment: "replace \bysame with previous author name globally"
%%
%% Annotation 86, 12/14 on page 66 [ ]
%% Replace: "semigroup bounds, <SEL>arXiv:1001.4171</SEL>, (2010)."
%% Comment: "\arXiv{1001.4171}"
%%
%⭣ ⭣ ⭣
bounds}, arXiv:1001.4171,  (2010).

\bibitem{HeSj21_01}
\bysame, {\em Improving semigroup %%
%⭡ ⭡ ⭡ END of annotations 76, 86
bounds with resolvent estimates}, Integral Equations Operator %%
%% Annotation 77, 3/14 on page 66 [ ]
%% Replace: "Operator Theory, <SEL>93</SEL> (2021), pp."
%% Comment: "\textbf{93}"
%%
%⭣ ⭣ ⭣
Theory, 93 (2021), %%
%⭡ ⭡ ⭡ END of annotation 77
pp.~Paper No. 36, 41.

\bibitem{HeSjVi24}
{\sc B.~Helffer, J.~Sj\"ostrand, and J.~Viola}, {\em Discussing semigroup bounds with resolvent estimates}, Integral Equations %%
%% Annotation 85, 11/14 on page 66 [ ]
%% Replace: "Operator Theory<SEL>, 96</SEL> (2024). [17]"
%% Comment: " \textbf{96} "
%%
%⭣ ⭣ ⭣
Operator Theory, 96 (2024). %%
%⭡ ⭡ ⭡ END of annotation 85

\bibitem{HeHiSj08-2}
{\sc F.~H\'erau, M.~Hitrik, and J.~Sj\"ostrand}, {\em Tunnel effect for {K}ramers-{F}okker-{P}lanck type operators}, Ann. Henri %%
%% Annotation 83, 9/14 on page 66 [ ]
%% Delete: "9 (2008), <SEL>pp.</SEL> 209--274. [18]"
%% Comment: ""
%%
%% Annotation 84, 10/14 on page 66 [ ]
%% Replace: "Henri Poincar´e<SEL>, 9</SEL> (2008), pp."
%% Comment: " \textbf{9} "
%%
%⭣ ⭣ ⭣
Poincar\'e, 9 (2008), pp.~209--274. %%
%⭡ ⭡ ⭡ END of annotations 83, 84

\bibitem{HeHiSj11_01}
\bysame, {\em Tunnel effect and symmetries for {K}ramers-{F}okker-{P}lanck type operators}, J. Inst. Math. %%
%% Annotation 78, 4/14 on page 66 [ ]
%% Replace: "Math. Jussieu, <SEL>10</SEL> (2011), pp."
%% Comment: "\textbf{10}"
%%
%% Annotation 82, 8/14 on page 66 [ ]
%% Delete: "10 (2011), <SEL>pp.</SEL> 567--634. [19]"
%% Comment: ""
%%
%⭣ ⭣ ⭣
Jussieu, 10 (2011), pp.~567--634. %%
%⭡ ⭡ ⭡ END of annotations 78, 82

\bibitem{HeNi04}
{\sc F.~H\'erau and F.~Nier}, {\em Isotropic hypoellipticity and trend to equilibrium for the {F}okker-{P}lanck equation with a high-degree potential}, Arch. Ration. %%
%% Annotation 79, 5/14 on page 66 [ ]
%% Replace: "Mech. Anal., <SEL>171</SEL> (2004), pp."
%% Comment: "\textbf{171}"
%%
%% Annotation 80, 6/14 on page 66 [ ]
%% Delete: "Mech. Anal.<SEL>,</SEL> 171 (2004),"
%% Comment: ""
%%
%% Annotation 81, 7/14 on page 66 [ ]
%% Delete: "171 (2004), <SEL>pp.</SEL> 151--218.  "
%% Comment: ""
%%
%⭣ ⭣ ⭣
Mech. Anal., 171 (2004), pp.~151--218. %%
%⭡ ⭡ ⭡ END of annotations 79, 80, 81

\bibitem{HeSjSt05}
{\sc F.~H\'erau, J.~Sj\"ostrand, and C.~Stolk}, {\em Semiclassical analysis for the {K}ramers–{F}okker–{P}lanck equation}, Communications in Partial Differential Equations, 30 (2005), pp.~689--760.

\bibitem{He06}
{\sc F.~Hérau}, {\em Hypocoercivity and exponential time decay for the linear inhomogeneous relaxation {B}oltzmann equation}, Asymptot. Anal., 46 (2006), pp.~349--359.

\bibitem{Ho67}
{\sc L.~Hörmander}, {\em {Hypoelliptic second order differential equations}}, Acta Mathematica, 119 (1967), pp.~147 -- 171.

\bibitem{Kh12}
{\sc R.~Khasminskii}, {\em Stochastic stability of differential equations}, vol.~66, Springer Berlin, Heidelberg, 01 2012.

\bibitem{LePMi20}
{\sc D.~Le~Peutrec and L.~Michel}, {\em Sharp asymptotics for non-reversible diffusion processes}, Probability and Mathematical Physics, 1 (2020), pp.~3--53.

\bibitem{LeSaSt20}
{\sc B.~Leimkuhler, M.~Sachs, and G.~Stoltz}, {\em Hypocoercivity properties of adaptive {L}angevin dynamics}, {SIAM Journal on Applied Mathematics}, 80 (2020), pp.~1197--1222.

\bibitem{LeLeLaHi24}
{\sc A.~Leroy, B.~Leimkuhler, J.~Latz, and D.~J. Higham}, {\em Adaptive stepsize algorithms for {L}angevin dynamics}.
\newblock arXiv:2403.11993, 2024.

\bibitem{Li09}
{\sc W.-X. Li}, {\em Global hypoellipticity and compactness of resolvent for {F}okker-{P}lanck operator}, Ann. Sc. Norm. Super. Pisa Cl. Sci., 11 (2009).

\bibitem{Mi19}
{\sc L.~Michel}, {\em {A}bout small eigenvalues of the {W}itten {L}aplacian}, Pure Appl. Anal., 1 (2019), pp.~149--206.

\bibitem{NiSaWh24}
{\sc F.~Nier, X.~Sang, and F.~White}, {\em The {G}rushin problem for {B}ismut's hypoelliptic {L}aplacian}.
\newblock arxiv:2405.08389, 2024.

\bibitem{No23}
{\sc T.~Normand}, {\em Metastability results for a class of linear {B}oltzmann equations}, Ann. Henri Poincar\'{e}, 24 (2023), pp.~4013--4067.

\bibitem{No24}
\bysame, {\em Spectral analysis of a semiclassical random walk associated to a general confining potential}, 2024.

\bibitem{Sc11}
{\sc J.~Schur}, {\em Bemerkungen zur theorie der beschränkten bilinearformen mit unendlich vielen veränderlichen.}, Journal für die reine und angewandte Mathematik, 140 (1911), pp.~1--28.

\bibitem{Vi09}
{\sc C.~Villani}, {\em Hypocoercivity}, Mem. Amer. Math. Soc., 202 (2009), pp.~iv+141.
\end{thebibliography}
\end{document}
